%% Volume I, Chapter 22: Canonical Forms and Matrix Functions
\chapter{Canonical Forms and Matrix Functions}
\label{ch:canonical-forms}

\abstract{Not every matrix can be diagonalised, and for those that cannot, one wants to know how close to diagonal a change of coordinates can bring them. There are two answers. Unitary changes of coordinates, which never magnify errors, reach a triangular matrix; this is the form that numerical computation uses. Arbitrary changes of coordinates reach a matrix that is diagonal except for a few ones just above the diagonal; this form classifies matrices completely, but it is so fragile that no computer can find it reliably. Built on these two forms, functions of matrices receive a meaning, among them the exponential, which solves systems of linear differential equations.}

This chapter completes the part of the book on linear algebra. \Cref{ch:linear-maps} asked in which basis the matrix of an operator is as simple as possible, observed that the rank and the trace do not decide whether two matrices are similar, and prepared the ground with the chains of kernels $\ker T\subseteq\ker T^2\subseteq\cdots$ and the Fitting decomposition. \Cref{ch:eigenvalues} found that the diagonal form is reached exactly when every eigenvalue has as many independent eigenvectors as its multiplicity, met the Jordan blocks $J_n(\lambda)$ as the simplest matrices for which this fails, and proved that $A^k\to0$ exactly when the spectral radius is less than $1$; it left several statements to this chapter, among them that diagonalisable matrices are dense, that the exponent of $x-\mu$ in the minimal polynomial is the size of the largest Jordan block, and that the entries of $A^k$ are sums of polynomials in $k$ times powers of the eigenvalues. \Cref{ch:symmetric-matrices} proved that normal matrices are exactly those that a unitary matrix diagonalises and defined $f(A)$ for Hermitian matrices, and \cref{ch:svd} found that eigenvalues and singular values can differ greatly and stated Weyl's inequalities between them without proof. \Cref{ch:determinants} announced the identity $\det e^B=e^{\operatorname{tr}B}$, and \cref{ch:transcendental} proved that the solutions of $f'=\lambda f$ are the functions $ce^{\lambda x}$, a statement that asks to be extended to systems.

The first answer to the question of the simplest form is Schur's theorem: every complex square matrix is unitarily similar to an upper triangular matrix, $A=QTQ^{\mathsf H}$, with the eigenvalues on the diagonal of $T$ in any prescribed order. The proof is a short induction, and the theorem has many consequences: a second proof of the spectral theorem for normal matrices, short proofs of the Cayley--Hamilton theorem and of the spectral mapping theorem, the simultaneous triangularisation of commuting matrices, the solvability of the Sylvester equation, the density of diagonalisable matrices, a measure of the distance from normality, Weyl's inequalities, and norms that come arbitrarily close to the spectral radius. For real matrices there is a real form with blocks of order $2$ on the diagonal, and for real normal matrices a real block diagonal form, which describes orthogonal and antisymmetric matrices completely.

The second answer is the Jordan canonical form. Every complex square matrix is similar to a block diagonal matrix of Jordan blocks $J_k(\lambda)$, unique up to the order of the blocks. The construction splits $\mathbb C^n$ into the generalised eigenspaces $\ker(A-\lambda I)^n$, which have the dimensions $\operatorname{am}(\lambda)$, and finds in each of them a basis of chains for the nilpotent operator $A-\lambda I$. The numbers of blocks of each size are read off from the ranks of the powers of $A-\lambda I$, so that these ranks decide similarity. The Jordan form gives the powers of a matrix explicitly, and with them a complete proof of the form of the solutions of linear recurrences; but it is a discontinuous function of the matrix, and the transformation to it can be arbitrarily ill-conditioned, so that it is a tool of theory and not of computation.

The two forms give meaning to $f(A)$ for functions $f$ more general than polynomials. Three natural definitions, by diagonalisation, by power series and by the Jordan form, agree whenever they apply, and all of them make $f(A)$ a polynomial in $A$, the polynomial that interpolates $f$ and its derivatives at the eigenvalues. The most important matrix function is the exponential $e^{A}=\sum_kA^k/k!$, which solves the system of differential equations $\mathbf x'=A\mathbf x$, satisfies $e^{A+B}=e^Ae^B$ when $A$ and $B$ commute and only then in general, and turns antisymmetric matrices into rotations. The last section asks how $e^A$ is computed, finds that the obvious methods can fail badly, and introduces the method of scaling and squaring that the stardust library uses.

As in the preceding chapters, $\mathbb F$ denotes $\mathbb R$ or $\mathbb C$, and the scalars are complex unless stated otherwise: a real matrix is regarded as a complex matrix whenever its eigenvalues are needed, and the real forms of \cref{sec:real-schur} and of \cref{sec:jordan-form} are stated separately. We keep the notation of \cref{ch:eigenvalues}: $\operatorname{spec}A$ is the set of eigenvalues, $\operatorname{am}(\lambda)$ and $\operatorname{gm}(\lambda)$ are the algebraic and the geometric multiplicity, $m_A$ is the minimal polynomial, and $C_{a,b}$ is the real $2\times2$ matrix with the rows $(a,-b)$ and $(b,a)$. Singular values are numbered in decreasing order, as in \cref{ch:svd}. The Jordan block $J_k(\lambda)$ of order $k\ge1$ has $\lambda$ on the diagonal, $1$ just above it and zeros elsewhere, with $J_1(\lambda)=[\lambda]$, and we write
\[
N_k\coloneqq J_k(0),\qquad J_k(\lambda)=\lambda I_k+N_k ,
\]
so that $N_k$ maps $\mathbf e_1$ to $\mathbf 0$ and $\mathbf e_j$ to $\mathbf e_{j-1}$ for $j\ge2$; its powers $N_k^j$ have ones on the $j$th diagonal above the main diagonal, and $N_k^k=0$. A block diagonal matrix with the square blocks $A_1,\ldots,A_r$ on the diagonal is written $\operatorname{diag}(A_1,\ldots,A_r)$, as in \cref{ch:symmetric-matrices}.

The chapter is organised as follows. \Cref{sec:schur} proves Schur's theorem and its first consequences, \cref{sec:schur-consequences} the consequences that concern perturbations, norms and singular values, and \cref{sec:real-schur} the real Schur form and the real normal matrices. \Cref{sec:nilpotent} constructs bases of chains for nilpotent operators, \cref{sec:jordan-form} the Jordan form, and \cref{sec:powers-again} applies it to powers of matrices and to linear recurrences. \Cref{sec:jordan-unstable} explains why the Jordan form is not computed. \Cref{sec:matrix-functions} defines functions of matrices, \cref{sec:matrix-exponential} studies the exponential, and \cref{sec:computing-matrix-functions} looks at the computation of the exponential with the stardust library.

\section{The Schur Decomposition}
\label{sec:schur}

\subsection{Unitary Triangularisation}

\Cref{prop:block-triangular} showed that an upper triangular matrix records a chain of invariant subspaces $\operatorname{span}\{\mathbf e_1\}\subset\operatorname{span}\{\mathbf e_1,\mathbf e_2\}\subset\cdots$, and \cref{thm:gm-am} used one eigenvector to split off a block of a matrix. Repeating this step, and keeping the basis orthonormal at every step, gives the first canonical form.

\begin{theorem}[Schur decomposition]
\label{thm:schur}
Let $A\in\mathbb C^{n\times n}$ have the eigenvalues $\lambda_1,\ldots,\lambda_n$, counted with multiplicity and listed in any order. Then there are a unitary $Q\in\mathbb C^{n\times n}$ and an upper triangular $T\in\mathbb C^{n\times n}$ with
\[
A=QTQ^{\mathsf H},\qquad t_{ii}=\lambda_i\quad(i=1,\ldots,n).
\]
If $A$ is real and all its eigenvalues are real, $Q$ and $T$ can be chosen real, with $Q$ orthogonal.
\end{theorem}

\begin{proof}
By induction on $n$; for $n=1$ take $Q=[1]$ and $T=A$. Let $n\ge2$, and let $\mathbf q_1$ be a unit eigenvector of $A$ for $\lambda_1$, which exists because $\lambda_1$ is an eigenvalue. Extend $\mathbf q_1$ to an orthonormal basis $\mathbf q_1,\mathbf w_2,\ldots,\mathbf w_n$ of $\mathbb C^n$ (\cref{cor:orthonormal-basis}), and let $W$ be the $n\times(n-1)$ matrix with the columns $\mathbf w_j$, so that $Q_1\coloneqq[\,\mathbf q_1\ W\,]$ is unitary. Since $A\mathbf q_1=\lambda_1\mathbf q_1$, $\mathbf q_1^{\mathsf H}\mathbf q_1=1$ and $W^{\mathsf H}\mathbf q_1=\mathbf 0$,
\[
Q_1^{\mathsf H}AQ_1=\begin{bmatrix}\mathbf q_1^{\mathsf H}A\mathbf q_1&\mathbf q_1^{\mathsf H}AW\\W^{\mathsf H}A\mathbf q_1&W^{\mathsf H}AW\end{bmatrix}=\begin{bmatrix}\lambda_1&\mathbf b^{\mathsf H}\\\mathbf 0&A'\end{bmatrix},\qquad A'\coloneqq W^{\mathsf H}AW\in\mathbb C^{(n-1)\times(n-1)},
\]
with $\mathbf b^{\mathsf H}\coloneqq\mathbf q_1^{\mathsf H}AW$. By \cref{prop:charpoly-similar,thm:det-block-triangular}, $p_A(x)=(x-\lambda_1)\,p_{A'}(x)$, so the eigenvalues of $A'$, counted with multiplicity, are $\lambda_2,\ldots,\lambda_n$, by the uniqueness of the factorisation into linear factors (\cref{thm:eigenvalues-exist}). By the induction hypothesis $A'=Q'T'Q'^{\mathsf H}$ with a unitary $Q'$ and an upper triangular $T'$ with the diagonal $\lambda_2,\ldots,\lambda_n$ in this order. The matrix $Q\coloneqq Q_1\operatorname{diag}(1,Q')$ is unitary, as a product of unitary matrices (\cref{prop:orthonormal-columns}\,(iv)), and
\[
Q^{\mathsf H}AQ=\begin{bmatrix}1&\mathbf 0^{\mathsf T}\\\mathbf 0&Q'^{\mathsf H}\end{bmatrix}\begin{bmatrix}\lambda_1&\mathbf b^{\mathsf H}\\\mathbf 0&A'\end{bmatrix}\begin{bmatrix}1&\mathbf 0^{\mathsf T}\\\mathbf 0&Q'\end{bmatrix}=\begin{bmatrix}\lambda_1&\mathbf b^{\mathsf H}Q'\\\mathbf 0&T'\end{bmatrix}\eqqcolon T ,
\]
which is upper triangular with the diagonal $\lambda_1,\ldots,\lambda_n$. If $A$ and its eigenvalues are real, the eigenvector $\mathbf q_1$ can be chosen real (\cref{thm:eigenvalues-exist}\,(iii)), the orthonormal basis can be completed in $\mathbb R^n$, and $A'$ is real with real eigenvalues, so every matrix in the proof is real.
\end{proof}

The factorisation $A=QTQ^{\mathsf H}$ is a \term{Schur decomposition} of $A$, and $T$ is a \term{Schur form} of $A$. In the language of operators, the theorem says that the columns $\mathbf q_1,\ldots,\mathbf q_n$ of $Q$ form an orthonormal basis for which every subspace $\operatorname{span}\{\mathbf q_1,\ldots,\mathbf q_k\}$ is invariant under $A$: the $k$th column of $AQ=QT$ says that $A\mathbf q_k=\sum_{i\le k}t_{ik}\mathbf q_i$. The first column $\mathbf q_1$ is an eigenvector, but the others in general are not.

A small example shows the construction. The matrix
\begin{equation}
A=\begin{bmatrix}3&1\\-1&1\end{bmatrix}
\label{eq:schur-example}
\end{equation}
has the trace $4$ and the determinant $4$, so its characteristic polynomial is $x^2-4x+4=(x-2)^2$, with the double eigenvalue $2$. The matrix $A-2I$ has the rows $(1,1)$ and $(-1,-1)$ and rank $1$, so $\operatorname{gm}(2)=1$, and $A$ is defective; its eigenvectors are the multiples of $(1,-1)$. With $\mathbf q_1=(1,-1)/\sqrt2$ and $\mathbf q_2=(1,1)/\sqrt2$, which completes it to an orthonormal basis, $A\mathbf q_2=(4,0)/\sqrt2$, so that $\mathbf q_1^{\mathsf T}A\mathbf q_2=2$ and $\mathbf q_2^{\mathsf T}A\mathbf q_2=2$, and
\[
A=QTQ^{\mathsf T},\qquad Q=\frac1{\sqrt2}\begin{bmatrix}1&1\\-1&1\end{bmatrix},\qquad T=\begin{bmatrix}2&2\\0&2\end{bmatrix}.
\]
The Schur form is not unique. The diagonal can be arranged in any order, and even for a fixed order the columns of $Q$ can be multiplied by numbers of modulus $1$: if $D$ is diagonal and unitary, $A=(QD)(D^{\mathsf H}TD)(QD)^{\mathsf H}$ is another Schur decomposition, whose triangular factor has the entries $\bar d_it_{ij}d_j$. Replacing $\mathbf q_1$ by $-\mathbf q_1$ in the example changes the sign of $t_{12}$. When eigenvalues are repeated, the freedom is larger still. \Cref{prop:schur-inequality} below identifies a quantity that does not depend on the choice.

Schur's theorem applies to operators as well. If $T$ is an operator on a complex inner product space $V$ of finite dimension and $\mathcal B$ an orthonormal basis, a Schur decomposition $[T]_{\mathcal B}=QRQ^{\mathsf H}$ defines a new basis $\mathcal B'$ whose $j$th vector has the coordinates $q_{1j},\ldots,q_{nj}$ with respect to $\mathcal B$. It is orthonormal, because the coordinate map of an orthonormal basis preserves inner products, and $[T]_{\mathcal B'}=Q^{-1}[T]_{\mathcal B}Q=R$ by \cref{thm:change-of-basis-map}.

\begin{corollary}[Schur's theorem for operators]
\label{cor:schur-operator}
Every operator on a complex inner product space of finite dimension has an upper triangular matrix with respect to some orthonormal basis.
\end{corollary}

\subsection{Normal Matrices Again}

\Cref{thm:normal-unitary} proved that a matrix is unitarily diagonalisable if and only if it is normal, by an induction that split off one eigenvector at a time. Schur's theorem gives a second proof, which reduces the question to triangular matrices.

\begin{lemma}[Normal triangular matrices]
\label{lem:triangular-normal}
A triangular matrix is normal if and only if it is diagonal.
\end{lemma}

\begin{proof}
A diagonal matrix $D$ is normal, since $D^{\mathsf H}D$ and $DD^{\mathsf H}$ are both the diagonal matrix of the numbers $|d_{ii}|^2$. Conversely, let $T$ be upper triangular and normal; the lower triangular case follows by taking conjugate transposes. We show by induction on $i$ that the $i$th row of $T$ has no nonzero entry off the diagonal. Suppose that this holds for the rows $1,\ldots,i-1$; then the entries $t_{ki}$ with $k<i$ vanish, so the $i$th column of $T$ has the single nonzero entry $t_{ii}$. Compare the $(i,i)$ entries of $T^{\mathsf H}T$ and $TT^{\mathsf H}$. The first is the squared length of the $i$th column, $\sum_k|t_{ki}|^2=|t_{ii}|^2$, and the second is the squared length of the $i$th row, $\sum_{j\ge i}|t_{ij}|^2$. They are equal, so $t_{ij}=0$ for every $j>i$. For $i=1$ the hypothesis about the earlier rows is empty, and the argument starts.
\end{proof}

\begin{corollary}[The spectral theorems from Schur's theorem]
\label{cor:schur-normal}
Let $A=QTQ^{\mathsf H}$ be a Schur decomposition of $A\in\mathbb C^{n\times n}$.
\begin{enumerate}
\item $A$ is normal if and only if $T$ is diagonal; every Schur decomposition of a normal matrix is a unitary diagonalisation.
\item $A$ is Hermitian if and only if $T$ is diagonal and real; skew-Hermitian if and only if $T$ is diagonal with imaginary diagonal; unitary if and only if $T$ is diagonal with entries of modulus $1$.
\end{enumerate}
\end{corollary}

\begin{proof}
(i) If $A^{\mathsf H}A=AA^{\mathsf H}$, then $T^{\mathsf H}T=Q^{\mathsf H}A^{\mathsf H}AQ=Q^{\mathsf H}AA^{\mathsf H}Q=TT^{\mathsf H}$, so $T$ is normal and, being triangular, diagonal by \cref{lem:triangular-normal}. Conversely, if $T$ is diagonal, $A$ is normal by the first part of the proof of \cref{thm:normal-unitary}. (ii) If $A^{\mathsf H}=A$, then $T^{\mathsf H}=T$; an upper triangular matrix that equals its conjugate transpose, which is lower triangular, is diagonal, and its diagonal is real. If $A^{\mathsf H}=-A$, the same argument gives an imaginary diagonal, and if $A^{\mathsf H}A=I$, then $A$ is normal and $T$ is a diagonal unitary matrix. The converses follow by reversing the computations, as at the end of the proof of \cref{thm:normal-unitary}.
\end{proof}

The two proofs of the spectral theorem for normal matrices are close relatives. Both split off an eigenvector with a unitary change of basis; the first proof shows directly that the off-diagonal block $\mathbf b$ vanishes, using the normality of $A$ at each step, while the second lets the induction run to the end without any hypothesis and uses normality once, on the finished triangle. The second proof also yields the Hermitian, skew-Hermitian and unitary cases of \cref{prop:hermitian-eigen,prop:unitary-eigen} together with their diagonalisability. With \cref{cor:schur-operator} and the fact that the matrix of the adjoint with respect to an orthonormal basis is the conjugate transpose (\cref{prop:adjoint-matrix}), the theorem extends to operators, as \cref{ch:inner-products} announced.

\begin{corollary}[Normal operators]
\label{cor:normal-operator}
An operator $T$ on a complex inner product space $V$ of finite dimension is normal if and only if $V$ has an orthonormal basis consisting of eigenvectors of $T$. It is then self-adjoint if and only if all eigenvalues are real, and unitary if and only if all eigenvalues have modulus $1$.
\end{corollary}

\begin{proof}
Let $\mathcal B'$ be an orthonormal basis with $[T]_{\mathcal B'}=R$ upper triangular (\cref{cor:schur-operator}). By \cref{prop:adjoint-matrix}, $[T^*]_{\mathcal B'}=R^{\mathsf H}$, so $T$ is normal if and only if $R$ is, if and only if $R$ is diagonal (\cref{lem:triangular-normal}), that is, if and only if the vectors of $\mathcal B'$ are eigenvectors. Conversely, if an orthonormal basis of eigenvectors exists, the matrix with respect to it is diagonal and hence normal, and so is $T$. The last statement is \cref{cor:schur-normal}\,(ii) for the matrix $R$.
\end{proof}

\subsection{Trace, Determinant, and Two Theorems Revisited}

The triangular form makes several results of \cref{ch:eigenvalues} transparent, because a triangular matrix shows its eigenvalues on its diagonal (\cref{prop:triangular-eigen}) and products and polynomials of upper triangular matrices are upper triangular, with diagonal entries computed from the diagonal entries alone: the $(i,i)$ entry of $ST$ is $\sum_ks_{ik}t_{ki}=s_{ii}t_{ii}$, since the terms with $k<i$ have $s_{ik}=0$ and those with $k>i$ have $t_{ki}=0$.

\begin{proposition}[Consequences of the Schur form]
\label{prop:schur-consequences}
Let $A\in\mathbb C^{n\times n}$ have the eigenvalues $\lambda_1,\ldots,\lambda_n$, counted with multiplicity.
\begin{enumerate}
\item $\operatorname{tr}A=\lambda_1+\cdots+\lambda_n$ and $\det A=\lambda_1\cdots\lambda_n$.
\item For every polynomial $q$, the eigenvalues of $q(A)$, counted with multiplicity, are $q(\lambda_1),\ldots,q(\lambda_n)$.
\item $p_A(A)=0$.
\end{enumerate}
\end{proposition}

\begin{proof}
Let $A=QTQ^{\mathsf H}$ with $t_{ii}=\lambda_i$. (i) The trace and the determinant are similarity invariants (\cref{prop:similarity-invariants}, \cref{cor:det-product}), and for the triangular $T$ they are the sum and the product of the diagonal entries (\cref{prop:det-triangular}). (ii) By \cref{prop:similarity-invariants}\,(ii), $q(A)=Qq(T)Q^{\mathsf H}$, and $q(T)$ is upper triangular with the diagonal entries $q(t_{ii})$, by the remark before the proposition applied to the powers of $T$ and their combinations. Its characteristic polynomial is therefore $\prod_i(x-q(\lambda_i))$, by \cref{prop:det-triangular} applied to $xI-q(T)$, and so is that of $q(A)$ (\cref{prop:charpoly-similar}). (iii) Let $V_k\coloneqq\operatorname{span}\{\mathbf e_1,\ldots,\mathbf e_k\}$, with $V_0\coloneqq\{\mathbf 0\}$, and $P_k\coloneqq(T-t_{11}I)(T-t_{22}I)\cdots(T-t_{kk}I)$. We show by induction on $k$ that $P_k$ maps $V_k$ to $\mathbf 0$. The matrix $T-t_{kk}I$ is upper triangular with a zero in position $(k,k)$, so it maps $\mathbf e_k$ into $V_{k-1}$, and it maps $V_{k-1}$ into itself; hence it maps $V_k$ into $V_{k-1}$. Since the factors of $P_k$ are polynomials in $T$, they commute, and $P_k=P_{k-1}(T-t_{kk}I)$ maps $V_k$ first into $V_{k-1}$ and then, by the induction hypothesis, to $\mathbf 0$; for $k=1$ the hypothesis is the statement $P_0V_0=\{\mathbf 0\}$ with $P_0=I$. For $k=n$, $P_n=p_A(T)$ vanishes on $V_n=\mathbb C^n$, and $p_A(A)=Qp_A(T)Q^{\mathsf H}=0$.
\end{proof}

Part (i) is \cref{thm:trace-det-eigen}, part (ii) the spectral mapping theorem (\cref{thm:spectral-mapping}), and part (iii) the Cayley--Hamilton theorem (\cref{thm:cayley-hamilton}). The earlier proofs worked with determinants: the coefficients of the characteristic polynomial, the factorisation of $x-q(t)$, and the adjugate of $\lambda I-A$. Here the work is done once, in the proof of Schur's theorem, and what remains is the observation that a triangular matrix carries its eigenvalues on its diagonal. The proof of (iii) is a matrix version of the following picture: the factor $T-t_{kk}I$ pushes the $k$th basis vector down into the span of its predecessors, and $n$ such factors push everything to zero.

\subsection{Commuting Matrices}

\Cref{thm:simultaneous} showed that commuting diagonalisable matrices can be diagonalised by one similarity. Without the hypothesis of diagonalisability, commuting matrices can still be triangularised together, and even by one unitary matrix.

\begin{theorem}[Simultaneous triangularisation]
\label{thm:simultaneous-triangular}
If $A,B\in\mathbb C^{n\times n}$ commute, there is a unitary $Q$ for which $Q^{\mathsf H}AQ$ and $Q^{\mathsf H}BQ$ are both upper triangular.
\end{theorem}

\begin{proof}
First, $A$ and $B$ have a common eigenvector. Let $\lambda$ be an eigenvalue of $A$. If $A\mathbf x=\lambda\mathbf x$, then $A(B\mathbf x)=BA\mathbf x=\lambda B\mathbf x$, so the eigenspace $E_\lambda(A)$ is invariant under $B$, and the restriction of $B$ to it, an operator on a nonzero complex space, has an eigenvector $\mathbf x$ (\cref{thm:eigenvalues-exist}, applied to its matrix in a basis). This $\mathbf x$ is an eigenvector of $B$ that lies in $E_\lambda(A)$. Now argue by induction on $n$, as in the proof of \cref{thm:schur}, with a unit common eigenvector $\mathbf q_1$ and $Q_1=[\,\mathbf q_1\ W\,]$ unitary. Both $Q_1^{\mathsf H}AQ_1$ and $Q_1^{\mathsf H}BQ_1$ have a first column $(\lambda,0,\ldots,0)$ and $(\mu,0,\ldots,0)$, with the lower right blocks $A'=W^{\mathsf H}AW$ and $B'=W^{\mathsf H}BW$. The product of two matrices of this block triangular form has the lower right block $A'B'$, so comparing the lower right blocks of the two sides of $(Q_1^{\mathsf H}AQ_1)(Q_1^{\mathsf H}BQ_1)=Q_1^{\mathsf H}ABQ_1=Q_1^{\mathsf H}BAQ_1=(Q_1^{\mathsf H}BQ_1)(Q_1^{\mathsf H}AQ_1)$ gives $A'B'=B'A'$. The induction hypothesis gives a unitary $Q'$ that triangularises $A'$ and $B'$, and $Q\coloneqq Q_1\operatorname{diag}(1,Q')$ triangularises $A$ and $B$, as in the proof of \cref{thm:schur}.
\end{proof}

If $Q^{\mathsf H}AQ$ and $Q^{\mathsf H}BQ$ are upper triangular with the diagonals $\lambda_1,\ldots,\lambda_n$ and $\mu_1,\ldots,\mu_n$, then $Q^{\mathsf H}(A+B)Q$ and $Q^{\mathsf H}ABQ$ are upper triangular with the diagonals $\lambda_i+\mu_i$ and $\lambda_i\mu_i$. So the eigenvalues of the sum and of the product of commuting matrices are obtained by pairing the eigenvalues of the two factors suitably, and in particular $\rho(A+B)\le\rho(A)+\rho(B)$ and $\rho(AB)\le\rho(A)\rho(B)$ for commuting matrices. For matrices that do not commute both inequalities fail: the matrices $N_2$ and $N_2^{\mathsf T}$ have the spectral radius $0$, while their sum, with the rows $(0,1)$ and $(1,0)$, has the eigenvalues $\pm1$, and their product $N_2N_2^{\mathsf T}=\operatorname{diag}(1,0)$ has the eigenvalue $1$. The converse of the theorem is false, since two upper triangular matrices need not commute; and the same proof triangularises any finite family of pairwise commuting matrices: the common eigenspace of the first $k$ matrices for given eigenvalues is invariant under each further matrix of the family, by the argument at the start of the proof, and so contains an eigenvector of the next one.

\subsection{The Sylvester Equation}

The Schur form also decides a question about linear equations whose unknown is a matrix. For $A\in\mathbb C^{m\times m}$, $B\in\mathbb C^{n\times n}$ and $C\in\mathbb C^{m\times n}$, the \term{Sylvester equation} is
\begin{equation}
AX+XB=C,\qquad X\in\mathbb C^{m\times n}.
\label{eq:sylvester}
\end{equation}
The map $X\mapsto AX+XB$ is linear on the space $\mathbb C^{m\times n}$ of dimension $mn$, so the equation is a linear system of $mn$ equations in $mn$ unknowns; the note on the Kronecker product in \cref{sec:matrix-of-map} wrote its matrix as $I_n\otimes A+B^{\mathsf T}\otimes I_m$. Whether it is uniquely solvable depends only on the eigenvalues.

\begin{theorem}[Solvability of the Sylvester equation]
\label{thm:sylvester-equation}
The equation \eqref{eq:sylvester} has a unique solution for every $C$ if and only if $A$ and $-B$ have no common eigenvalue.
\end{theorem}

\begin{proof}
Suppose first that $\lambda$ is an eigenvalue of $A$ and $-\lambda$ an eigenvalue of $B$. Then $A\mathbf x=\lambda\mathbf x$ for some $\mathbf x\ne\mathbf 0$, and since $B^{\mathsf T}$ has the same eigenvalues as $B$, there is $\mathbf y\ne\mathbf 0$ with $\mathbf y^{\mathsf T}B=-\lambda\mathbf y^{\mathsf T}$. The nonzero matrix $X\coloneqq\mathbf x\mathbf y^{\mathsf T}$ satisfies $AX+XB=\lambda\mathbf x\mathbf y^{\mathsf T}-\lambda\mathbf x\mathbf y^{\mathsf T}=0$, so the linear map $X\mapsto AX+XB$ is not injective, and the solution, when it exists, is not unique. Conversely, suppose that $A$ and $-B$ have no common eigenvalue. By \cref{cor:injective-surjective} it suffices to show that $AX+XB=0$ only for $X=0$. Let $A=URU^{\mathsf H}$ and $B=VSV^{\mathsf H}$ be Schur decompositions, and $Y\coloneqq U^{\mathsf H}XV$; multiplying the equation by $U^{\mathsf H}$ from the left and by $V$ from the right gives $RY+YS=0$. The $k$th column of this equation reads, since $S$ is upper triangular,
\[
R\mathbf y_k+\sum_{j\le k}s_{jk}\mathbf y_j=\mathbf 0,\qquad\text{that is,}\qquad(R+s_{kk}I)\mathbf y_k=-\sum_{j<k}s_{jk}\mathbf y_j .
\]
The matrix $R+s_{kk}I$ is upper triangular with the diagonal entries $r_{ii}+s_{kk}$, which are nonzero because $-s_{kk}$, an eigenvalue of $-B$, is not an eigenvalue $r_{ii}$ of $A$; so it is invertible (\cref{thm:triangular}\,(ii), which holds verbatim for complex matrices). For $k=1$ the right side is $\mathbf 0$, so $\mathbf y_1=\mathbf 0$, and by induction on $k$ every column vanishes. Hence $Y=0$ and $X=UYV^{\mathsf H}=0$.
\end{proof}

The proof is an algorithm. For a right-hand side $C$, the transformed equation $RY+YS=U^{\mathsf H}CV$ is solved column by column, each column by back substitution with a triangular matrix, at a cost of $O(m^2n+mn^2)$ operations once the two Schur decompositions are known, instead of the $O(m^3n^3)$ operations of elimination applied to the system of order $mn$. This is the method of Bartels and Stewart~\cite{bartels1972}, and the routine \code{linalg.solve\_sylvester} of the stardust library solves equations of the form \eqref{eq:sylvester}; Volume~VIII describes its algorithm and the computation of Schur forms, which starts by reducing the matrix to Hessenberg form (\code{linalg.hessenberg}) and then applies the QR iteration. The special case $B=-A$ shows that the commutator $X\mapsto AX-XA$ is never invertible, since $A$ and $A$ share all eigenvalues; indeed every polynomial in $A$ lies in its kernel.

\section{Consequences of Schur's Theorem}
\label{sec:schur-consequences}

\subsection{Diagonalisable Matrices Are Dense}

The note in \cref{sec:diagonalisation} stated that every complex matrix can be made diagonalisable by an arbitrarily small change, and outlined the proof with the Schur form. We now give it.

\begin{theorem}[Density of diagonalisable matrices]
\label{thm:diagonalisable-dense}
For every $A\in\mathbb C^{n\times n}$ and every $\varepsilon>0$ there is a matrix $A'\in\mathbb C^{n\times n}$ with $n$ pairwise different eigenvalues and $\|A-A'\|_{\mathrm F}<\varepsilon$. In particular, every complex matrix is the limit of a sequence of diagonalisable matrices.
\end{theorem}

\begin{proof}
Let $A=QTQ^{\mathsf H}$ be a Schur decomposition. Choose numbers $e_1,\ldots,e_n$ one after the other, with $|e_i|<\varepsilon/\sqrt n$ and $t_{ii}+e_i$ different from the $i-1$ numbers $t_{jj}+e_j$, $j<i$; this is possible because only finitely many values are excluded and the disc $|e|<\varepsilon/\sqrt n$ contains infinitely many numbers. Let $E\coloneqq\operatorname{diag}(e_1,\ldots,e_n)$ and $A'\coloneqq Q(T+E)Q^{\mathsf H}$. The triangular matrix $T+E$ has the pairwise different diagonal entries $t_{ii}+e_i$, which are its eigenvalues and those of $A'$, and $A'$ is diagonalisable by \cref{cor:distinct-diagonalisable}. Finally $\|A-A'\|_{\mathrm F}=\|QEQ^{\mathsf H}\|_{\mathrm F}=\|E\|_{\mathrm F}<\varepsilon$, because multiplication by a unitary matrix on either side does not change the Frobenius norm (\cref{prop:orthonormal-columns}\,(i), applied to $E$ and to $(QE)^{\mathsf H}$).
\end{proof}

The theorem makes a useful method of proof available: to establish an identity between continuous functions of the entries of a matrix, it suffices to establish it for diagonalisable matrices, or even for matrices with distinct eigenvalues. If two continuous functions $F,G\colon\mathbb C^{n\times n}\to\mathbb C$ agree on a dense set, they agree everywhere, since for every $A$ there is a sequence $A_k\to A$ from the set, and $F(A)=\lim F(A_k)=\lim G(A_k)=G(A)$. Two examples show the method at work.

The Cayley--Hamilton theorem asserts that the entries of $p_A(A)$ vanish. Each entry is a polynomial in the entries of $A$, since the coefficients of $p_A$ are polynomials in the entries (\cref{thm:charpoly-coefficients}) and so are the entries of the powers of $A$; it is therefore a continuous function of $A$. For a diagonalisable matrix, $p_A(A)=Sp_A(\Lambda)S^{-1}=0$, as \cref{sec:cayley-hamilton} observed. By the theorem, the entries vanish for every complex matrix, and so also for every real one. This is a third proof of the Cayley--Hamilton theorem, after the one with the adjugate and the one in \cref{prop:schur-consequences}.

For square matrices $A,B\in\mathbb C^{n\times n}$, \cref{thm:ab-ba} showed $p_{AB}=p_{BA}$ with Sylvester's determinant identity. The same follows from a density argument with a simpler dense set, the invertible matrices. If $A$ is invertible, $AB=A(BA)A^{-1}$ is similar to $BA$, and similar matrices have the same characteristic polynomial (\cref{prop:charpoly-similar}). For a general $A$, the matrix $A+tI$ is invertible for all but at most $n$ values of $t$, namely those with $-t\in\operatorname{spec}A$. Each coefficient of $p_{(A+tI)B}$ and of $p_{B(A+tI)}$ is a polynomial in $t$, and the two polynomials agree for infinitely many $t$, so they are equal (\cref{cor:factor}); at $t=0$ this gives $p_{AB}=p_{BA}$.

The method transfers \emph{identities} between continuous functions. It does not transfer properties that are not preserved under limits: diagonalisability itself is the obvious example, since every matrix is a limit of diagonalisable ones, and so are the properties of having distinct eigenvalues or of having a basis of eigenvectors with a bounded condition number. The density also does not make defective matrices harmless. A matrix close to a defective one is diagonalisable, but by \cref{sec:computing-eigenvalues} its eigenvalues are extremely sensitive, and \cref{sec:jordan-unstable} shows that its eigenvectors are nearly dependent.

\subsection{The Departure from Normality}

A Schur decomposition $A=QTQ^{\mathsf H}$ splits the triangular factor into its diagonal and its strictly upper triangular part, $T=D+N$, with $D=\operatorname{diag}(\lambda_1,\ldots,\lambda_n)$. By \cref{cor:schur-normal}, $A$ is normal exactly when $N=0$, and the size of $N$ measures how far $A$ is from normal.

\begin{proposition}[Schur's inequality and the departure from normality]
\label{prop:schur-inequality}
Let $A\in\mathbb C^{n\times n}$ have the eigenvalues $\lambda_1,\ldots,\lambda_n$, counted with multiplicity, and the singular values $\sigma_1,\ldots,\sigma_n$, and let $A=Q(D+N)Q^{\mathsf H}$ be a Schur decomposition with $D$ diagonal and $N$ strictly upper triangular. Then
\[
\sum_{i=1}^n\sigma_i^2=\|A\|_{\mathrm F}^2=\sum_{i=1}^n|\lambda_i|^2+\|N\|_{\mathrm F}^2 .
\]
In particular $\|N\|_{\mathrm F}$ is the same for every Schur decomposition of $A$, and
\[
\sum_{i=1}^n|\lambda_i|^2\le\|A\|_{\mathrm F}^2,
\]
with equality if and only if $A$ is normal.
\end{proposition}

\begin{proof}
The first equation is \cref{prop:singular-eigen}\,(v). The Frobenius norm is unitarily invariant, as in the proof of \cref{thm:diagonalisable-dense}, so $\|A\|_{\mathrm F}^2=\|D+N\|_{\mathrm F}^2$, and since $D$ and $N$ have no nonzero entries in common positions, this is $\sum_i|\lambda_i|^2+\|N\|_{\mathrm F}^2$. The left side and the sum of the $|\lambda_i|^2$ do not depend on the decomposition, so neither does $\|N\|_{\mathrm F}$. Equality in the inequality means $N=0$, which by \cref{cor:schur-normal} means that $A$ is normal.
\end{proof}

The inequality $\sum_i|\lambda_i|^2\le\sum_{i,j}|a_{ij}|^2$ is due to Schur, in the paper of 1909 that also contains the triangularisation~\cite{schur1909}. The number $\operatorname{dep}(A)\coloneqq\|N\|_{\mathrm F}=(\|A\|_{\mathrm F}^2-\sum_i|\lambda_i|^2)^{1/2}$ is Henrici's \term{departure from normality}~\cite{henrici1962}; it can be computed from the entries and the eigenvalues, without a Schur decomposition. For the defective matrix \eqref{eq:schur-example}, $\|A\|_{\mathrm F}^2=12$ and $\sum|\lambda_i|^2=8$, so $\operatorname{dep}(A)=2$, which is the entry $t_{12}$ found above, up to its sign. For the matrix $A_2$ of \cref{sec:matrix-powers}, with the rows $(4/5,2)$ and $(0,4/5)$, which is its own Schur form, $\operatorname{dep}(A_2)=2$ as well; \cref{ch:svd} found its singular values $1+\sqrt{41}/5=2.2806\ldots$ and $\sqrt{41}/5-1=0.2806\ldots$, whose squares add up to $\|A_2\|_{\mathrm F}^2=4+32/25$, while the squared moduli of its eigenvalues add up only to $32/25$. The singular values and the moduli of the eigenvalues coincide exactly for normal matrices (\cref{prop:singular-eigen}\,(iii)); the proposition shows that the difference between the two sums of squares is $\operatorname{dep}(A)^2$, so that the Schur form and the singular value decomposition carry the same information about normality. The departure from normality reappears in \cref{sec:computing-matrix-functions}, where it governs the transient behaviour of $e^{tA}$.

\subsection{Weyl's Inequalities}

\Cref{prop:singular-eigen}\,(iv) showed that the moduli of the eigenvalues of a square matrix lie between the smallest and the largest singular value, and that the product of all moduli equals the product of all singular values, $|\det A|$. The note on Weyl's inequalities in \cref{sec:svd-existence} stated a chain of inequalities between these two facts. With the Schur form and the interlacing of singular values it has a short proof.

\begin{theorem}[Weyl's inequalities]
\label{thm:weyl-product}
Let $A\in\mathbb C^{n\times n}$ have the singular values $\sigma_1\ge\cdots\ge\sigma_n$ and the eigenvalues $\lambda_1,\ldots,\lambda_n$, counted with multiplicity and ordered by decreasing modulus, $|\lambda_1|\ge\cdots\ge|\lambda_n|$. Then
\[
|\lambda_1\lambda_2\cdots\lambda_k|\le\sigma_1\sigma_2\cdots\sigma_k\qquad(k=1,\ldots,n),
\]
with equality for $k=n$.
\end{theorem}

\begin{proof}
By \cref{thm:schur} there is a Schur decomposition $A=QTQ^{\mathsf H}$ with the diagonal $t_{ii}=\lambda_i$ in the given order, and $T$ has the same singular values as $A$ (\cref{prop:singular-eigen}\,(ii)). Fix $k$, and let $T_k$ be the leading principal $k\times k$ submatrix of $T$, which is upper triangular with the diagonal $\lambda_1,\ldots,\lambda_k$. By \cref{prop:singular-eigen}\,(iv), applied to $T_k$, and \cref{prop:det-triangular},
\[
|\lambda_1\cdots\lambda_k|=|\det T_k|=\sigma_1(T_k)\cdots\sigma_k(T_k).
\]
The matrix $T_k$ arises from $T$ by deleting the last $n-k$ columns, one at a time, and then the last $n-k$ rows, one at a time. By \cref{thm:singular-interlacing}, each deletion of a column or a row does not increase any singular value, $\sigma_i(B)\le\sigma_i(C)$ when $B$ arises from $C$ by such a deletion and $i$ is at most the number of singular values of $B$. Hence $\sigma_i(T_k)\le\sigma_i(T)=\sigma_i(A)$ for $i=1,\ldots,k$, and the product of the first $k$ singular values of $T_k$ is at most $\sigma_1\cdots\sigma_k$. For $k=n$ both sides are $|\det A|$, by \cref{prop:singular-eigen}\,(iv).
\end{proof}

For $k=1$ the theorem is $\rho(A)\le\sigma_1=\|A\|_2$, the case $p=2$ of \cref{prop:spectral-radius-norm}. Dividing the case $k=n$ by the case $k=n-1$ gives the other end of \cref{prop:singular-eigen}\,(iv), $|\lambda_n|\ge\sigma_n$, when $A$ is invertible. The theorem was proved by Weyl in 1949~\cite{weyl1949}; the proof above, through a triangular submatrix, avoids the compound matrices that the note mentioned. For the matrix $A_2$, the case $k=1$ reads $4/5\le2.2806$, and the case $k=2$ is the equation $16/25=\sigma_1\sigma_2$. For a normal matrix all the inequalities are equalities, since the singular values are the moduli of the eigenvalues; for a matrix far from normal, the first eigenvalues can be much smaller in modulus than the first singular values, and the last ones correspondingly larger.

\subsection{Norms Close to the Spectral Radius}

Every operator norm bounds the spectral radius from above (\cref{prop:spectral-radius-norm}), and for $A_2$ the bound was poor: $\rho(A_2)=4/5$, while all the norms of \cref{ch:vectors-matrices} exceed $2$. The Schur form shows that a suitable norm brings the bound as close to $\rho(A)$ as one wishes. For an invertible $M\in\mathbb C^{n\times n}$, the function $\|\mathbf x\|_M\coloneqq\|M^{-1}\mathbf x\|_\infty$ is a norm on $\mathbb C^n$, since $M^{-1}$ is linear and injective, and for every matrix $A$ and every vector $\mathbf x$,
\begin{equation}
\|A\mathbf x\|_M=\|(M^{-1}AM)(M^{-1}\mathbf x)\|_\infty\le\|M^{-1}AM\|_\infty\,\|\mathbf x\|_M .
\label{eq:similarity-norm}
\end{equation}

\begin{proposition}[A norm close to the spectral radius]
\label{prop:norm-near-radius}
For every $A\in\mathbb C^{n\times n}$ and every $\varepsilon>0$ there is an invertible $M\in\mathbb C^{n\times n}$ with $\|M^{-1}AM\|_\infty\le\rho(A)+\varepsilon$, so that $\|A\mathbf x\|_M\le(\rho(A)+\varepsilon)\|\mathbf x\|_M$ for every $\mathbf x$.
\end{proposition}

\begin{proof}
Let $A=QTQ^{\mathsf H}$ be a Schur decomposition, and for $0<\tau\le1$ let $D_\tau\coloneqq\operatorname{diag}(1,\tau,\tau^2,\ldots,\tau^{n-1})$. The matrix $D_\tau^{-1}TD_\tau$ has the entries $\tau^{j-i}t_{ij}$: the diagonal is unchanged, and the entries above it are multiplied by positive powers of $\tau$. By \cref{thm:norm-formulas}, which holds verbatim for complex matrices,
\[
\|D_\tau^{-1}TD_\tau\|_\infty=\max_i\Bigl(|t_{ii}|+\sum_{j>i}\tau^{j-i}|t_{ij}|\Bigr)\le\rho(A)+\tau\max_i\sum_{j>i}|t_{ij}| ,
\]
since $|t_{ii}|\le\rho(A)$ and $\tau^{j-i}\le\tau$ for $j>i$. Choosing $\tau$ small enough and $M\coloneqq QD_\tau$, so that $M^{-1}AM=D_\tau^{-1}TD_\tau$, gives the claim, and \eqref{eq:similarity-norm} gives the last statement.
\end{proof}

The proposition gives a third proof that $\rho(A)<1$ implies $A^k\to0$, after the proof with the Cayley--Hamilton theorem in \cref{cor:powers-general} and the one with the Jordan form in \cref{sec:powers-again}. Choose $\varepsilon$ with $q\coloneqq\rho(A)+\varepsilon<1$. Applying the inequality $k$ times gives $\|A^k\mathbf x\|_M\le q^k\|\mathbf x\|_M$, and the entries of $A^k$ are the components of the vectors $A^k\mathbf e_j$, which satisfy $\|A^k\mathbf e_j\|_\infty=\|MM^{-1}A^k\mathbf e_j\|_\infty\le\|M\|_\infty\|A^k\mathbf e_j\|_M\le\|M\|_\infty q^k\|\mathbf e_j\|_M$, so they tend to zero geometrically. In matrix form, $A^k=M(M^{-1}AM)^kM^{-1}$ and the submultiplicativity of the $\infty$-norm give $\|A^k\|_\infty\le Cq^k$ with $C\coloneqq\|M\|_\infty\|M^{-1}\|_\infty$. This bound shows that the Neumann series $\sum_kA^k$ converges when $\rho(A)<1$, by comparison with the geometric series, which is the sufficiency in \cref{cor:neumann-spectral}, and it gives the upper half of Gelfand's formula (\cref{thm:gelfand}) once more: $\|A^k\|_\infty^{1/k}\le C^{1/k}(\rho(A)+\varepsilon)$, whose right side tends to $\rho(A)+\varepsilon$. A norm with $\|A\mathbf x\|\le\rho(A)\|\mathbf x\|$ for all $\mathbf x$ does not exist in general: for $N_2$, $\rho(N_2)=0$, while $N_2\mathbf e_2=\mathbf e_1\ne\mathbf 0$. The factor $\tau^{j-i}$ in the proof makes this visible: the off-diagonal part can be made small, but it does not disappear, and the price is the condition number $\|D_\tau\|_\infty\|D_\tau^{-1}\|_\infty=\tau^{1-n}$ of the scaling, which grows as $\tau$ decreases.

\section{The Real Schur Form}
\label{sec:real-schur}

\subsection{Quasi-Triangular Matrices}

For a real matrix with nonreal eigenvalues, the Schur decomposition of \cref{thm:schur} has complex factors, even though the matrix is real. Complex arithmetic costs several real operations per complex operation, and the conjugate pairs of eigenvalues of a real matrix suggest that a real form should exist in which each pair occupies a real block of order $2$. \Cref{prop:real-2x2-form} showed that a pair $a\pm ib$ comes with a real invariant plane, on which the matrix acts as $C_{a,b}$ in a suitable basis, and \cref{thm:real-block-diagonal} assembled these planes into a real block diagonal form for diagonalisable matrices. Using the planes one at a time, as Schur's theorem uses eigenvectors, gives a real triangular form for every real matrix.

\begin{theorem}[Real Schur decomposition]
\label{thm:real-schur}
Let $A\in\mathbb R^{n\times n}$. There is a real orthogonal $Q\in\mathbb R^{n\times n}$ for which
\[
Q^{\mathsf T}AQ=T=\begin{bmatrix}T_{11}&T_{12}&\cdots&T_{1m}\\&T_{22}&\cdots&T_{2m}\\&&\ddots&\vdots\\&&&T_{mm}\end{bmatrix}
\]
is block upper triangular, with zeros below the diagonal blocks, and each diagonal block $T_{jj}$ is either of order $1$, a real eigenvalue of $A$, or of order $2$, a real matrix whose eigenvalues are a pair $a\pm ib$ of nonreal eigenvalues of $A$. The real eigenvalues and the conjugate pairs can be made to appear along the diagonal in any prescribed order.
\end{theorem}

\begin{proof}
By induction on $n$; for $n=1$ take $Q=[1]$, and for $n=2$ with nonreal eigenvalues take $Q=I_2$, since $A$ is then a single block of order $2$. Let $n\ge2$, apart from this case. If the first eigenvalue in the prescribed order is real, proceed as in the proof of \cref{thm:schur} with a real unit eigenvector and a real orthonormal completion: $Q_1^{\mathsf T}AQ_1$ has the first column $(\lambda,0,\ldots,0)$ and a real lower right block $A'$ of order $n-1$. Otherwise the first entry in the order is a pair $a\pm ib$ with $b>0$. Let $\mathbf x=\mathbf u+i\mathbf v$ be an eigenvector for $a-ib$. By \cref{prop:real-2x2-form}, $\mathbf u$ and $\mathbf v$ are independent, and $A[\,\mathbf u\ \mathbf v\,]=[\,\mathbf u\ \mathbf v\,]C_{a,b}$. The Gram--Schmidt process (\cref{thm:gram-schmidt}) turns $\mathbf u,\mathbf v$ into an orthonormal basis $\mathbf q_1,\mathbf q_2$ of their span, with $[\,\mathbf q_1\ \mathbf q_2\,]=[\,\mathbf u\ \mathbf v\,]R$ for an invertible upper triangular $R\in\mathbb R^{2\times2}$, the inverse of the triangular factor of the QR factorisation of $[\,\mathbf u\ \mathbf v\,]$. Then
\[
A[\,\mathbf q_1\ \mathbf q_2\,]=[\,\mathbf u\ \mathbf v\,]C_{a,b}R=[\,\mathbf q_1\ \mathbf q_2\,]T_{11},\qquad T_{11}\coloneqq R^{-1}C_{a,b}R ,
\]
a real matrix similar to $C_{a,b}$, with the eigenvalues $a\pm ib$. Extend $\mathbf q_1,\mathbf q_2$ to an orthonormal basis of $\mathbb R^n$, with the remaining vectors as the columns of $W$, and let $Q_1\coloneqq[\,\mathbf q_1\ \mathbf q_2\ W\,]$. Since $W^{\mathsf T}A[\,\mathbf q_1\ \mathbf q_2\,]=W^{\mathsf T}[\,\mathbf q_1\ \mathbf q_2\,]T_{11}=0$, the matrix $Q_1^{\mathsf T}AQ_1$ has the block $T_{11}$ in its upper left corner, zeros below it, and a real lower right block $A'=W^{\mathsf T}AW$ of order $n-2$. In both cases $p_A=p_{T_{11}}\,p_{A'}$ by \cref{thm:det-block-triangular}, so the eigenvalues of $A'$ are those of $A$ with the first entry of the order removed. The induction hypothesis, applied to $A'$ with the rest of the order, and the same assembly as in the proof of \cref{thm:schur} complete the proof.
\end{proof}

A block upper triangular matrix with diagonal blocks of orders $1$ and $2$ is called \term{quasi-triangular}, and $A=QTQ^{\mathsf T}$ is a \term{real Schur decomposition}. The eigenvalues of $T$, and so of $A$, are the diagonal entries of order $1$ together with the eigenvalues of the blocks of order $2$, by \cref{prop:block-triangular-eigen}, and they can be computed from each block by the formula \eqref{eq:eigen-2x2}. This is the form in which computations with real matrices deliver eigenvalues: the routine \code{linalg.schur} of the stardust library, applied to a real matrix, returns a real quasi-triangular $T$ and a real orthogonal $Z$ with $A=ZTZ^{\mathsf T}$; for the matrix with the rows $(0,-2)$ and $(1,0)$, whose eigenvalues are $\pm i\sqrt2$, it returned the matrix itself as $T$ and the identity as $Z$. \Cref{thm:real-block-diagonal} and \cref{thm:real-schur} complement each other. The first reaches a block diagonal form with the blocks $C_{a,b}$, but only for diagonalisable matrices and with a similarity that may be far from orthogonal; the second applies to every real matrix and uses an orthogonal matrix, at the price of the blocks above the diagonal.

The real Schur form also proves the real version of the density theorem, which the note in \cref{sec:diagonalisation} derived from the discriminant: every real matrix is the limit of real matrices with $n$ different complex eigenvalues. Indeed, in $T$ one may add a small real number to each diagonal block of order $1$, and a small real multiple $\delta I_2$ to each block of order $2$, which moves its eigenvalues $a\pm ib$ to $a+\delta\pm ib$; choosing these numbers one after the other, small and avoiding the finitely many values that would create a coincidence with the eigenvalues of the earlier blocks, gives a real quasi-triangular $T'$ near $T$ with $n$ different eigenvalues, and $QT'Q^{\mathsf T}$ is near $A$.

\subsection{Real Normal Matrices}

For a real normal matrix the blocks above the diagonal disappear, as the off-diagonal part did in \cref{cor:schur-normal}, and the blocks of order $2$ take the form $C_{a,b}$. The first step generalises \cref{lem:triangular-normal}.

\begin{lemma}[Normal block triangular matrices]
\label{lem:block-triangular-normal}
Let
\[
T=\begin{bmatrix}T_{11}&T_{12}\\0&T_{22}\end{bmatrix}
\]
with square diagonal blocks be normal. Then $T_{12}=0$, and $T_{11}$ and $T_{22}$ are normal.
\end{lemma}

\begin{proof}
The upper left blocks of $T^{\mathsf H}T$ and $TT^{\mathsf H}$ are $T_{11}^{\mathsf H}T_{11}$ and $T_{11}T_{11}^{\mathsf H}+T_{12}T_{12}^{\mathsf H}$. They are equal, and the traces of $T_{11}^{\mathsf H}T_{11}$ and $T_{11}T_{11}^{\mathsf H}$ are both $\|T_{11}\|_{\mathrm F}^2$, so $\operatorname{tr}(T_{12}T_{12}^{\mathsf H})=\|T_{12}\|_{\mathrm F}^2=0$, and $T_{12}=0$. Then $T$ is block diagonal, and comparing the diagonal blocks of $T^{\mathsf H}T=TT^{\mathsf H}$ shows that $T_{11}$ and $T_{22}$ are normal.
\end{proof}

A real $2\times2$ matrix $M$ with the rows $(p,q)$ and $(r,s)$ and nonreal eigenvalues is normal only in a special way. The $(1,1)$ entries of $M^{\mathsf T}M$ and $MM^{\mathsf T}$ are $p^2+r^2$ and $p^2+q^2$, so normality forces $r=\pm q$. If $r=q$, the matrix is symmetric and its eigenvalues are real; so $r=-q$, and $q\ne0$, since otherwise $M$ is diagonal. The $(1,2)$ entries are $pq+rs=q(p-s)$ and $pr+qs=-q(p-s)$, and their equality forces $p=s$. So $M=C_{p,-q}$, and conjugating with the orthogonal matrix $\operatorname{diag}(1,-1)$, which changes $C_{a,b}$ into $C_{a,-b}$, turns it into $C_{p,|q|}$. Conversely $C_{a,b}^{\mathsf T}C_{a,b}=(a^2+b^2)I=C_{a,b}C_{a,b}^{\mathsf T}$, so every $C_{a,b}$ is normal.

\begin{theorem}[Real normal matrices]
\label{thm:real-normal}
A real matrix $A\in\mathbb R^{n\times n}$ is normal if and only if there is a real orthogonal $Q$ with
\[
Q^{\mathsf T}AQ=\operatorname{diag}\bigl(\lambda_1,\ldots,\lambda_r,C_{a_1,b_1},\ldots,C_{a_s,b_s}\bigr),\qquad b_j>0 ,
\]
where $\lambda_1,\ldots,\lambda_r$ are the real eigenvalues and $a_j\pm ib_j$ the pairs of nonreal eigenvalues of $A$, counted with multiplicity, and $r+2s=n$.
\end{theorem}

\begin{proof}
A block diagonal matrix with normal blocks is normal, and so is every matrix orthogonally similar to a normal matrix, as in the proof of \cref{cor:schur-normal}; this gives the sufficiency. Conversely, let $A$ be normal and $A=QTQ^{\mathsf T}$ a real Schur decomposition. Then $T$ is normal. Applying \cref{lem:block-triangular-normal} with $T_{11}$ the first diagonal block, and then again to the normal lower right block, and so on, shows that $T$ is block diagonal with normal diagonal blocks. A normal block of order $2$ with nonreal eigenvalues is $C_{a,b}$ or becomes $C_{a,b}$ with $b>0$ after conjugation with $\operatorname{diag}(1,-1)$, by the computation above; including these sign changes in $Q$, which remains orthogonal, gives the form of the theorem.
\end{proof}

The theorem completes \cref{thm:normal-unitary} for real matrices: a real normal matrix cannot in general be diagonalised by a real orthogonal matrix, but it can be brought to a block diagonal form in which each block of order $2$ is a rotation combined with a scaling, $C_{a,b}=rR_\theta$ with $a+ib=re^{i\theta}$. The three classes of real normal matrices that matter most are described by the possible eigenvalues. A real symmetric matrix has only real eigenvalues, so $s=0$ and the theorem is the real spectral theorem. For orthogonal and antisymmetric matrices we obtain the following forms.

\begin{corollary}[Orthogonal and antisymmetric matrices]
\label{cor:orthogonal-form}
\leavevmode
\begin{enumerate}
\item For every orthogonal $Q\in\mathrm O(n)$ there is an orthogonal $U$ with
\[
U^{\mathsf T}QU=\operatorname{diag}\bigl(I_p,-I_q,R_{\varphi_1},\ldots,R_{\varphi_s}\bigr),\qquad0<\varphi_j<\pi,\quad p+q+2s=n,
\]
and $\det Q=(-1)^q$.
\item For every antisymmetric $K\in\mathbb R^{n\times n}$, $K^{\mathsf T}=-K$, there is an orthogonal $U$ with
\[
U^{\mathsf T}KU=\operatorname{diag}\bigl(0,\ldots,0,C_{0,b_1},\ldots,C_{0,b_s}\bigr),\qquad b_j>0 ;
\]
in particular the rank of $K$ is even, $2s$.
\end{enumerate}
\end{corollary}

\begin{proof}
(i) An orthogonal matrix is normal, and its eigenvalues have modulus $1$ (\cref{prop:unitary-eigen}). The real eigenvalues are therefore $\pm1$; order them so that the eigenvalues $1$ come first. A block $C_{a,b}$ with $a^2+b^2=1$ and $b>0$ is the rotation $R_\varphi$ with $\cos\varphi=a$ and $\sin\varphi=b$, so $0<\varphi<\pi$. The determinant is the product of the determinants of the blocks, and $\det R_\varphi=1$. (ii) An antisymmetric matrix is normal, and its eigenvalues are imaginary (\cref{prop:hermitian-eigen}\,(iv)); its real eigenvalues are therefore $0$, and its blocks of order $2$ have $a=0$. The rank of a block diagonal matrix is the sum of the ranks of the blocks, and $C_{0,b}$ is invertible.
\end{proof}

Part (i) describes every orthogonal transformation of $\mathbb R^n$: there are mutually orthogonal lines and planes, filling the space, such that the transformation fixes the vectors of $p$ lines, reverses those of $q$ lines, and rotates each of $s$ planes by an angle $\varphi_j$. Since $R_\pi=-I_2$, two reversed lines can be combined into a plane rotated by $\pi$, and a rotation, with $\det Q=1$ and $q$ even, rotates mutually orthogonal planes and fixes the complement. For $n=3$ and $\det Q=1$ the possibilities are $p=3$, the identity; $p=1$ and $q=2$, a rotation by $\pi$ about an axis; and $p=1$ and $s=1$, a rotation by $\varphi_1$; this recovers \cref{prop:rotation-axis}\,(iii), with the angle chosen in $[0,\pi]$ by orienting the axis. The cyclic permutation matrix $P$ that maps $\mathbf e_1\mapsto\mathbf e_2\mapsto\mathbf e_3\mapsto\mathbf e_1$, which \cref{sec:unitary-matrices} identified as a rotation by a third of a turn about the diagonal $(1,1,1)$, illustrates the computation. With $\mathbf u_1=(1,1,1)/\sqrt3$, $\mathbf u_2=(1,-1,0)/\sqrt2$ and $\mathbf u_3=(1,1,-2)/\sqrt6$, one finds $P\mathbf u_2=(0,1,-1)/\sqrt2=-\mathbf u_2/2+(\sqrt3/2)\mathbf u_3$ and $P\mathbf u_3=(-2,1,1)/\sqrt6=-(\sqrt3/2)\mathbf u_2-\mathbf u_3/2$, so that
\[
U^{\mathsf T}PU=\begin{bmatrix}1&0&0\\0&-1/2&-\sqrt3/2\\0&\sqrt3/2&-1/2\end{bmatrix}=\operatorname{diag}\bigl(1,R_{2\pi/3}\bigr),\qquad U=[\,\mathbf u_1\ \mathbf u_2\ \mathbf u_3\,].
\]
Part (ii) describes the antisymmetric matrices. In $\mathbb R^3$, the matrix $W$ with $W\mathbf x=\mathbf w\times\mathbf x$ of \cref{sec:complex-eigenvalues} has the eigenvalues $0$ and $\pm i\|\mathbf w\|$, so its form is $\operatorname{diag}(0,C_{0,\|\mathbf w\|})$, with the first column of $U$ along $\mathbf w$; \cref{sec:matrix-exponential} turns it into a rotation about $\mathbf w$.

\section{Nilpotent Matrices and Jordan Chains}
\label{sec:nilpotent}

\subsection{Chains}

Schur's theorem achieves a triangular form with a unitary similarity; it does not say how far the triangle can be cleared by a general similarity. The answer begins with the nilpotent operators, the operators $N$ with $N^k=0$ for some $k$, whose only eigenvalue is $0$ (\cref{sec:cayley-hamilton}). For them the triangular form is strictly upper triangular, and the question is how few nonzero entries it can be made to have. The Jordan block $N_k$ is the model. It acts on the basis $\mathbf e_1,\ldots,\mathbf e_k$ as a chain of steps $\mathbf e_k\mapsto\mathbf e_{k-1}\mapsto\cdots\mapsto\mathbf e_1\mapsto\mathbf 0$, and the differentiation operator on $\mathcal P_n$, with the basis $1,x,x^2/2!,\ldots,x^n/n!$, acts in the same way (\cref{sec:similarity}).

\begin{definition}[Jordan chain]
\label{def:jordan-chain}
Let $N$ be an operator on a vector space $V$. A \term{Jordan chain} of length $k\ge1$ for $N$ is a list $N^{k-1}\mathbf v,\ldots,N\mathbf v,\mathbf v$ of vectors of $V$ with $N^{k-1}\mathbf v\ne\mathbf 0$ and $N^k\mathbf v=\mathbf 0$. The vector $\mathbf v$ is the \term{top} of the chain, and $N^{k-1}\mathbf v$, an eigenvector of $N$ for the eigenvalue $0$, is its \term{bottom}.
\end{definition}

The vectors of a chain are linearly independent. If $\sum_{j=0}^{k-1}c_jN^j\mathbf v=\mathbf 0$ and not all $c_j$ vanish, let $c_i$ be the first nonzero coefficient; applying $N^{k-1-i}$ kills every term with $j>i$, because $N^{k-1-i+j}\mathbf v=\mathbf 0$ for $j>i$, and leaves $c_iN^{k-1}\mathbf v\ne\mathbf 0$, a contradiction. The span $U$ of a chain is invariant under $N$, and with respect to the basis $\mathbf b_1=N^{k-1}\mathbf v,\ldots,\mathbf b_k=\mathbf v$, ordered from the bottom to the top, $N\mathbf b_1=\mathbf 0$ and $N\mathbf b_j=\mathbf b_{j-1}$ for $j\ge2$: the matrix of $N|_U$ is $N_k$. A basis of $V$ that is the union of disjoint Jordan chains therefore gives a block diagonal matrix with blocks $N_k$. The central result of this section is that such a basis always exists.

\subsection{Bases of Chains}

Let $N$ be a nilpotent operator on a space $V$ of dimension $n\ge1$, and let
\[
\begin{gathered}
K_j\coloneqq\ker N^j,\qquad\{\mathbf 0\}=K_0\subset K_1\subset\cdots\subset K_q=V,\\
d_j\coloneqq\dim K_j-\dim K_{j-1}\quad(j\ge1),
\end{gathered}
\]
where $q$ is the smallest exponent with $N^q=0$. By \cref{prop:kernel-chain}, the inclusions are strict up to $q$, so $d_1,\ldots,d_q\ge1$ and $d_j=0$ for $j>q$, and the increments decrease, $d_1\ge d_2\ge\cdots$; their sum is $\dim V=n$. A vector of $K_j$ that is not in $K_{j-1}$ is said to have \term{height} $j$: it is annihilated by $N^j$ but not by $N^{j-1}$. The top of a chain of length $k$ has height $k$, and $N$ lowers the height of a vector of height $j\ge1$ by exactly one. We say that vectors $\mathbf w_1,\ldots,\mathbf w_s$ are \term{independent modulo} a subspace $U$ if $\sum_ic_i\mathbf w_i\in U$ only for $c_1=\cdots=c_s=0$; for $U=\{\mathbf 0\}$ this is ordinary independence.

\begin{lemma}[Applying $N$ to independent vectors]
\label{lem:chain-step}
Let $j\ge1$. If $\mathbf w_1,\ldots,\mathbf w_s\in K_{j+1}$ are independent modulo $K_j$, then $N\mathbf w_1,\ldots,N\mathbf w_s\in K_j$ are independent modulo $K_{j-1}$. Every list of vectors of $K_j$ that is independent modulo $K_{j-1}$ can be extended by vectors of $K_j$ to a list that is independent modulo $K_{j-1}$ and has $d_j$ elements; such a list, together with a basis of $K_{j-1}$, is a basis of $K_j$.
\end{lemma}

\begin{proof}
If $\sum_ic_iN\mathbf w_i\in K_{j-1}$, then $N^{j-1}N\sum_ic_i\mathbf w_i=\mathbf 0$, so $\sum_ic_i\mathbf w_i\in K_j$, and all $c_i$ vanish; the vectors $N\mathbf w_i$ lie in $K_j$ because $N^jN\mathbf w_i=N^{j+1}\mathbf w_i=\mathbf 0$. For the second statement, let $\mathbf u_1,\ldots,\mathbf u_l$ be a basis of $K_{j-1}$ and $\mathbf w_1,\ldots,\mathbf w_s\in K_j$ independent modulo $K_{j-1}$. The list $\mathbf u_1,\ldots,\mathbf u_l,\mathbf w_1,\ldots,\mathbf w_s$ is independent: if $\sum a_m\mathbf u_m+\sum c_i\mathbf w_i=\mathbf 0$, then $\sum c_i\mathbf w_i\in K_{j-1}$, so all $c_i=0$, and then all $a_m=0$. Extend it to a basis of $K_j$ (\cref{thm:basis-extension}) by vectors $\mathbf w_{s+1},\ldots,\mathbf w_t$; the same argument, read backwards, shows that $\mathbf w_1,\ldots,\mathbf w_t$ are independent modulo $K_{j-1}$, and counting gives $t=\dim K_j-\dim K_{j-1}=d_j$.
\end{proof}

\begin{theorem}[Bases of Jordan chains]
\label{thm:chain-basis}
Let $N$ be a nilpotent operator on a space $V$ of finite dimension $n\ge1$. Then $V$ has a basis that is the union of disjoint Jordan chains for $N$. In every such basis, the number of chains of length at least $j$ is $d_j=\dim\ker N^j-\dim\ker N^{j-1}$, and the number of chains of length exactly $j$ is $d_j-d_{j+1}$.
\end{theorem}

\begin{proof}
We choose lists $L_q,L_{q-1},\ldots,L_1$, where $L_j$ consists of $d_j$ vectors of height $j$ that are independent modulo $K_{j-1}$. For $L_q$, extend the empty list by \cref{lem:chain-step}: $L_q$ is a list of $d_q$ vectors of $K_q=V$, independent modulo $K_{q-1}$; call it $B_q$ as well. If $L_{j+1}$ has been chosen, the vectors $N\mathbf w$, $\mathbf w\in L_{j+1}$, are $d_{j+1}$ vectors of $K_j$, independent modulo $K_{j-1}$, by the lemma; extend them by a list $B_j$ of $d_j-d_{j+1}$ further vectors of $K_j$ to a list $L_j$ of $d_j$ vectors that are independent modulo $K_{j-1}$. In particular $d_j-d_{j+1}\ge0$, which is \cref{prop:kernel-chain}\,(iv) again.

By construction, every vector of $L_i$ has the form $N^{j-i}\mathbf b$ for exactly one $j\ge i$ and one $\mathbf b\in B_j$, and for each $\mathbf b\in B_j$ the vectors $\mathbf b\in L_j$, $N\mathbf b\in L_{j-1}$, \ldots, $N^{j-1}\mathbf b\in L_1$ form a Jordan chain of length $j$, since $\mathbf b\in K_j$ gives $N^j\mathbf b=\mathbf 0$ and $N^{j-1}\mathbf b\in L_1$ is not zero. So the union of the lists $L_1,\ldots,L_q$ is a union of disjoint chains, one of length $j$ for each vector of $B_j$. It has $d_1+\cdots+d_q=n$ elements, and it is independent: in a vanishing combination of its vectors, let $j$ be the largest index for which a vector of $L_j$ has a nonzero coefficient. All vectors of $L_1,\ldots,L_{j-1}$ lie in $K_{j-1}$, so the part of the combination that comes from $L_j$ lies in $K_{j-1}$, and its coefficients vanish because $L_j$ is independent modulo $K_{j-1}$, a contradiction. Hence the union is a basis.

Finally, let any basis of disjoint chains be given, and let $c_k$ be the number of its chains of length $k$. In the bottom-to-top ordering the matrix of $N$ is block diagonal with $c_k$ blocks $N_k$ for each $k$. The block $N_k^j$ has rank $\max(k-j,0)$, since its nonzero columns are $\mathbf e_1,\ldots,\mathbf e_{k-j}$, so $\operatorname{rank}N^j=\sum_kc_k\max(k-j,0)$, and
\[
d_j=\operatorname{rank}N^{j-1}-\operatorname{rank}N^j=\sum_kc_k\bigl(\max(k-j+1,0)-\max(k-j,0)\bigr)=\sum_{k\ge j}c_k ,
\]
by \cref{thm:rank-nullity}. This is the number of chains of length at least $j$, and $c_j=d_j-d_{j+1}$.
\end{proof}

The proof chooses the chains from the top: the vectors of the greatest height first, then their images together with new vectors of the next height, and so on down to the eigenvectors. Choosing from the bottom fails in general, because an eigenvector need not be the bottom of a long chain: for $N_2\oplus N_1=\operatorname{diag}(N_2,N_1)$, the eigenvector $\mathbf e_3$ is the bottom of no chain of length $2$, since $\mathbf e_3$ is not in the image of $N$. \Cref{fig:dot-diagram} shows the structure that the theorem describes.

\begin{figure}
\sdfloat{%
\sdpanel{A}{\begin{tikzpicture}[x=0.9cm, y=0.62cm, font=\scriptsize]
% chains of lengths 4, 2, 2, 1 as rows; column j holds the vectors of height j
\foreach \y/\k in {4/4, 3/2, 2/2, 1/1} {
  \foreach \j in {1,...,\k} { \fill (\j,\y) circle (2.2pt); }
  \ifnum\k>1 \foreach \j in {2,...,\k} { \draw[-{Latex[length=1.5mm]}, line width=0.4pt, shorten >=4pt, shorten <=4pt] (\j,\y) -- (\j-1,\y); } \fi
  \draw[-{Latex[length=1.5mm]}, line width=0.4pt, shorten <=4pt] (1,\y) -- (0.35,\y);
}
\foreach \y in {1,...,4} { \node at (0.12,\y) {$\mathbf 0$}; }
\foreach \j/\d in {1/4, 2/3, 3/1, 4/1} { \node at (\j,0.15) {$d_{\j}=\d$}; }
\foreach \j in {1,...,4} { \node at (\j,4.85) {$j=\j$}; }
\node[anchor=west] at (4.35,4) {$k_1=4$};
\node[anchor=west] at (2.35,3) {$k_2=2$};
\node[anchor=west] at (2.35,2) {$k_3=2$};
\node[anchor=west] at (1.35,1) {$k_4=1$};
\end{tikzpicture}}\hfill
\sdpanel{B}{\begin{tikzpicture}[x=0.36cm, y=0.36cm, font=\scriptsize]
% the matrix diag(N_4, N_2, N_2, N_1) of order 9; block (r,c) cell centres at (c-0.5, 9.5-r)
\fill[sdfill] (0,9) rectangle (4,5);
\fill[sdfill] (4,5) rectangle (6,3);
\fill[sdfill] (6,3) rectangle (8,1);
\fill[sdfill] (8,1) rectangle (9,0);
\draw[line width=0.4pt] (0,0) rectangle (9,9);
\foreach \r/\c in {1/2, 2/3, 3/4, 5/6, 7/8} { \node at (\c-0.5, 9.5-\r) {$1$}; }
\end{tikzpicture}}}{The structure of a nilpotent operator (\cref{thm:chain-basis}). (A) A basis of four Jordan chains of lengths $4$, $2$, $2$ and $1$, one per row; the dots in column $j$ are the basis vectors of height $j$, and $N$ moves each dot one place to the left, the dots of the first column to $\mathbf 0$. Column $j$ contains $d_j=\dim\ker N^j-\dim\ker N^{j-1}$ dots, so the column lengths $4,3,1,1$ are read off from the kernels, and the row lengths from the column lengths. (B) The matrix of $N$ in this basis, ordered chain by chain from the bottom to the top: the block diagonal matrix $\operatorname{diag}(N_4,N_2,N_2,N_1)$ of order $9$, with the blocks shaded and the five ones shown; all other entries are zero.\label{fig:dot-diagram}}
\end{figure}

\begin{corollary}[Nilpotent matrices]
\label{cor:nilpotent-form}
Every nilpotent $N\in\mathbb F^{n\times n}$ is similar over $\mathbb F$ to a block diagonal matrix $\operatorname{diag}(N_{k_1},\ldots,N_{k_r})$ with $k_1\ge\cdots\ge k_r\ge1$, $k_1+\cdots+k_r=n$ and $r=\dim\mathcal N(N)$. The sizes $k_i$ are uniquely determined by $N$: the number of $i$ with $k_i\ge j$ is $\operatorname{rank}N^{j-1}-\operatorname{rank}N^j$. Two nilpotent matrices are similar if and only if $\operatorname{rank}N^j$ is the same for both, for every $j$.
\end{corollary}

\begin{proof}
Apply \cref{thm:chain-basis} to $\mathbf x\mapsto N\mathbf x$ on $\mathbb F^n$ and take the basis of chains, ordered by decreasing length and each chain from the bottom to the top, as the columns of $S$; then $S^{-1}NS$ is the block diagonal matrix, by \cref{sec:similarity} and the remarks after \cref{def:jordan-chain}. The number of chains is $d_1=\dim\mathcal N(N)$. The counts are those of the theorem, and they are expressed by ranks, which similar matrices share (\cref{prop:similarity-invariants}). If two nilpotent matrices have the same ranks of all powers, they have the same sizes $k_i$ and are similar to the same block diagonal matrix, hence to each other.
\end{proof}

The numbers $k_1\ge k_2\ge\cdots$ and $d_1\ge d_2\ge\cdots$ are the lengths of the rows and of the columns of the dot diagram, and each list determines the other. For example, the matrix
\[
N=\begin{bmatrix}0&1&0\\0&0&0\\0&1&0\end{bmatrix}
\]
has $N\mathbf x=(x_2,0,x_2)$, so $N^2=0$, $\operatorname{rank}N=1$ and $\operatorname{rank}N^2=0$; hence $d_1=2$ and $d_2=1$, and there are two chains, of lengths $2$ and $1$. Following the proof, a vector of height $2$ is any vector outside $\ker N=\{x_2=0\}$, say $\mathbf e_2$, with $N\mathbf e_2=(1,0,1)$; the second chain is a vector of $\ker N$ independent of $(1,0,1)$, say $\mathbf e_1$. With the columns $(1,0,1)$, $\mathbf e_2$ and $\mathbf e_1$,
\[
S=\begin{bmatrix}1&0&1\\0&1&0\\1&0&0\end{bmatrix},\qquad S^{-1}NS=\begin{bmatrix}0&1&0\\0&0&0\\0&0&0\end{bmatrix}=\operatorname{diag}(N_2,N_1),
\]
as the products $NS=[\,\mathbf 0\ (1,0,1)\ \mathbf 0\,]=S\operatorname{diag}(N_2,N_1)$ confirm.

\begin{note}{The Weyr and the Segre characteristic}
The two lists of the dot diagram have names. The block sizes $k_1\ge k_2\ge\cdots$, the lengths of the chains, form the \emph{Segre characteristic} of $N$, and the increments $d_1\ge d_2\ge\cdots$, the dimensions added by the successive kernels, form its \emph{Weyr characteristic}. Both are partitions of $n$, and they are \emph{conjugate} partitions: each is obtained from the other by exchanging the rows and the columns of the diagram. Ordering the basis of \cref{fig:dot-diagram} by columns instead of rows, so that the vectors of height $1$ come first, then those of height $2$, and so on, gives a block upper triangular matrix whose blocks just above the diagonal are identity matrices with rows of zeros appended, $[\,I_{d_j}\ 0\,]^{\mathsf T}$ in the block position $(j-1,j)$, and whose other blocks are zero; together with the eigenvalue on the diagonal it is the \emph{Weyr canonical form}~\cite[Sect.~3.4]{horn2013}.
\end{note}

\section{The Jordan Canonical Form}
\label{sec:jordan-form}

\subsection{Generalised Eigenspaces}

A general matrix is reduced to nilpotent pieces by shifting. If $\lambda$ is an eigenvalue of $A$, the matrix $A-\lambda I$ is singular, but it is not nilpotent unless $\lambda$ is the only eigenvalue; the part of the space on which it is nilpotent is found with the Fitting decomposition of \cref{thm:fitting}.

\begin{definition}[Generalised eigenspace]
\label{def:generalised-eigenspace}
Let $A\in\mathbb C^{n\times n}$ and $\lambda\in\mathbb C$. A nonzero vector $\mathbf x$ with $(A-\lambda I)^k\mathbf x=\mathbf 0$ for some $k\ge1$ is a \term{generalised eigenvector} of $A$ for $\lambda$, and
\[
G_\lambda(A)\coloneqq\ker(A-\lambda I)^n
\]
is the \term{generalised eigenspace} of $\lambda$.
\end{definition}

By \cref{prop:kernel-chain}\,(iii), the kernels of the powers of $A-\lambda I$ stop growing at an exponent at most $n$, so $G_\lambda(A)$ contains every generalised eigenvector, and it contains the eigenspace $E_\lambda=\ker(A-\lambda I)$. A generalised eigenvector for $\lambda$ exists only if $\lambda$ is an eigenvalue, since $(A-\lambda I)^k$ is invertible otherwise. For the Jordan block $J_n(\lambda)$, the generalised eigenspace of $\lambda$ is the whole space, while the eigenspace is a line.

\begin{theorem}[Generalised eigenspaces]
\label{thm:generalised-eigenspaces}
Let $A\in\mathbb C^{n\times n}$ have the pairwise different eigenvalues $\mu_1,\ldots,\mu_k$.
\begin{enumerate}
\item Each $G_{\mu_i}(A)$ is invariant under $A$, the restriction of $A-\mu_iI$ to it is nilpotent, and $\mu_i$ is the only eigenvalue of the restriction of $A$ to it.
\item $\dim G_{\mu_i}(A)=\operatorname{am}(\mu_i)$.
\item $\mathbb C^n=G_{\mu_1}(A)\oplus G_{\mu_2}(A)\oplus\cdots\oplus G_{\mu_k}(A)$.
\end{enumerate}
\end{theorem}

\begin{proof}
(i) $G_\mu=\ker p(A)$ for the polynomial $p(x)=(x-\mu)^n$, and kernels of polynomials in $A$ are invariant under $A$ (\cref{sec:invariant-quotient-dual}). The restriction $N$ of $A-\mu I$ satisfies $N^n=0$. If $A\mathbf x=\nu\mathbf x$ with $\mathbf x\in G_\mu$, $\mathbf x\ne\mathbf 0$, then $(A-\mu I)\mathbf x=(\nu-\mu)\mathbf x$, so $\mathbf 0=(A-\mu I)^n\mathbf x=(\nu-\mu)^n\mathbf x$ and $\nu=\mu$.

(ii) and (iii) We use induction on $n$; for $n=1$ both are clear. Let $\mu\coloneqq\mu_1$ and apply \cref{thm:fitting} to the operator $A-\mu I$ on $\mathbb C^n$ with $m=n$: $\mathbb C^n=G_\mu\oplus R$ with $R\coloneqq\operatorname{im}(A-\mu I)^n$, both subspaces are invariant under $A-\mu I$ and hence under $A=(A-\mu I)+\mu I$, and $A-\mu I$ restricted to $R$ is invertible, so that $\mu$ is not an eigenvalue of $A|_R$. In a basis made of a basis of $G_\mu$ and a basis of $R$, the matrix of $A$ is block diagonal (\cref{prop:block-triangular}), and by \cref{prop:charpoly-similar,thm:det-block-triangular}
\[
p_A=p_{A|_{G_\mu}}\,p_{A|_R}=(x-\mu)^{\dim G_\mu}\,p_{A|_R} ,
\]
where the second equation holds because, by (i), $\mu$ is the only root of the characteristic polynomial of $A|_{G_\mu}$. Since $\mu$ is not a root of $p_{A|_R}$, the multiplicity of the factor $x-\mu$ in $p_A$ is $\dim G_\mu$; this is (ii) for $\mu_1$. The roots of $p_{A|_R}$ are $\mu_2,\ldots,\mu_k$, with the same multiplicities as in $p_A$.

Next, $G_\nu(A)\subseteq R$ for every eigenvalue $\nu\ne\mu$. Let $\mathbf x\in G_\nu(A)$ and write $\mathbf x=\mathbf g+\mathbf r$ with $\mathbf g\in G_\mu$ and $\mathbf r\in R$. Both subspaces are invariant under $(A-\nu I)^n$, so $\mathbf 0=(A-\nu I)^n\mathbf g+(A-\nu I)^n\mathbf r$ is a decomposition of $\mathbf 0$ in the direct sum, and $(A-\nu I)^n\mathbf g=\mathbf 0$. But $\nu$ is not an eigenvalue of $A|_{G_\mu}$, by (i), so $(A-\nu I)|_{G_\mu}$ and its powers are injective, and $\mathbf g=\mathbf 0$. Hence $\mathbf x\in R$. Moreover $G_\nu(A)=G_\nu(A|_R)$: a vector of $R$ annihilated by a power of $A-\nu I$ is annihilated by its $(\dim R)$th power, by \cref{prop:kernel-chain}\,(iii) applied in $R$, and conversely. If $R=\{\mathbf 0\}$, then $\mu$ is the only eigenvalue and $\mathbb C^n=G_\mu$, and the proof is complete. Otherwise $\dim R<n$, since $\dim G_\mu\ge1$, and the induction hypothesis applies to $A|_R$, through its matrix in a basis of $R$, and gives $R=G_{\mu_2}(A)\oplus\cdots\oplus G_{\mu_k}(A)$ with $\dim G_{\mu_i}(A)=\operatorname{am}(\mu_i)$ for $i\ge2$. Together with $\mathbb C^n=G_{\mu_1}\oplus R$, this proves (ii) and (iii).
\end{proof}

Part (ii) is the statement announced at the end of \cref{sec:eigenspaces}: the kernels $\ker(A-\lambda I)^j$ grow until they reach the dimension $\operatorname{am}(\lambda)$. Since each step adds at least one dimension, they reach it no later than at $j=\operatorname{am}(\lambda)$, and $G_\lambda(A)=\ker(A-\lambda I)^{\operatorname{am}(\lambda)}$. Part (iii) generalises the decomposition $\mathbb F^n=E_{\mu_1}\oplus\cdots\oplus E_{\mu_k}$ of a diagonalisable matrix in \cref{thm:diagonalisable}\,(iv): every matrix splits the space into its generalised eigenspaces, and it is diagonalisable exactly when they are eigenspaces. The projections onto the summands are polynomials in $A$, as in \cref{lem:spectral-projections}; \cref{prop:spectral-projections-general} below constructs them. The decomposition can also be proved from the Cayley--Hamilton theorem: the factors $(x-\mu_i)^{\operatorname{am}(\mu_i)}$ of $p_A$ are pairwise coprime polynomials, and a polynomial identity $\sum_iu_i(x)\prod_{j\ne i}(x-\mu_j)^{\operatorname{am}(\mu_j)}=1$, obtained by the Euclidean algorithm for polynomials, provides the projections directly~\cite[Sect.~6.8]{hoffman1971}.

\subsection{The Theorem}

\begin{theorem}[Jordan canonical form]
\label{thm:jordan-form}
Every $A\in\mathbb C^{n\times n}$ is similar to a block diagonal matrix
\[
J=\operatorname{diag}\bigl(J_{k_1}(\lambda_1),J_{k_2}(\lambda_2),\ldots,J_{k_r}(\lambda_r)\bigr),
\]
in which $\lambda_1,\ldots,\lambda_r$ are eigenvalues of $A$, not necessarily different, and $k_1+\cdots+k_r=n$. The matrix $J$ is unique up to the order of the blocks: for every $\lambda\in\mathbb C$ and $j\ge1$, the number of blocks $J_k(\lambda)$ with $k\ge j$ is
\[
\operatorname{rank}(A-\lambda I)^{j-1}-\operatorname{rank}(A-\lambda I)^j .
\]
In particular the number of blocks for an eigenvalue $\lambda$ is $\operatorname{gm}(\lambda)$, and their sizes add up to $\operatorname{am}(\lambda)$. If $A$ is real and all its eigenvalues are real, the similarity can be chosen real.
\end{theorem}

\begin{proof}
For each eigenvalue $\mu$, the restriction $N_\mu$ of $A-\mu I$ to $G_\mu$ is nilpotent (\cref{thm:generalised-eigenspaces}\,(i)), and by \cref{thm:chain-basis} the space $G_\mu$ has a basis of disjoint Jordan chains for $N_\mu$. By \cref{thm:generalised-eigenspaces}\,(iii), the union of these bases over all eigenvalues is a basis of $\mathbb C^n$. Take its vectors, chain by chain and each chain from the bottom to the top, as the columns of $S$. For a chain $\mathbf b_1,\ldots,\mathbf b_k$ of $N_\mu$, ordered in this way, $A\mathbf b_1=\mu\mathbf b_1$ and $A\mathbf b_j=\mu\mathbf b_j+\mathbf b_{j-1}$ for $j\ge2$, which says that the corresponding columns of $S^{-1}AS$ form the block $J_k(\mu)$. So $S^{-1}AS$ is a block diagonal matrix of Jordan blocks.

For the uniqueness, let $J$ be any block diagonal matrix of Jordan blocks similar to $A$, and fix $\lambda$ and $j\ge1$. The matrix $(J-\lambda I)^j$ is block diagonal with the blocks $(J_{k_i}(\lambda_i)-\lambda I)^j$. If $\lambda_i\ne\lambda$, the block $J_{k_i}(\lambda_i)-\lambda I$ is triangular with the nonzero diagonal $\lambda_i-\lambda$, hence invertible, and its powers have rank $k_i$. If $\lambda_i=\lambda$, the block is $N_{k_i}$, whose $j$th power has rank $\max(k_i-j,0)$. Since ranks are similarity invariants,
\[
\operatorname{rank}(A-\lambda I)^j=\operatorname{rank}(J-\lambda I)^j=\sum_{i:\,\lambda_i\ne\lambda}k_i+\sum_{i:\,\lambda_i=\lambda}\max(k_i-j,0),
\]
and subtracting this from the same expression for $j-1$ gives the number of $i$ with $\lambda_i=\lambda$ and $k_i\ge j$, as in the proof of \cref{thm:chain-basis}. These numbers, for all $j$, determine how many blocks $J_k(\lambda)$ of each size occur, so every such $J$ has the same blocks. For $j=1$ the number of blocks is $n-\operatorname{rank}(A-\lambda I)=\operatorname{gm}(\lambda)$, and the sizes of the blocks for $\lambda$ add up to the dimension of the span of their chains, which is $\dim G_\lambda=\operatorname{am}(\lambda)$. If $A$ and its eigenvalues are real, the kernels $\ker(A-\mu I)^j$ have real bases, and the construction of \cref{thm:chain-basis} can be carried out in $\mathbb R^n$.
\end{proof}

The matrix $J$ is the \term{Jordan canonical form}, or Jordan normal form, of $A$, and the columns of $S$ form a \term{Jordan basis}. The blocks are usually grouped by eigenvalues, and within each group ordered by decreasing size. The Jordan basis is not unique, far from it: the choices in \cref{lem:chain-step} are arbitrary, and even for a single block $J_k(\lambda)$, every polynomial $c_0I+c_1N_k+\cdots$ with $c_0\ne0$ commutes with $J_k(\lambda)$ and transforms one Jordan basis into another. For an operator on a complex space of finite dimension the theorem gives, through any of its matrices, a basis in which its matrix is a Jordan matrix; differentiation on $\mathcal P_n$ consists of a single chain, $J_{n+1}(0)$.

\subsection{A Worked Example}

For the matrix
\[
A=\begin{bmatrix}2&0&3&-3\\1&2&-2&2\\1&0&0&2\\1&0&1&1\end{bmatrix},
\]
the characteristic polynomial is $p_A(x)=(x-2)^3(x+1)$, as one checks by evaluating $\det(xI-A)$ at four points, for instance $p_A(0)=-8$, $p_A(1)=-2$, $p_A(3)=4$ and $p_A(5)=162$, which determine the monic polynomial of degree $4$ (\cref{thm:interpolation}). So $\operatorname{am}(2)=3$ and $\operatorname{am}(-1)=1$. For the eigenvalue $2$,
\[
A-2I=\begin{bmatrix}0&0&3&-3\\1&0&-2&2\\1&0&-2&2\\1&0&1&-1\end{bmatrix},\qquad(A-2I)^2=\begin{bmatrix}0&0&-9&9\\0&0&9&-9\\0&0&9&-9\\0&0&0&0\end{bmatrix},
\]
with the ranks $2$ and $1$; the rank of $(A-2I)^3$ is $1$ again. So $\dim\ker(A-2I)=2$ and $\dim\ker(A-2I)^2=3=\operatorname{am}(2)$, the increments are $d_1=2$ and $d_2=1$, and by \cref{thm:jordan-form} the eigenvalue $2$ has two blocks, one of size at least $2$; since the sizes add up to $3$, they are $2$ and $1$. The eigenvalue $-1$ is simple and has a block of size $1$. Hence
\[
J=\operatorname{diag}\bigl(J_2(2),J_1(2),J_1(-1)\bigr)=\begin{bmatrix}2&1&0&0\\0&2&0&0\\0&0&2&0\\0&0&0&-1\end{bmatrix}.
\]
To find a Jordan basis, follow the proof of \cref{thm:chain-basis}. The first row of $A-2I$ gives $x_3=x_4$ on $\ker(A-2I)$, and then the second row gives $x_1=0$, so $\ker(A-2I)$ is spanned by $\mathbf e_2$ and $(0,0,1,1)$; the kernel of $(A-2I)^2$ is the space $x_3=x_4$, spanned by these two vectors and $\mathbf e_1$. The vector $\mathbf e_1$ has height $2$, and it is the top of the chain $(A-2I)\mathbf e_1=(0,1,1,1)$, $\mathbf e_1$. The chain of length $1$ is a vector of $\ker(A-2I)$ independent of $(0,1,1,1)$, say $\mathbf e_2$. For the eigenvalue $-1$, solving $(A+I)\mathbf x=\mathbf 0$ gives the eigenvector $(-1,1,1,0)$. With these four vectors as columns,
\[
S=\begin{bmatrix}0&1&0&-1\\1&0&1&1\\1&0&0&1\\1&0&0&0\end{bmatrix},\qquad S^{-1}=\begin{bmatrix}0&0&0&1\\1&0&1&-1\\0&1&-1&0\\0&0&1&-1\end{bmatrix},\qquad S^{-1}AS=J ,
\]
as multiplying out confirms. The minimal polynomial, by \cref{thm:jordan-minimal} below, is $(x-2)^2(x+1)$, of degree $3$, and the powers of $A$ are $A^m=SJ^mS^{-1}$, whose entries are combinations of $2^m$, $m2^{m-1}$ and $(-1)^m$ (\cref{sec:powers-again}).

\subsection{Consequences}

The rank formula of \cref{thm:jordan-form} answers the question of \cref{ch:linear-maps} which numbers decide similarity.

\begin{corollary}[Complete invariants of similarity]
\label{cor:similarity-complete}
Two matrices $A,B\in\mathbb C^{n\times n}$ are similar if and only if
\[
\operatorname{rank}(A-\lambda I)^j=\operatorname{rank}(B-\lambda I)^j
\]
for every $\lambda\in\mathbb C$ and every $j=1,\ldots,n$. It suffices to check the eigenvalues $\lambda$ of $A$ and of $B$.
\end{corollary}

\begin{proof}
If $B=P^{-1}AP$, then $(B-\lambda I)^j=P^{-1}(A-\lambda I)^jP$, and ranks are preserved. Conversely, the ranks determine the Jordan forms by \cref{thm:jordan-form}, since for $j>n$ they no longer change; if they agree, $A$ and $B$ are similar to the same Jordan matrix, and hence to each other (\cref{prop:similarity}). For $\lambda$ not an eigenvalue of either matrix, both ranks are $n$.
\end{proof}

The identity and the shear $K=J_2(1)$ of \cref{sec:similarity} differ already in $\operatorname{rank}(A-I)$, which is $0$ and $1$. The matrices $E_{11}$ and $N_2$ of \cref{prop:similarity}\,(iii) have the same rank $1$, but $\operatorname{rank}E_{11}^2=1$ and $\operatorname{rank}N_2^2=0$. The higher powers are needed in general. The nilpotent matrices $\operatorname{diag}(N_3,N_3,N_1)$ and $\operatorname{diag}(N_3,N_2,N_2)$ of order $7$ have the same characteristic polynomial $x^7$, the same minimal polynomial $x^3$ and the same geometric multiplicity $3$, and $\operatorname{rank}N=4$ for both; only $\operatorname{rank}N^2$, which is $2$ and $1$, tells them apart.

The Jordan form also explains the minimal polynomial. For an eigenvalue $\mu$, the size of the largest Jordan block for $\mu$ is called the \term{index} of $\mu$, $\operatorname{ind}(\mu)$.

\begin{theorem}[Minimal polynomial and Jordan blocks]
\label{thm:jordan-minimal}
Let $A\in\mathbb C^{n\times n}$ have the pairwise different eigenvalues $\mu_1,\ldots,\mu_k$. Then
\[
m_A(x)=\prod_{i=1}^k(x-\mu_i)^{\operatorname{ind}(\mu_i)} ,
\]
and $\operatorname{ind}(\mu)$ is the smallest exponent $q$ with $\ker(A-\mu I)^q=\ker(A-\mu I)^{q+1}$. It satisfies $1\le\operatorname{ind}(\mu)\le\operatorname{am}(\mu)-\operatorname{gm}(\mu)+1$. Consequently, $A$ is diagonalisable if and only if all its Jordan blocks have size $1$, if and only if $m_A$ has no repeated roots.
\end{theorem}

\begin{proof}
Let $A=SJS^{-1}$. For a polynomial $q$, $q(A)=Sq(J)S^{-1}$, and $q(J)$ is block diagonal with the blocks $q(J_{k_i}(\lambda_i))$, since sums and products of block diagonal matrices are computed block by block. So $q(A)=0$ if and only if $q(J_{k_i}(\lambda_i))=0$ for every $i$. The minimal polynomial of $J_k(\lambda)$ is $(x-\lambda)^k$ (\cref{sec:cayley-hamilton}), so by \cref{prop:minimal-polynomial}\,(i), $q(J_k(\lambda))=0$ if and only if $(x-\lambda)^k$ divides $q$. Hence $q(A)=0$ if and only if $q$ is divisible by $(x-\mu_i)^{\operatorname{ind}(\mu_i)}$ for every $i$, and the monic polynomial of least degree with this property is the product in the theorem. For the second statement, $\dim\ker(A-\mu I)^j=\dim\ker(J-\mu I)^j=\sum_{i:\,\lambda_i=\mu}\min(k_i,j)$, by the rank computation in the proof of \cref{thm:jordan-form} and \cref{thm:rank-nullity}; this sum increases with $j$ as long as $j<\max k_i$ and is constant from $j=\max k_i$ on. The $\operatorname{gm}(\mu)$ blocks for $\mu$ have sizes adding up to $\operatorname{am}(\mu)$, and each has size at least $1$, so the largest has size at most $\operatorname{am}(\mu)-(\operatorname{gm}(\mu)-1)$. Finally, $A$ is similar to a diagonal matrix if and only if its Jordan form is diagonal, by the uniqueness of the Jordan form, and this happens if and only if every index is $1$.
\end{proof}

The last statement is \cref{thm:diagonalisable-minimal}, now with an explanation: the exponent of $x-\mu$ in $m_A$ is the length of the longest chain for $\mu$. A matrix whose minimal and characteristic polynomials coincide is called \term{nonderogatory}; the Fibonacci matrix and every matrix with distinct eigenvalues are examples, and the identity $I_n$ with $n\ge2$ is not.

\begin{proposition}[Nonderogatory matrices]
\label{prop:nonderogatory}
For $A\in\mathbb C^{n\times n}$ the following are equivalent:
\begin{enumerate}
\item $m_A=p_A$;
\item every eigenvalue of $A$ has exactly one Jordan block, that is, $\operatorname{gm}(\lambda)=1$ for every eigenvalue;
\item $A$ is similar to the companion matrix $C_{p_A}$ of \cref{prop:companion}.
\end{enumerate}
\end{proposition}

\begin{proof}
(i)$\Leftrightarrow$(ii) By \cref{thm:jordan-minimal}, $m_A=p_A$ if and only if $\operatorname{ind}(\mu)=\operatorname{am}(\mu)$ for every eigenvalue, that is, if and only if the largest block for $\mu$ is the only one. (iii)$\Rightarrow$(ii) For every $\lambda$, the first $n-1$ columns of $C_q-\lambda I$ are $\mathbf e_{j+1}-\lambda\mathbf e_j$, $j=1,\ldots,n-1$; their entries in the rows $2,\ldots,n$ form an upper triangular matrix with ones on the diagonal, so they are independent, and $\operatorname{rank}(C_q-\lambda I)\ge n-1$. Hence every eigenvalue of $C_q$ has geometric multiplicity $1$, and similar matrices have the same geometric multiplicities. (ii)$\Rightarrow$(iii) The companion matrix $C_{p_A}$ has the characteristic polynomial $p_A$ (\cref{prop:companion}), so by the implication just proved, it has exactly one block for each eigenvalue $\mu$, of size $\operatorname{am}(\mu)$. If $A$ satisfies (ii), it has the same blocks, so $A$ and $C_{p_A}$ have the same Jordan form and are similar.
\end{proof}

\begin{proposition}[A matrix and its transpose]
\label{prop:transpose-similar}
Every $A\in\mathbb C^{n\times n}$ is similar to $A^{\mathsf T}$.
\end{proposition}

\begin{proof}
$(A^{\mathsf T}-\lambda I)^j=((A-\lambda I)^j)^{\mathsf T}$, and a matrix and its transpose have the same rank. By \cref{cor:similarity-complete}, $A$ and $A^{\mathsf T}$ are similar.
\end{proof}

The similarity can also be seen block by block: if $P$ is the permutation matrix that reverses the order of the coordinates, $P\mathbf e_j=\mathbf e_{k+1-j}$, then $PJ_k(\lambda)P=J_k(\lambda)^{\mathsf T}$, since $P$ turns the ones above the diagonal into ones below it. The proposition contrasts with the conjugate transpose. Since $A^{\mathsf H}=\bar A^{\mathsf T}$ is similar to $\bar A$, whose Jordan form is the conjugate of that of $A$, the matrices $A$ and $A^{\mathsf H}$ are similar if and only if every eigenvalue $\lambda$ of $A$ has the same Jordan blocks as $\bar\lambda$; a real matrix satisfies this, and $\operatorname{diag}(i,0)$ does not.

\subsection{The Real Jordan Form}

A real matrix with nonreal eigenvalues has a complex Jordan form, and, as with the Schur form, there is a real variant in which each conjugate pair of eigenvalues is represented by real blocks of even order. For $a,b\in\mathbb R$ with $b>0$ and $k\ge1$, let $C_k(a,b)$ be the real block matrix of order $2k$ with $k$ diagonal blocks $C_{a,b}$, the identity $I_2$ in the $k-1$ positions just above them and zeros elsewhere; for $k=2$,
\[
C_2(a,b)=\begin{bmatrix}C_{a,b}&I_2\\0&C_{a,b}\end{bmatrix}=\begin{bmatrix}a&-b&1&0\\b&a&0&1\\0&0&a&-b\\0&0&b&a\end{bmatrix}.
\]

\begin{theorem}[Real Jordan form]
\label{thm:real-jordan}
Every $A\in\mathbb R^{n\times n}$ is similar, by a real invertible matrix, to a block diagonal matrix whose blocks are Jordan blocks $J_k(\lambda)$ with real eigenvalues $\lambda$ and blocks $C_k(a,b)$ with $b>0$. For each real eigenvalue the blocks $J_k(\lambda)$ are those of the complex Jordan form of $A$, and for each pair $a\pm ib$ there is one block $C_k(a,b)$ for each block $J_k(a-ib)$ of the complex Jordan form. The form is unique up to the order of the blocks.
\end{theorem}

\begin{proof}
The complex Jordan form has, for each nonreal eigenvalue $\mu$, the same blocks as for $\bar\mu$, because $\operatorname{rank}(A-\bar\mu I)^j=\operatorname{rank}\overline{(A-\mu I)^j}=\operatorname{rank}(A-\mu I)^j$ for a real $A$. Moreover, the conjugates of the vectors of a Jordan chain for $\mu$ form a Jordan chain for $\bar\mu$, so the Jordan basis can be chosen with the chains for $\bar\mu$ conjugate to those for $\mu$. For a chain $\mathbf b_1,\ldots,\mathbf b_k$ for $\bar\mu=a-ib$, write $\mathbf b_j=\mathbf u_j+i\mathbf v_j$; separating real and imaginary parts in $A\mathbf b_j=\bar\mu\mathbf b_j+\mathbf b_{j-1}$, as in \cref{prop:real-2x2-form}, gives $A[\,\mathbf u_j\ \mathbf v_j\,]=[\,\mathbf u_j\ \mathbf v_j\,]C_{a,b}+[\,\mathbf u_{j-1}\ \mathbf v_{j-1}\,]$, which is the block $C_k(a,b)$ in the real basis $\mathbf u_1,\mathbf v_1,\ldots,\mathbf u_k,\mathbf v_k$. Replacing each pair of conjugate chains by these $2k$ real vectors does not change the complex span, as in the proof of \cref{thm:real-block-diagonal}, so the real vectors together with real chains for the real eigenvalues form a basis of $\mathbb R^n$. The uniqueness follows from that of the complex form, since $C_k(a,b)$ is similar over $\mathbb C$ to $\operatorname{diag}(J_k(a-ib),J_k(a+ib))$. Horn and Johnson treat the real Jordan form in the same way~\cite[Sect.~3.4]{horn2013}.
\end{proof}

\begin{note}{The rational canonical form}
The Jordan form needs the eigenvalues, which lie in $\mathbb C$ even for a matrix with rational entries. There is a canonical form that needs no extension of the field of scalars: every matrix over a field $\mathbb F$ is similar over $\mathbb F$ to a block diagonal matrix of companion matrices $C_{q_1},\ldots,C_{q_s}$ of monic polynomials with $q_{i+1}$ dividing $q_i$, and these polynomials, the \emph{invariant factors}, are uniquely determined; $q_1$ is the minimal polynomial and $q_1\cdots q_s$ the characteristic polynomial. Two matrices with entries in $\mathbb F$ are therefore similar over $\mathbb F$ if and only if they are similar over any larger field. This \emph{rational canonical form}, or Frobenius normal form, is computed by exact operations on polynomials, and its theory is that of modules over the ring of polynomials~\cite[Chap.~7]{hoffman1971}. It is the right tool for questions of exact algebra; for computation with numbers, it shares the defects of the Jordan form discussed in \cref{sec:jordan-unstable}.
\end{note}

\section{Powers, Again}
\label{sec:powers-again}

\subsection{Powers of a Jordan Block}

\Cref{sec:matrix-powers} computed the powers of diagonalisable matrices from their eigenvalues, and \cref{sec:cayley-hamilton} proved that $A^k\to0$ when $\rho(A)<1$ without describing $A^k$. With the Jordan form, $A^m=SJ^mS^{-1}$, and $J^m$ is block diagonal with the blocks $J_k(\lambda)^m$, so everything comes down to the powers of a single Jordan block.

\begin{proposition}[Powers of a Jordan block]
\label{prop:jordan-block-powers}
For $\lambda\in\mathbb C$, $k\ge1$ and $m\in\mathbb N$,
\[
J_k(\lambda)^m=\sum_{j=0}^{k-1}\binom mj\lambda^{m-j}N_k^j=\begin{bmatrix}\lambda^m&\binom m1\lambda^{m-1}&\binom m2\lambda^{m-2}&\cdots&\binom m{k-1}\lambda^{m-k+1}\\&\lambda^m&\binom m1\lambda^{m-1}&\cdots&\binom m{k-2}\lambda^{m-k+2}\\&&\ddots&\ddots&\vdots\\&&&\lambda^m&\binom m1\lambda^{m-1}\\&&&&\lambda^m\end{bmatrix},
\]
where a term with $j>m$ is zero, because the binomial coefficient $m^{\underline j}/j!$ vanishes for $j>m$. The entry in position $(i,i+j)$ is $(m^{\underline j}/j!)\lambda^{m-j}$.
\end{proposition}

\begin{proof}
$J_k(\lambda)=\lambda I+N_k$, and $\lambda I$ commutes with $N_k$. The binomial theorem (\cref{thm:binomial}) holds for two commuting matrices, since its proof uses only the distributive law and the commutativity of the two summands, so
\[
(\lambda I+N_k)^m=\sum_{j=0}^m\binom mj\lambda^{m-j}N_k^j ,
\]
and the terms with $j\ge k$ vanish because $N_k^j=0$. The matrix $N_k^j$ has ones in the positions $(i,i+j)$ and zeros elsewhere.
\end{proof}

The binomial coefficient $m^{\underline j}/j!=m(m-1)\cdots(m-j+1)/j!$ is a polynomial in $m$ of degree $j$, so each entry of $J_k(\lambda)^m$ is a polynomial in $m$ of degree less than $k$, times $\lambda^m$, when $\lambda\ne0$: the entry is $(m^{\underline j}\lambda^{-j}/j!)\,\lambda^m$. For $\lambda=0$ the block is nilpotent, and $N_k^m=0$ for $m\ge k$. For $k=2$ and $\lambda=4/5$ the proposition gives the powers of the matrix $A_2$ of \cref{sec:matrix-powers}, which is $(4/5)I+2N_2$ and therefore similar to $J_2(4/5)$: $A_2^m=(4/5)^mI+2m(4/5)^{m-1}N_2$.

\subsection{The Entries of \texorpdfstring{$A^m$}{Am} and Linear Recurrences}

\begin{theorem}[Entries of the powers of a matrix]
\label{thm:power-entries}
Let $A\in\mathbb C^{n\times n}$, let $\mu_1,\ldots,\mu_l$ be its pairwise different nonzero eigenvalues, and let $m_0\coloneqq\operatorname{ind}(0)$ if $0$ is an eigenvalue and $m_0\coloneqq0$ otherwise. For each position $(i,j)$ there are polynomials $P_1,\ldots,P_l$, with $\deg P_t<\operatorname{ind}(\mu_t)$, such that
\[
(A^m)_{ij}=\sum_{t=1}^lP_t(m)\,\mu_t^m\qquad\text{for all }m\ge m_0 .
\]
\end{theorem}

\begin{proof}
Let $A=SJS^{-1}$ with the Jordan form $J$. For $m\ge m_0$, every block of $J^m$ for the eigenvalue $0$ is zero, since its size is at most $\operatorname{ind}(0)$, and by \cref{prop:jordan-block-powers} every entry of a block $J_k(\mu_t)^m$ has the form $c\,(m^{\underline j}/j!)\,\mu_t^{-j}\cdot\mu_t^m$ with $j<k\le\operatorname{ind}(\mu_t)$ and $c\in\{0,1\}$. The entry $(A^m)_{ij}=\sum_{a,b}s_{ia}(J^m)_{ab}(S^{-1})_{bj}$ is a fixed linear combination of the entries of $J^m$, and collecting the terms with the same $\mu_t$ gives the polynomials $P_t$.
\end{proof}

This is the form that \cref{sec:cayley-hamilton} derived from \cref{prop:linear-recurrence-k}, with a sharper bound on the degrees: the degree of $P_t$ is less than the index of $\mu_t$, not only less than its algebraic multiplicity. For a diagonalisable matrix all indices are $1$ and the $P_t$ are constants. The argument can be reversed, and it gives a complete proof of \cref{prop:linear-recurrence-k}, which \cref{ch:counting} proved only in outline.

\begin{corollary}[Linear recurrences with constant coefficients]
\label{cor:linear-recurrence}
Let $c_1,\ldots,c_k\in\mathbb C$ with $c_k\ne0$, and let the polynomial $q(x)=x^k-c_1x^{k-1}-\cdots-c_k$ have the pairwise different roots $r_1,\ldots,r_s$ with the multiplicities $m_1,\ldots,m_s$. Then the sequences $(a_n)_{n\ge0}$ with
\[
a_n=c_1a_{n-1}+c_2a_{n-2}+\cdots+c_ka_{n-k}\qquad(n\ge k)
\]
are exactly the sequences $a_n=\sum_{i=1}^sP_i(n)r_i^n$ with polynomials $P_i$ of degree less than $m_i$, and the $k$ coefficients of $P_1,\ldots,P_s$ are uniquely determined by $a_0,\ldots,a_{k-1}$.
\end{corollary}

\begin{proof}
Let $\mathcal S$ be the set of solutions, a subspace of the space of all sequences. A solution is determined by $a_0,\ldots,a_{k-1}$, which can be prescribed arbitrarily, and the map $(a_n)\mapsto(a_0,\ldots,a_{k-1})$ is linear and bijective from $\mathcal S$ onto $\mathbb C^k$; so $\dim\mathcal S=k$. Write $q(x)=x^k+\gamma_{k-1}x^{k-1}+\cdots+\gamma_0$ with $\gamma_{k-j}=-c_j$, and let $C\coloneqq C_q^{\mathsf T}$ be the transposed companion matrix of \cref{prop:companion}. As observed after that proposition, the vectors $\mathbf s_n\coloneqq(a_n,\ldots,a_{n+k-1})$ of a solution satisfy $\mathbf s_{n+1}=C\mathbf s_n$, so $\mathbf s_n=C^n\mathbf s_0$ and $a_n=\mathbf e_1^{\mathsf T}C^n\mathbf s_0$ for all $n\ge0$. The eigenvalues of $C$ are the roots of $p_C=p_{C_q}=q$ (\cref{prop:companion,prop:transpose-similar}); they are nonzero, since $q(0)=-c_k\ne0$, so $m_0=0$ in \cref{thm:power-entries}. By \cref{prop:nonderogatory}, $C$ is nonderogatory, so $\operatorname{ind}(r_i)=\operatorname{am}(r_i)=m_i$. By \cref{thm:power-entries}, each entry of $C^n$, and hence $a_n$, has the form $\sum_iP_i(n)r_i^n$ with $\deg P_i<m_i$, for all $n\ge0$. So $\mathcal S$ is contained in the span $\mathcal T$ of the $k$ sequences $(n^jr_i^n)_{n\ge0}$, $0\le j<m_i$. Since $\dim\mathcal S=k$, the span $\mathcal T$ has dimension at least $k$, so these $k$ sequences are linearly independent and span a space of dimension exactly $k$, which contains $\mathcal S$ and therefore equals it (\cref{prop:subspace-dimension}). Hence every sequence $\sum_iP_i(n)r_i^n$ with $\deg P_i<m_i$ is a solution, and since the $k$ sequences are a basis of $\mathcal S$, the coefficients of a solution are determined by it, and so by its first $k$ terms.
\end{proof}

The proof never verifies directly that $n^jr_i^n$ solves the recurrence; it follows from counting dimensions, as in \cref{sec:rank-nullity-applications}. For the recurrence $a_n=4a_{n-1}-4a_{n-2}$, with $q(x)=(x-2)^2$, the solutions are $a_n=(\alpha+\beta n)2^n$; the factor $n$ is the shadow of the Jordan block $J_2(2)$ of the companion matrix, whose powers carry the entry $m2^{m-1}$. \Cref{ch:difference-equations} returns to linear recurrences, their generating functions and the stability of their numerical solution, and \cref{ch:linear-higher-order} meets the same phenomenon for differential equations with constant coefficients, where a repeated root of the characteristic equation produces solutions $t^je^{\lambda t}$.

\subsection{Decay, Boundedness and Growth}

The explicit powers turn the qualitative question of \cref{sec:matrix-powers}, whether $A^m$ tends to zero, into a precise description of the size of $A^m$.

\begin{theorem}[The size of the powers]
\label{thm:powers-size}
Let $A\in\mathbb C^{n\times n}$.
\begin{enumerate}
\item $A^m\to0$ as $m\to\infty$ if and only if $\rho(A)<1$.
\item The powers $A^m$, $m\in\mathbb N$, are bounded if and only if $\rho(A)\le1$ and every eigenvalue of modulus $1$ has index $1$, that is, $\operatorname{gm}(\lambda)=\operatorname{am}(\lambda)$ for every eigenvalue with $|\lambda|=1$.
\item Let $\rho(A)>0$, and let $s$ be the largest index of an eigenvalue of modulus $\rho(A)$. Then there are constants $0<c\le C$ such that
\[
c\,m^{s-1}\rho(A)^m\le\|A^m\|_\infty\le C\,m^{s-1}\rho(A)^m\qquad\text{for all }m\ge2s .
\]
\end{enumerate}
\end{theorem}

\begin{proof}
Let $A=SJS^{-1}$ and $\kappa\coloneqq\|S\|_\infty\|S^{-1}\|_\infty$. From $A^m=SJ^mS^{-1}$ and $J^m=S^{-1}A^mS$, $\|J^m\|_\infty/\kappa\le\|A^m\|_\infty\le\kappa\|J^m\|_\infty$, so it suffices to prove the statements for $J$, whose powers are block diagonal. The norm of a block diagonal matrix is the largest norm of a block, by \cref{thm:norm-formulas}, and by \cref{prop:jordan-block-powers} the largest row sum of the moduli of $J_k(\mu)^m$ is the first one,
\[
\|J_k(\mu)^m\|_\infty=\sum_{j=0}^{k-1}\binom mj|\mu|^{m-j}.
\]
(i) If $|\mu|<1$, each term $(m^{\underline j}/j!)|\mu|^{m-j}\le m^j|\mu|^{m-j}$ tends to zero, since a power of $|\mu|<1$ beats every polynomial (\cref{thm:standard-limits}); for $\mu=0$ the terms vanish for $m\ge k$. So $J^m\to0$ if $\rho(A)<1$. Conversely, if $A^m\to0$, then $\rho(A)<1$ by \cref{thm:powers-to-zero}. (ii) If $|\mu|>1$, the diagonal entry $\mu^m$ is unbounded, and if $|\mu|=1$ and $k\ge2$, the entry $m\mu^{m-1}$ in position $(1,2)$ has modulus $m$. Conversely, if $\rho(A)\le1$ and the blocks for the eigenvalues of modulus $1$ have size $1$, these blocks are numbers of modulus $1$, and the others tend to zero by (i), so $J^m$ is bounded. (iii) Let $\rho\coloneqq\rho(A)$. For a block with $|\mu|=\rho$ and $k\le s$, and $m\ge1$,
\[
\sum_{j=0}^{k-1}\binom mj\rho^{m-j}\le\rho^m\sum_{j=0}^{s-1}m^j\rho^{-j}\le s\max(1,\rho^{1-s})\,m^{s-1}\rho^m .
\]
For a block with $0<|\mu|<\rho$, each term $(m^{\underline j}/j!)|\mu|^{m-j}$, divided by $m^{s-1}\rho^m$, is at most $m^{j}|\mu|^{-j}(|\mu|/\rho)^m$, which tends to zero and is therefore bounded; blocks with $\mu=0$ vanish for $m\ge n$ and are bounded before. This gives the upper bound for $J$. For the lower bound, take a block $J_s(\mu)$ with $|\mu|=\rho$. Its entry in position $(1,s)$ is $(m^{\underline{s-1}}/(s-1)!)\,\mu^{m-s+1}$, and for $m\ge2s$ every factor of $m(m-1)\cdots(m-s+2)$ is at least $m/2$, so
\[
\|J^m\|_\infty\ge\binom m{s-1}\rho^{m-s+1}\ge\frac{\rho^{1-s}}{2^{s-1}(s-1)!}\,m^{s-1}\rho^m .
\]
Dividing by $\kappa$ gives the lower bound for $A$.
\end{proof}

Gelfand's formula follows from part (iii) as well, since $(cm^{s-1})^{1/m}\to1$ and $(Cm^{s-1})^{1/m}\to1$, but part (iii) says more: the norms of the powers grow or decay like $\rho^m$ times a polynomial whose degree is one less than the size of the largest Jordan block on the circle $|\lambda|=\rho$. Part (ii) is the condition under which an iteration $\mathbf x_{m+1}=A\mathbf x_m$ is \term{stable}, in the sense that its iterates stay bounded for every starting vector; a matrix with a Jordan block of size $2$ for the eigenvalue $1$, such as the shear $K$, has powers that grow linearly although its spectral radius is $1$. \Cref{ch:difference-equations} uses this condition for recurrences, and \cref{ch:linear-ode-systems} meets its continuous counterpart for systems of differential equations.

\Cref{fig:neumann} showed the norms of the powers of $A_2=(4/5)I+2N_2$ growing to $4.5$ before they decay, and the theorem explains the shape. By \cref{prop:jordan-block-powers}, $\|A_2^m\|_\infty=(4/5)^m+2m(4/5)^{m-1}=(1+5m/2)(4/5)^m$, the product of a polynomial of degree $s-1=1$, which grows, and a geometric sequence, which decays. The ratio of consecutive values is $(4/5)(7/2+5m/2)/(1+5m/2)$, which exceeds $1$ exactly when $m<18/5$, so the largest value occurs at $m=4$ and equals $11\cdot(4/5)^4=4.5056$. The larger the Jordan block, the higher and later the peak. \Cref{fig:jordan-powers} shows the norms $\|J_k(0.9)^m\|_\infty$ for $k=1,2,3,4$. For $k=1$ they decay from the start; for $k=4$ they rise to $261$ at $m=28$ and are still about $6.06$ at $m=100$, although every eigenvalue has modulus $0.9$; they fall below $1$ only at $m=123$.

\begin{figure}
\sdfloat{%
\begin{tikzpicture}
\begin{axis}[sdplot, width=0.62\textwidth, height=0.42\textwidth, ymode=log,
  xmin=0, xmax=200, ymin=1e-10, ymax=3e3, xtick={0,50,100,150,200},
  ytick={1e-10,1e-8,1e-6,1e-4,1e-2,1,1e2}, xlabel={$m$}, legend pos=south west]
\addplot[black, line width=0.7pt, samples at={0,2,...,200}] {0.9^x};
\addlegendentry{$k=1$}
\addplot[black, dashed, line width=0.7pt, samples at={0,1,...,200}] {0.9^x+x*0.9^(x-1)};
\addlegendentry{$k=2$}
\addplot[black, densely dotted, line width=0.8pt, samples at={0,1,...,200}] {0.9^x+x*0.9^(x-1)+x*(x-1)/2*0.9^(x-2)};
\addlegendentry{$k=3$}
\addplot[black, dashdotted, line width=0.7pt, samples at={0,1,...,200}] {0.9^x+x*0.9^(x-1)+x*(x-1)/2*0.9^(x-2)+x*(x-1)*(x-2)/6*0.9^(x-3)};
\addlegendentry{$k=4$}
\end{axis}
\end{tikzpicture}}{The norms $\|J_k(0.9)^m\|_\infty=\sum_{j<k}(m^{\underline j}/j!)\,0.9^{m-j}$ of the powers of Jordan blocks with the eigenvalue $0.9$, for $k=1,2,3,4$ (\cref{prop:jordan-block-powers}); the values at integer $m$ are joined by lines. All four tend to zero, but the block of size $k$ first grows like $m^{k-1}$: the maxima are $4.26$ at $m=9$, $31.5$ at $m=19$ and $261$ at $m=28$ for $k=2,3,4$.\label{fig:jordan-powers}}
\end{figure}

\section{Why the Jordan Form Is Not Computed}
\label{sec:jordan-unstable}

\subsection{A Discontinuous Form}

The Jordan form is a complete answer to the question of similarity, but it is not a continuous function of the matrix, and this makes it unsuitable for computation with numbers that carry rounding errors. The example of \cref{sec:computing-eigenvalues} shows it. For $\varepsilon>0$, the matrix $J_\varepsilon$ obtained from $N_n$ by adding $\varepsilon$ in the lower left corner has the characteristic polynomial $x^n-\varepsilon$ and the $n$ different eigenvalues $\varepsilon^{1/n}e^{2\pi ij/n}$. Its Jordan form is therefore the diagonal matrix of these eigenvalues, which tends to the zero matrix as $\varepsilon\to0$, while $J_\varepsilon$ tends to $N_n$, whose Jordan form is $N_n$ itself. In order $2$ the same happens without any spreading of eigenvalues: the matrix with the rows $(0,1)$ and $(0,\varepsilon)$ has the eigenvalues $0$ and $\varepsilon$ and the Jordan form $\operatorname{diag}(0,\varepsilon)$ for $\varepsilon\ne0$, which tends to $0$, and the limit $N_2$ is not similar to $0$. Every neighbourhood of a defective matrix contains diagonalisable matrices (\cref{thm:diagonalisable-dense}), so the Jordan form jumps, in every neighbourhood of a defective matrix, between different block structures.

The block structure is decided by ranks, $\operatorname{rank}(A-\lambda I)^j$ for an eigenvalue $\lambda$, and ranks are discontinuous as well: a small perturbation can increase the rank, never decrease it near a given matrix, because a nonzero minor stays nonzero. In floating-point arithmetic the matrix $A-\lambda I$ is not even known exactly, since $\lambda$ is a computed eigenvalue, and deciding its rank is the problem of the numerical rank of \cref{sec:low-rank}: the matrix is within the distance $\sigma_{r+1}$ of a matrix of rank $r$ (\cref{thm:eckart-young}), and when the singular values decrease gradually there is no meaningful rank to report. For $J_\varepsilon$ with a tiny $\varepsilon$, the matrix $J_\varepsilon-0\cdot I$ has the singular values $1$, $n-1$ times, and $\varepsilon$ (\cref{sec:svd-perturbation}), and whether its rank is $n-1$ or $n$ depends on whether $\varepsilon$ is regarded as an error.

\subsection{Ill-Conditioned Similarities}

The similarity that achieves the Jordan form, or a diagonal form near a defective matrix, can be arbitrarily ill-conditioned. For $\delta\ne0$ let
\[
A_\delta\coloneqq\begin{bmatrix}1&1\\0&1+\delta\end{bmatrix},
\]
with the distinct eigenvalues $1$ and $1+\delta$ and the eigenvectors $\mathbf e_1$ and $(1,\delta)$, which make the angle $\theta$ with $\cos\theta=1/\sqrt{1+\delta^2}$; as $\delta\to0$ the matrix tends to the shear $K=J_2(1)$, and the two eigenvector lines merge.

\begin{proposition}[Condition of a diagonalising matrix]
\label{prop:diagonalising-condition}
Every invertible $S$ for which $S^{-1}A_\delta S$ is diagonal satisfies
\[
\kappa_2(S)\ge\frac{\sqrt{1+\delta^2}+1}{|\delta|}>\frac2{|\delta|},
\]
and equality holds when the columns of $S$ have equal lengths.
\end{proposition}

\begin{proof}
By \cref{thm:diagonalisable}, the columns of $S$ are eigenvectors, so they are nonzero multiples of $\mathbf e_1$ and of $\mathbf v\coloneqq(1,\delta)/\sqrt{1+\delta^2}$, in some order. The condition number $\kappa_2(S)=\sigma_1/\sigma_2$ does not change when the columns are exchanged, when $S$ is multiplied by a nonzero number, or when a column is multiplied by a number of modulus $1$ (\cref{prop:singular-eigen}\,(ii), with a permutation or a diagonal unitary matrix on the right). So we may assume $S=[\,\mathbf e_1\ t\mathbf v\,]$ with $t>0$. With $c\coloneqq\langle\mathbf e_1,\mathbf v\rangle=\cos\theta$,
\[
S^{\mathsf T}S=\begin{bmatrix}1&tc\\tc&t^2\end{bmatrix},
\]
whose eigenvalues $\sigma_1^2$ and $\sigma_2^2$ have the sum $1+t^2$ and the product $t^2(1-c^2)=t^2\sin^2\theta$. Hence
\[
\kappa_2(S)+\frac1{\kappa_2(S)}=\frac{\sigma_1^2+\sigma_2^2}{\sigma_1\sigma_2}=\frac{1+t^2}{t\sin\theta}\ge\frac2{\sin\theta},
\]
with equality for $t=1$. The function $\kappa+1/\kappa$ is increasing for $\kappa\ge1$, and $\kappa+1/\kappa=2/\sin\theta$ has the solution $\kappa=(1+\cos\theta)/\sin\theta\ge1$, so $\kappa_2(S)\ge(1+\cos\theta)/\sin\theta$. With $\cos\theta=1/\sqrt{1+\delta^2}$ and $\sin\theta=|\delta|/\sqrt{1+\delta^2}$ this is the bound.
\end{proof}

For $\delta=10^{-8}$ every diagonalising matrix has a condition number above $2\times10^8$, although $A_\delta$ is a perfectly harmless upper triangular matrix of norm about $1.6$, whose Schur form is the matrix itself with $Q=I$. By the Bauer--Fike theorem (\cref{thm:bauer-fike}), such a condition number permits large movements of the eigenvalues under perturbations, and any computation that passes through $S$ and $S^{-1}$ can amplify the rounding errors by it; \cref{sec:computing-matrix-functions} measures this for the computation of $e^{A}$. At $\delta=0$ there is no diagonalising matrix at all, and a Jordan basis exists, but it cannot be found as a limit of the eigenvector bases of the nearby diagonalisable matrices, which degenerate.

\subsection{The Schur Form in Contrast}

A Schur decomposition avoids both difficulties. The transformation is unitary, so it has the condition number $1$ and does not magnify errors (\cref{sec:unitary-matrices}), and the computed decomposition of a matrix $A$ is the exact decomposition of a nearby matrix: the algorithms of Volume~VIII, which reduce $A$ to Hessenberg form and then apply the QR iteration, return a unitary $Q$ and a triangular $T$ with $A+E=QTQ^{\mathsf H}$ for a perturbation $E$ whose norm is a modest multiple of $u\|A\|_2$, up to rounding errors in the orthogonality of the computed $Q$~\cite[Sect.~7.5]{golub2013}. The eigenvalues on the diagonal of $T$ are those of $A+E$, and so they inherit the sensitivity of the eigenvalues of $A$; but the decomposition itself is computed as accurately as the data permit, and the triangle $T$ carries in its off-diagonal part the departure from normality of \cref{prop:schur-inequality}, which the Jordan form would replace by ones.

We tested this with the routines \code{linalg.schur} and \code{linalg.eig} of the stardust library, on the machine and with the version of the library described in \cref{sec:systems-practice}. For the defective matrix \eqref{eq:schur-example}, stored exactly in binary64, \code{schur} returned a triangular $T$ with the diagonal entries $2+1.65\times10^{-8}$ and $2-1.65\times10^{-8}$ and the entry $t_{12}=2.0000$ to all digits shown, and an orthogonal $Z$ whose first column agrees with $(1,-1)/\sqrt2$ to about $6\times10^{-9}$: the eigenvalues have moved by about $\sqrt u$, as \cref{tab:jordan-sensitivity} leads one to expect for a block of size $2$, while the off-diagonal entry, the departure from normality, is accurate. \Cref{tab:schur-defective} repeats the test with the matrices $HJ_n(1)H$ of \cref{sec:computing-eigenvalues}, formed in binary64 from the Householder reflection $H=I-2\mathbf v\mathbf v^{\mathsf T}/\mathbf v^{\mathsf T}\mathbf v$ with $\mathbf v=(1,2,\ldots,n)$, whose entries were rounded to binary64, by two matrix products from left to right. The returned $T$ is real and quasi-triangular, as \cref{thm:real-schur} describes; its eigenvalues, computed from the blocks by the formula \eqref{eq:eigen-2x2}, have the same largest distance from $1$ as the values that \code{eigvals} returned for \cref{tab:jordan-sensitivity}, of the order of $u^{1/n}$ or smaller. Yet the computed $Z$ is orthogonal to within $5\times10^{-15}$, a few dozen units of $u$, and $ZTZ^{\mathsf T}$ reproduces the matrix to the same accuracy, relative to its norm. The eigenvector matrix $V$ returned by \code{eig}, with columns of unit length, has a condition number, computed by \code{linalg.cond}, that grows from $8.5\times10^7$ to values above $10^{17}$, which means that for $n=8$ and $n=16$ the computed eigenvectors are linearly dependent to working precision.

\begin{table}
\sdfloat{%
\begin{sdtable}{colspec={rllll}, row{1}={font=\small\bfseries}}
\toprule
$n$ & $\max|\hat\lambda-1|$ & $\|Z^{\mathsf T}Z-I\|_2$ & $\|A-ZTZ^{\mathsf T}\|_2/\|A\|_2$ & $\kappa_2(V)$ \\
\midrule
$2$ & $1.2\times10^{-8}$ & $2.4\times10^{-16}$ & $3.1\times10^{-16}$ & $8.5\times10^{7}$ \\
$4$ & $1.2\times10^{-4}$ & $3.0\times10^{-15}$ & $3.3\times10^{-15}$ & $6.4\times10^{11}$ \\
$8$ & $2.9\times10^{-3}$ & $2.5\times10^{-15}$ & $2.1\times10^{-15}$ & $3.3\times10^{17}$ \\
$16$ & $7.6\times10^{-2}$ & $4.7\times10^{-15}$ & $5.3\times10^{-15}$ & $1.8\times10^{18}$ \\
\bottomrule
\end{sdtable}}{The Schur decomposition $A=ZTZ^{\mathsf T}$ computed by \code{linalg.schur} for the matrices $A=HJ_n(1)H$ formed in binary64, which are similar to the Jordan block $J_n(1)$: the largest distance from $1$ of the eigenvalues $\hat\lambda$ of the computed $T$, the departure of the computed $Z$ from orthogonality, the relative residual, and the condition number of the eigenvector matrix $V$ computed by \code{linalg.eig}. All norms were computed by \code{linalg.matrix\_norm} and \code{linalg.cond} in binary64.\label{tab:schur-defective}}
\end{table}

\Cref{tab:schur-jordan} summarises the comparison. Both forms are canonical in a sense, but in different senses: the Jordan form is a complete invariant of similarity and is unique, while the Schur form is far from unique and its off-diagonal part depends on choices, but it is reached by transformations that do not amplify errors. The Jordan form is the language in which theorems about powers, functions and differential equations are stated and proved, and the Schur form the one in which they are computed. Algorithms that determine a Jordan structure numerically exist, but they must decide ranks with tolerances and can only find the structure of some matrix near the given one; Golub and Wilkinson discuss what can and cannot be achieved~\cite{golubwilkinson1976}.

\begin{table}
\sdfloat{%
\begin{sdtable}{colspec={lll}, row{1}={font=\small\bfseries}}
\toprule
 & Schur form $A=QTQ^{\mathsf H}$ & Jordan form $A=SJS^{-1}$ \\
\midrule
exists for & every $A\in\mathbb C^{n\times n}$ & every $A\in\mathbb C^{n\times n}$ \\
transformation & unitary, $\kappa_2(Q)=1$ & any invertible $S$, $\kappa_2(S)$ unbounded \\
result & upper triangular & block diagonal, Jordan blocks \\
uniqueness & no; $\|N\|_{\mathrm F}$ and $\operatorname{spec}A$ fixed & up to the order of the blocks \\
decides similarity & no & yes (\cref{cor:similarity-complete}) \\
real version & quasi-triangular (\cref{thm:real-schur}) & blocks $C_k(a,b)$ (\cref{thm:real-jordan}) \\
normal $A$ & diagonal & diagonal \\
in floating point & backward stable algorithms & discontinuous at defective $A$ \\
used for & computation & theory \\
\bottomrule
\end{sdtable}}{The Schur form and the Jordan form compared (\cref{sec:schur,sec:jordan-form,sec:jordan-unstable}); $N$ is the strictly upper triangular part of $T$.\label{tab:schur-jordan}}
\end{table}

\begin{note}{Pseudospectra}
The eigenvalues of a matrix far from normal say little about the matrix in the presence of perturbations, and Trefethen and Embree~\cite{trefethen2005} advocate a set that does: the $\varepsilon$-pseudospectrum $\Lambda_\varepsilon(A)\coloneqq\{z\in\mathbb C:\sigma_n(zI-A)<\varepsilon\}$, the set of points at which $zI-A$ is within the distance $\varepsilon$ of a singular matrix. By \cref{cor:distance-singular}, $z\in\Lambda_\varepsilon(A)$ if and only if $z$ is an eigenvalue of $A+E$ for some $E$ with $\|E\|_2<\varepsilon$, so the pseudospectrum is the union of the spectra of all perturbations of size less than $\varepsilon$. For a normal matrix it is the union of the discs of radius $\varepsilon$ about the eigenvalues, by \cref{prop:singular-eigen}\,(iii) applied to the normal matrix $zI-A$; for $J_n(0)$ it contains the open disc of radius $\varepsilon^{1/n}$ about $0$, as the perturbations $J_\delta$ of \cref{sec:computing-eigenvalues} with complex $\delta$, $|\delta|<\varepsilon$, show. Conversely, $(zI-J_n(0))^{-1}=\sum_{j<n}N_n^j/z^{j+1}$ for $z\ne0$, whose $2$-norm is at most $n|z|^{-n}$ for $|z|<1$ and at most $n$ for $|z|\ge1$, so that $\sigma_n(zI-J_n(0))$ is at least $|z|^n/n$ or $1/n$; hence for $n\varepsilon\le1$ the pseudospectrum lies in the disc of radius $(n\varepsilon)^{1/n}$. Both radii are of the order $\varepsilon^{1/n}$, the sensitivity of \cref{tab:jordan-sensitivity}. Pseudospectra explain the transient growth of $\|A^k\|$ and of $\|e^{tA}\|$ in \cref{fig:neumann,fig:jordan-powers,fig:transient}, and the behaviour of iterative methods for nonnormal matrices; Volume~VIII returns to them.
\end{note}

\section{Functions of Matrices}
\label{sec:matrix-functions}

\subsection{Three Definitions}

A polynomial $q(x)=\sum_kc_kx^k$ can be evaluated at a square matrix, $q(A)=\sum_kc_kA^k$, and \cref{ch:eigenvalues,ch:symmetric-matrices} used such values throughout. For other functions, such as $e^x$, $\sqrt x$ or $1/x$, three definitions suggest themselves. If $A=S\Lambda S^{-1}$ is diagonalisable, one may set $f(A)\coloneqq Sf(\Lambda)S^{-1}$ with $f(\Lambda)=\operatorname{diag}(f(\lambda_1),\ldots,f(\lambda_n))$, as \cref{ch:symmetric-matrices} did for Hermitian matrices with a unitary $S$. If $f$ is the sum of a power series $\sum_kc_kx^k$, one may set $f(A)\coloneqq\sum_kc_kA^k$ whenever the series converges, as \cref{ch:vectors-matrices} did for the Neumann series. And if $A=SJS^{-1}$ is in Jordan form, one may ask what $f(J_k(\lambda))$ ought to be. For a polynomial the answer is forced. Expanding $q$ about $\lambda$ gives a finite Taylor formula,
\begin{equation}
q(x)=\sum_{j=0}^{\deg q}\frac{q^{(j)}(\lambda)}{j!}(x-\lambda)^j ,
\label{eq:polynomial-taylor}
\end{equation}
which holds for real and complex $\lambda$: writing $x=(x-\lambda)+\lambda$ in each power $x^k$ and expanding by the binomial theorem gives $q(x)=\sum_jb_j(x-\lambda)^j$ for some coefficients $b_j$, and differentiating $j$ times, with the formal derivative of \cref{prop:hermite}, and setting $x=\lambda$ gives $q^{(j)}(\lambda)=j!\,b_j$. Substituting $J_k(\lambda)$ for $x$, and using $(J_k(\lambda)-\lambda I)^j=N_k^j$, which vanishes for $j\ge k$,
\begin{equation}
q(J_k(\lambda))=\sum_{j=0}^{k-1}\frac{q^{(j)}(\lambda)}{j!}N_k^j=\begin{bmatrix}q(\lambda)&q'(\lambda)&\frac{q''(\lambda)}{2!}&\cdots&\frac{q^{(k-1)}(\lambda)}{(k-1)!}\\&q(\lambda)&q'(\lambda)&\cdots&\frac{q^{(k-2)}(\lambda)}{(k-2)!}\\&&\ddots&\ddots&\vdots\\&&&q(\lambda)&q'(\lambda)\\&&&&q(\lambda)\end{bmatrix}.
\label{eq:polynomial-jordan}
\end{equation}
For $q(x)=x^m$ this is \cref{prop:jordan-block-powers}, since $q^{(j)}(\lambda)/j!=(m^{\underline j}/j!)\lambda^{m-j}$. The value of a polynomial at a Jordan block depends only on the values of the polynomial and of its first $k-1$ derivatives at $\lambda$. This suggests the general definition.

For a nonreal eigenvalue $\lambda$, the derivatives $f^{(j)}(\lambda)$ below are complex derivatives, defined by the difference quotients of \cref{ch:differentiation} with complex increments, which Volume~II studies; for polynomials and rational functions they are given by the familiar formulas, and for sums of power series by the series differentiated term by term (the note in \cref{sec:power-series}). Only the numbers $f^{(j)}(\lambda)$ enter the definition.

\begin{definition}[Function of a matrix]
\label{def:matrix-function}
Let $A\in\mathbb C^{n\times n}$ have the pairwise different eigenvalues $\mu_1,\ldots,\mu_k$ with the indices $s_i\coloneqq\operatorname{ind}(\mu_i)$. A function $f$ is \term{defined on the spectrum} of $A$ if the numbers $f^{(j)}(\mu_i)$ exist for $0\le j<s_i$ and $i=1,\ldots,k$. For such $f$ and a Jordan form $A=SJS^{-1}$, $J=\operatorname{diag}(J_{k_1}(\lambda_1),\ldots,J_{k_r}(\lambda_r))$, let
\[
f(J_k(\lambda))\coloneqq\sum_{j=0}^{k-1}\frac{f^{(j)}(\lambda)}{j!}N_k^j,\qquad f(A)\coloneqq S\operatorname{diag}\bigl(f(J_{k_1}(\lambda_1)),\ldots,f(J_{k_r}(\lambda_r))\bigr)S^{-1}.
\]
\end{definition}

The definition depends on the choice of $S$ and on the order of the blocks only in appearance. To see this, and to connect it with the other two definitions, we need Hermite interpolation with derivatives of any order, which extends \cref{prop:hermite}.

\begin{proposition}[Hermite interpolation]
\label{prop:hermite-general}
Let $t_1,\ldots,t_k\in\mathbb C$ be pairwise different, let $m_1,\ldots,m_k\ge1$ and $M\coloneqq m_1+\cdots+m_k$, and let numbers $y_{i,j}\in\mathbb C$ be given for $0\le j<m_i$. There is exactly one polynomial $p$ of degree less than $M$ with $p^{(j)}(t_i)=y_{i,j}$ for all $i$ and all $0\le j<m_i$.
\end{proposition}

\begin{proof}
The map $p\mapsto(p^{(j)}(t_i))_{i,j}$ from $\mathcal P_{M-1}$ to $\mathbb C^M$ is linear, and both spaces have dimension $M$, so by \cref{cor:injective-surjective} it suffices to show that it is injective. Let $p\in\mathcal P_{M-1}$ satisfy $p^{(j)}(t_i)=0$ for all $i$ and $j<m_i$. By \eqref{eq:polynomial-taylor} at $\lambda=t_i$, the first $m_i$ coefficients of $p$ in powers of $x-t_i$ vanish, so $p=(x-t_i)^{m_i}r_i$ for a polynomial $r_i$, and $t_i$ is a root of $p$ of multiplicity at least $m_i$. If $p\ne0$, the factorisation of $p$ into linear factors (\cref{cor:factor}), which is unique, contains at least $m_i$ factors $x-t_i$ for each $i$, so $\deg p\ge M$, a contradiction. Hence $p=0$.
\end{proof}

\begin{theorem}[Functions of matrices are polynomials in the matrix]
\label{thm:function-polynomial}
Let $f$ be defined on the spectrum of $A\in\mathbb C^{n\times n}$, with the notation of \cref{def:matrix-function}. If a polynomial $q$ satisfies $q^{(j)}(\mu_i)=f^{(j)}(\mu_i)$ for all $i$ and $0\le j<s_i$, then $f(A)=q(A)$. Such polynomials exist; the one of least degree has degree less than $\deg m_A=s_1+\cdots+s_k$. In particular $f(A)$ does not depend on the choice of the Jordan form and of $S$, and for a polynomial $f$ it is the ordinary value $\sum_kc_kA^k$.
\end{theorem}

\begin{proof}
Every block $J_k(\lambda)$ of $J$ has $\lambda=\mu_i$ for some $i$ and $k\le s_i$. By the definition and by \eqref{eq:polynomial-jordan}, $f(J_k(\lambda))$ and $q(J_k(\lambda))$ are formed from the same numbers $f^{(j)}(\mu_i)=q^{(j)}(\mu_i)$, $j<k$, so they are equal, and therefore
\[
f(A)=Sf(J)S^{-1}=Sq(J)S^{-1}=q(A)
\]
by \cref{prop:similarity-invariants}\,(ii). The right side does not involve $S$ or $J$. A polynomial $q$ with the required properties exists by \cref{prop:hermite-general}, with $m_i=s_i$, and its degree is less than $\sum_is_i$, which is $\deg m_A$ by \cref{thm:jordan-minimal}. If $f$ is a polynomial, $q=f$ is allowed.
\end{proof}

The polynomial of least degree in the theorem is the \term{Hermite interpolation polynomial} of $f$ on the spectrum of $A$. It depends on $f$ and on $A$, so the theorem does not say that $f(A)$ is one polynomial in $A$ for all $A$; it says that for each $A$, the matrix $f(A)$ is a polynomial in $A$, determined by finitely many values of $f$ and its derivatives. Two functions that agree to the required orders at the eigenvalues of $A$ give the same matrix, however different they are elsewhere.

\subsection{Properties}

\begin{theorem}[Properties of matrix functions]
\label{thm:function-properties}
Let $f$ and $g$ be defined on the spectrum of $A\in\mathbb C^{n\times n}$, whose eigenvalues, counted with multiplicity, are $\lambda_1,\ldots,\lambda_n$.
\begin{enumerate}
\item $f(A)$ commutes with $A$, and with every matrix that commutes with $A$.
\item $f(XAX^{-1})=Xf(A)X^{-1}$ for every invertible $X$; $f(A^{\mathsf T})=f(A)^{\mathsf T}$; and if $A=\operatorname{diag}(A_1,A_2)$ with square blocks, then $f(A)=\operatorname{diag}(f(A_1),f(A_2))$.
\item The eigenvalues of $f(A)$, counted with multiplicity, are $f(\lambda_1),\ldots,f(\lambda_n)$.
\item $(cf+g)(A)=cf(A)+g(A)$ for $c\in\mathbb C$, and $(fg)(A)=f(A)g(A)$.
\item If $A=S\Lambda S^{-1}$ with $\Lambda=\operatorname{diag}(\lambda_1,\ldots,\lambda_n)$, then $f(A)=Sf(\Lambda)S^{-1}=\sum_jf(\mu_j)P_j$, with the spectral projections $P_j$ of \cref{lem:spectral-projections}; if $A=Q\Lambda Q^{\mathsf H}$ with a unitary $Q$, then $f(A)=Qf(\Lambda)Q^{\mathsf H}$.
\end{enumerate}
\end{theorem}

\begin{proof}
Let $q$ be a polynomial as in \cref{thm:function-polynomial}. (i) $f(A)=q(A)$, and a matrix that commutes with $A$ commutes with its powers and with $q(A)$. (ii) $XAX^{-1}$ has the Jordan form of $A$, with $XS$ in place of $S$, so the same $q$ serves for it, and $q(XAX^{-1})=Xq(A)X^{-1}$. By \cref{prop:transpose-similar} and \cref{thm:jordan-form}, $A^{\mathsf T}$ has the same Jordan blocks as $A$, so again the same $q$ serves, and $q(A^{\mathsf T})=q(A)^{\mathsf T}$. The Jordan form of $\operatorname{diag}(A_1,A_2)$ consists of the blocks of the Jordan forms of $A_1$ and $A_2$, so every eigenvalue of $A_i$ is one of $A$, with an index at most its index in $A$; hence $q$ also serves for $A_1$ and $A_2$, and $q(A)=\operatorname{diag}(q(A_1),q(A_2))$. (iii) $f(J)$ is upper triangular with the diagonal entries $f(\lambda_i)$, and similar matrices have the same characteristic polynomial. (iv) The linearity is clear from the definition. For the product it suffices, by $f(A)g(A)=Sf(J)g(J)S^{-1}$, to consider one block. By Leibniz's rule $(fg)^{(j)}=\sum_{a+b=j}j!\,f^{(a)}g^{(b)}/(a!\,b!)$, which follows from the product rule by induction on $j$,
\[
\begin{aligned}
f(J_k(\lambda))\,g(J_k(\lambda))&=\sum_{a,b<k}\frac{f^{(a)}(\lambda)}{a!}\frac{g^{(b)}(\lambda)}{b!}N_k^{a+b}\\
&=\sum_{j=0}^{k-1}\Bigl(\sum_{a+b=j}\frac{f^{(a)}(\lambda)\,g^{(b)}(\lambda)}{a!\,b!}\Bigr)N_k^j=\sum_{j=0}^{k-1}\frac{(fg)^{(j)}(\lambda)}{j!}N_k^j ,
\end{aligned}
\]
since $N_k^{a+b}=0$ for $a+b\ge k$. (v) For a diagonalisable matrix all blocks have size $1$, so $f(\Lambda)$ is the diagonal matrix of the $f(\lambda_i)$ and only the values of $f$ are needed. With the Lagrange polynomial $q$ that takes the values $f(\mu_j)$ at the different eigenvalues, \cref{lem:spectral-projections}\,(ii) gives $f(A)=q(A)=\sum_jq(\mu_j)P_j=\sum_jf(\mu_j)P_j$. A unitary $Q$ is the special case $S=Q$, $S^{-1}=Q^{\mathsf H}$.
\end{proof}

Part (v) shows that the definition contains the first of the three definitions, and in particular the definition $f(A)=Qf(\Lambda)Q^{\mathsf H}$ of \cref{ch:symmetric-matrices} for Hermitian matrices. Writing out the spectral projections as the Lagrange polynomials \eqref{eq:lagrange-basis} gives \term{Sylvester's formula}
\[
f(A)=\sum_{j=1}^kf(\mu_j)\prod_{i\ne j}\frac{A-\mu_iI}{\mu_j-\mu_i}
\]
for a diagonalisable $A$ with the different eigenvalues $\mu_1,\ldots,\mu_k$, which Sylvester stated in 1883 for matrices with distinct eigenvalues~\cite{sylvester1883}. For general matrices the generalised eigenspaces take the place of the eigenspaces, and the matrix splits into a diagonalisable and a nilpotent part.

\begin{proposition}[Spectral projections of a general matrix]
\label{prop:spectral-projections-general}
Let $A\in\mathbb C^{n\times n}$ have the pairwise different eigenvalues $\mu_1,\ldots,\mu_k$ with the indices $s_1,\ldots,s_k$, and let $P_i\coloneqq e_i(A)$, where $e_i$ is a function equal to $1$ in a neighbourhood of $\mu_i$ and to $0$ in a neighbourhood of each other eigenvalue.
\begin{enumerate}
\item $P_i$ is the projection onto $G_{\mu_i}(A)$ along the sum of the other generalised eigenspaces; $P_1+\cdots+P_k=I$, $P_iP_l=0$ for $i\ne l$, $P_i^2=P_i$, and each $P_i$ is a polynomial in $A$.
\item The matrix $D\coloneqq\sum_i\mu_iP_i$ is diagonalisable, $N\coloneqq A-D$ is nilpotent, $A=D+N$, and $D$ and $N$ are polynomials in $A$, so that $DN=ND$.
\item For every $f$ defined on the spectrum of $A$,
\[
f(A)=\sum_{i=1}^k\sum_{j=0}^{s_i-1}\frac{f^{(j)}(\mu_i)}{j!}(A-\mu_iI)^jP_i .
\]
\end{enumerate}
\end{proposition}

\begin{proof}
Let $A=SJS^{-1}$, with the columns of $S$ a Jordan basis. (i) The derivatives of $e_i$ vanish at every eigenvalue, so $e_i(J)$ has the blocks $I$ for the blocks of $J$ with the eigenvalue $\mu_i$ and $0$ for the others. In the Jordan basis, $P_i$ therefore keeps the coordinates along the chains for $\mu_i$, which span $G_{\mu_i}$, and annihilates the others, which span the other generalised eigenspaces; this is the projection described. The identities follow from \cref{thm:function-properties}\,(iv), since $\sum_ie_i=1$, $e_ie_l=0$ and $e_i^2=e_i$ near every eigenvalue, and $P_i$ is a polynomial in $A$ by \cref{thm:function-polynomial}. (ii) In the Jordan basis, $D$ acts on the chains for $\mu_i$ as multiplication by $\mu_i$, so it is diagonal in this basis, and $N$ acts on each chain for $\mu_i$ as $J_k(\mu_i)-\mu_iI=N_k$, so $N^n=0$. Both are polynomials in $A$ by (i). (iii) In the Jordan basis, the right side acts on a chain of length $k$ for $\mu_i$ as $\sum_{j<s_i}f^{(j)}(\mu_i)N_k^j/j!$, which is $f(J_k(\mu_i))$ because $N_k^j=0$ for $k\le j<s_i$, and the terms with other eigenvalues vanish there.
\end{proof}

For a diagonalisable matrix, $N=0$ and (iii) is part (v) of \cref{thm:function-properties}. For a general matrix, (ii) is the decomposition of an operator into commuting \emph{semisimple} and \emph{nilpotent} parts, which is unique and, unlike the Jordan basis, involves no choices~\cite[Sect.~6.8]{hoffman1971}; (iii) shows that $f(A)$ is determined by the projections $P_i$, the nilpotent parts $(A-\mu_iI)P_i$ and finitely many derivatives of $f$.

\subsection{Power Series}

The second definition, by power series, gives the same matrices. Let $f(z)=\sum_{k\ge0}c_kz^k$ with complex coefficients and the radius of convergence $R>0$, defined as in \cref{def:radius} with $|c_k|r^k$; by the note in \cref{sec:power-series}, the series and all the series obtained from it by differentiating term by term converge absolutely for $|z|<R$, and their sums are the derivatives $f^{(j)}(z)=\sum_kc_kk^{\underline j}z^{k-j}$.

\begin{theorem}[Matrix power series]
\label{thm:power-series-matrix}
Let $f(z)=\sum_kc_kz^k$ have the radius of convergence $R>0$, and let $A\in\mathbb C^{n\times n}$. If $\rho(A)<R$, the series $\sum_kc_kA^k$ converges entrywise, and its sum is $f(A)$ in the sense of \cref{def:matrix-function}. If $\rho(A)>R$, the series diverges.
\end{theorem}

\begin{proof}
Let $\rho(A)<R$, and let $f_K(z)\coloneqq\sum_{k\le K}c_kz^k$ be the partial sums, which are polynomials. For each eigenvalue $\lambda$ and each $j$, $f_K^{(j)}(\lambda)=\sum_{k\le K}c_kk^{\underline j}\lambda^{k-j}\to f^{(j)}(\lambda)$ as $K\to\infty$, because $|\lambda|<R$. By \eqref{eq:polynomial-jordan}, $f_K(J_k(\lambda))=\sum_{j<k}f_K^{(j)}(\lambda)N_k^j/j!$, so $f_K(J)\to f(J)$ entrywise, and since the entries of $f_K(A)=Sf_K(J)S^{-1}$ are fixed linear combinations of those of $f_K(J)$, $f_K(A)\to Sf(J)S^{-1}=f(A)$. If $\rho(A)>R$, let $A\mathbf x=\lambda\mathbf x$ with $|\lambda|>R$ and $x_i\ne0$. Then the $i$th component of $f_K(A)\mathbf x$ is $f_K(\lambda)x_i$, and $f_K(\lambda)$ diverges by \cref{thm:radius}\,(ii); since the components of $f_K(A)\mathbf x$ are fixed combinations of the entries of $f_K(A)$, the partial sums $f_K(A)$ do not converge.
\end{proof}

The geometric series $1/(1-z)=\sum_kz^k$ has $R=1$, and the theorem says that $\sum_kA^k$ converges when $\rho(A)<1$, with the sum $f(A)$ for $f(z)=1/(1-z)$. Since $f(z)(1-z)=1$, \cref{thm:function-properties}\,(iv) gives $f(A)(I-A)=I$: the sum is $(I-A)^{-1}$, which is \cref{cor:neumann-spectral} again. The series $\sum_{k\ge1}(-1)^{k+1}z^k/k$ of $\operatorname{Log}(1+z)$ has $R=1$ and defines $\operatorname{Log}(I+A)$ for $\rho(A)<1$; the exponential series has $R=\infty$ and defines $e^A$ for every matrix, the subject of the next section.

\subsection{Examples}

For a block of size $2$ the definition reads
\[
f(J_2(\lambda))=\begin{bmatrix}f(\lambda)&f'(\lambda)\\0&f(\lambda)\end{bmatrix}.
\]
The derivative appears as a limit of difference quotients. The matrix with the rows $(\lambda,1)$ and $(0,\lambda+h)$, $h\ne0$, has the distinct eigenvalues $\lambda$ and $\lambda+h$ with the eigenvectors $\mathbf e_1$ and $(1,h)$, and \cref{thm:function-properties}\,(v) gives it the entry $(f(\lambda+h)-f(\lambda))/h$ in position $(1,2)$ of its function $f$, which tends to $f'(\lambda)$ as the matrix tends to $J_2(\lambda)$. The defective matrix \eqref{eq:schur-example}, with the double eigenvalue $2$ of index $2$, satisfies $(A-2I)^2=0$, and its Hermite polynomial is $q(x)=f(2)+f'(2)(x-2)$, so
\[
f(A)=f(2)\,I+f'(2)\,(A-2I)
\]
for every $f$ that is differentiable at $2$. For $f(x)=\sqrt x$ this gives the square root
\[
A^{1/2}=\sqrt2\,I+\frac{A-2I}{2\sqrt2}=\frac1{2\sqrt2}\begin{bmatrix}5&1\\-1&3\end{bmatrix},\qquad\Bigl(\frac1{2\sqrt2}\Bigr)^2\begin{bmatrix}5&1\\-1&3\end{bmatrix}^2=\frac18\begin{bmatrix}24&8\\-8&8\end{bmatrix}=A ,
\]
as \cref{thm:function-properties}\,(iv) predicts from $(\sqrt x)^2=x$.

Square roots illustrate what the definition does and does not cover. For a positive semidefinite $A$, $f(x)=\sqrt x$ gives the square root $A^{1/2}$ of \cref{thm:square-root}, by part (v). For an invertible $A$ and the principal square root $\sqrt z=e^{(\operatorname{Log}z)/2}$ of \cref{sec:complex-exp}, which is complex differentiable except on the half-line $(-\infty,0]$, the definition gives a square root of $A$ whenever $A$ has no eigenvalue in $(-\infty,0)$; other branches give others. But not every square root of a matrix is a function of it. The reflections $S_\alpha$ of \cref{sec:matrices-linear-maps} all satisfy $S_\alpha^2=I_2$, while every function of $I_2$ is a multiple of $I_2$, by (v); so $I_2$ has infinitely many square roots, of which only $\pm I_2$ are of the form $f(I_2)$. And some matrices have no square root at all: if $X^2=N_2$, then $X^4=0$, so $X$ is nilpotent and $X^2=0$ (\cref{sec:rank-inequalities}), which contradicts $X^2=N_2\ne0$. The definition is consistent with this, since $\sqrt x$ is not differentiable at $0$ and so is not defined on the spectrum of $N_2$.

The inverse is the function $f(z)=1/z$, defined on the spectrum exactly when $0$ is not an eigenvalue; since $f(z)z=1$, (iv) gives $f(A)A=I$, so $f(A)=A^{-1}$. By \cref{thm:function-polynomial}, $A^{-1}$ is a polynomial in $A$ of degree less than $\deg m_A$, which sharpens the consequence of the Cayley--Hamilton theorem in \cref{sec:cayley-hamilton}. The \term{sign function}, $\operatorname{sign}z=1$ for $\operatorname{Re}z>0$ and $-1$ for $\operatorname{Re}z<0$, is defined on the spectrum of every matrix without eigenvalues on the imaginary axis, with all its derivatives equal to $0$; by \cref{prop:spectral-projections-general}\,(iii), $\operatorname{sign}A=P_+-P_-$, where $P_+$ is the sum of the spectral projections of the eigenvalues in the right half-plane and $P_-$ of those in the left half-plane. So $(I+\operatorname{sign}A)/2=P_+$ is the projection onto the sum of the generalised eigenspaces of the eigenvalues with positive real part, a quantity that matters for the stability of differential equations. Finally, the \term{principal logarithm} $\operatorname{Log}A$ is defined, with the principal logarithm $\operatorname{Log}z$ of \cref{sec:complex-exp}, for every matrix without eigenvalues on $(-\infty,0]$, and the composition rule below shows $e^{\operatorname{Log}A}=A$; a real matrix without such eigenvalues has a real principal logarithm~\cite[Chap.~11]{higham2008}.

\begin{theorem}[Composition]
\label{thm:function-composition}
Let $A\in\mathbb C^{n\times n}$ have the pairwise different eigenvalues $\mu_i$ with the indices $s_i$, and let $h=f\circ g$, where $g^{(j)}(\mu_i)$ and $f^{(j)}(g(\mu_i))$ exist for $0\le j<s_i$. Then $h$ is defined on the spectrum of $A$, $f$ is defined on the spectrum of $g(A)$, and $h(A)=f(g(A))$.
\end{theorem}

\begin{proof}
Fix $i$ and write $\mu\coloneqq\mu_i$ and $s\coloneqq s_i$. Since $g^{(s-1)}(\mu)$ and $f^{(s-1)}(g(\mu))$ exist, the derivatives of $g$ of order less than $s-1$ exist in a neighbourhood $U$ of $\mu$, and those of $f$ of order less than $s-1$ in a neighbourhood of $g(\mu)$, into which $g$ maps $U$ if $U$ is small enough, because $g$ is continuous at $\mu$. We show that for $1\le j<s$ there are polynomials $\Phi_{j,1},\ldots,\Phi_{j,j}$ in $j$ variables with integer coefficients, the same for all functions $f$ and $g$, such that
\[
h^{(j)}=\sum_{l=1}^j\bigl(f^{(l)}\circ g\bigr)\,\Phi_{j,l}\bigl(g',\ldots,g^{(j)}\bigr),
\]
on $U$ for $j<s-1$ and at $\mu$ for $j=s-1$. For $j=1$ this is the chain rule $h'=(f'\circ g)\,g'$, with $\Phi_{1,1}(y_1)=y_1$. If the formula holds on $U$ for some $j<s-1$, its right side can be differentiated, on $U$ if $j+1<s-1$ and at $\mu$ if $j+1=s-1$: by the chain rule the derivative of $f^{(l)}\circ g$ is $(f^{(l+1)}\circ g)\,g'$, and by the product rule, applied to each monomial, the derivative of $\Phi_{j,l}(g',\ldots,g^{(j)})$ is a polynomial with integer coefficients in $g',\ldots,g^{(j+1)}$. Collecting terms gives the formula for $j+1$. (The polynomials $\Phi_{j,l}$ are given explicitly by Fa\`a di Bruno's formula, but only their existence is used here.) Hence $h^{(j)}(\mu_i)$ exists for $j<s_i$, so that $h$ is defined on the spectrum of $A$, and it is the same polynomial expression in the numbers $g^{(a)}(\mu_i)$ and $f^{(l)}(g(\mu_i))$, $1\le a,l\le j$, for all $f$ and $g$. Let $p$ be a polynomial with $p^{(j)}(\mu_i)=g^{(j)}(\mu_i)$ for $j<s_i$, so that $g(A)=p(A)$. The matrix $p(J_k(\mu_i))$ is upper triangular with the single eigenvalue $p(\mu_i)$, and its Jordan blocks have size at most $k$, since $(p(J_k(\mu_i))-p(\mu_i)I)$ is a polynomial in $N_k$ without constant term, whose $k$th power vanishes. So every eigenvalue $\nu$ of $g(A)$ is one of the numbers $g(\mu_i)$, with index at most $\max\{s_i:g(\mu_i)=\nu\}$, and $f$ is defined on the spectrum of $g(A)$. Let $r$ be a polynomial with $r^{(j)}(\nu)=f^{(j)}(\nu)$ for every such $\nu$ and $j<\max\{s_i:g(\mu_i)=\nu\}$ (\cref{prop:hermite-general}). Then $f(g(A))=r(p(A))=(r\circ p)(A)$, and the polynomial $r\circ p$ has, at each $\mu_i$, the same derivatives of order less than $s_i$ as $h$, because they are given by the same expressions in the same numbers: $p^{(a)}(\mu_i)=g^{(a)}(\mu_i)$ for $a<s_i$, and $r^{(l)}(p(\mu_i))=f^{(l)}(g(\mu_i))$ for $l<s_i$, since $p(\mu_i)=g(\mu_i)$. By \cref{thm:function-polynomial}, $(r\circ p)(A)=h(A)$.
\end{proof}

With $g=\operatorname{Log}$ and $f=\exp$, $h(z)=z$ near every eigenvalue, so $e^{\operatorname{Log}A}=A$; with $g(z)=\sqrt z$ and $f(z)=z^2$ the theorem confirms $(A^{1/2})^2=A$. Higham's book~\cite[Chap.~1]{higham2008} develops these definitions and properties systematically and proves the equivalence of the definitions in the form used here.

\begin{note}{The Cauchy integral and the holomorphic functional calculus}
A fourth definition uses complex integration. If $f$ is complex differentiable in an open set that contains the spectrum of $A$, and $\Gamma$ is a closed curve in that set that winds once around each eigenvalue, then
\[
f(A)=\frac1{2\pi i}\oint_\Gamma f(z)\,(zI-A)^{-1}\,dz ,
\]
the matrix version of Cauchy's integral formula; the matrix $(zI-A)^{-1}$ is the \emph{resolvent}, and its entries are rational functions of $z$ with poles at the eigenvalues, of orders up to the indices. This definition agrees with \cref{def:matrix-function}~\cite[Sect.~1.2.3]{higham2008}, and it extends to operators on infinite-dimensional spaces, for which there are no Jordan forms; Volume~II develops Cauchy's integral formula and the spectral theory of operators.
\end{note}

\section{The Matrix Exponential}
\label{sec:matrix-exponential}

\subsection{Definition and First Properties}

\begin{definition}[Matrix exponential]
\label{def:matrix-exponential}
The \term{exponential} of $A\in\mathbb C^{n\times n}$ is
\[
e^A=\exp(A)\coloneqq\sum_{k=0}^\infty\frac{A^k}{k!}=I+A+\frac{A^2}{2!}+\frac{A^3}{3!}+\cdots.
\]
\end{definition}

\begin{proposition}[Basic properties of the exponential]
\label{prop:exp-basic}
Let $A\in\mathbb C^{n\times n}$, and let $\|\cdot\|$ be one of the operator norms $\|\cdot\|_1$, $\|\cdot\|_2$, $\|\cdot\|_\infty$.
\begin{enumerate}
\item The series converges absolutely entrywise, $\|e^A\|\le e^{\|A\|}$, and $e^0=I$.
\item $e^A$ is the function $\exp$ of $A$ in the sense of \cref{def:matrix-function}. Consequently $e^A$ commutes with every matrix that commutes with $A$, $e^{XAX^{-1}}=Xe^AX^{-1}$, $(e^A)^{\mathsf T}=e^{A^{\mathsf T}}$, $(e^A)^{\mathsf H}=e^{A^{\mathsf H}}$, $e^{\operatorname{diag}(A_1,A_2)}=\operatorname{diag}(e^{A_1},e^{A_2})$, and the eigenvalues of $e^A$ are $e^{\lambda_1},\ldots,e^{\lambda_n}$.
\item $\det e^A=e^{\operatorname{tr}A}$. In particular $e^A$ is invertible.
\end{enumerate}
\end{proposition}

\begin{proof}
(i) By \cref{prop:operator-norms}\,(iii) and~(iv), which hold for complex matrices, $|(A^k)_{ij}|\le\|A^k\|\le\|A\|^k$, so the series of the moduli of the $(i,j)$ entries is dominated by $\sum_k\|A\|^k/k!=e^{\|A\|}$ (\cref{thm:exp-series}), and converges by comparison. The partial sums satisfy $\|\sum_{k\le K}A^k/k!\|\le\sum_{k\le K}\|A\|^k/k!\le e^{\|A\|}$, and the norm of the limit is the limit of the norms, since the norms are continuous functions of the entries. (ii) The exponential series has $R=\infty$, so \cref{thm:power-series-matrix} applies to every $A$, and \cref{thm:function-properties} gives the consequences; $(e^A)^{\mathsf H}=e^{A^{\mathsf H}}$ follows from the transpose rule and from conjugating the series term by term. (iii) By (ii) and \cref{thm:trace-det-eigen}, $\det e^A=\prod_ie^{\lambda_i}=e^{\sum_i\lambda_i}=e^{\operatorname{tr}A}$, by the functional equation of the complex exponential (\cref{thm:complex-exp-properties}). Since $e^z\ne0$, $\det e^A\ne0$.
\end{proof}

Part (iii) is the identity announced in \cref{ch:determinants}. It connects the trace and the determinant, the two invariants of \cref{prop:similarity-invariants,cor:det-product}, through the exponential, as $\det(I+tB)=1+t\operatorname{tr}B+O(t^2)$ did to first order in the remarks after \cref{thm:jacobi-formula}. Jacobi's formula gives a second proof for a real $B$, once the derivative of $e^{tB}$ and its inverse $e^{-tB}$ are known (\cref{thm:exp-derivative,cor:exp-laws} below): the function $\varphi(t)\coloneqq\det e^{tB}$ satisfies $\varphi'(t)=\varphi(t)\operatorname{tr}(e^{-tB}Be^{tB})=\varphi(t)\operatorname{tr}B$ by \cref{thm:jacobi-formula}, and $\varphi(0)=1$, so $\varphi(t)=e^{t\operatorname{tr}B}$ by \cref{thm:exp-ode}.

\subsection{The Derivative and the Differential Equation}

\begin{theorem}[Derivative of the exponential]
\label{thm:exp-derivative}
For every $A\in\mathbb C^{n\times n}$, the function $t\mapsto e^{tA}$ from $\mathbb R$ to $\mathbb C^{n\times n}$ is differentiable entrywise, and
\[
\frac d{dt}e^{tA}=Ae^{tA}=e^{tA}A .
\]
\end{theorem}

\begin{proof}
The $(i,j)$ entry of $e^{tA}$ is $\sum_k(A^k)_{ij}t^k/k!$, a power series in $t$ whose coefficients satisfy $|(A^k)_{ij}|r^k/k!\le(\|A\|r)^k/k!$, a bounded sequence for every $r>0$; so its radius of convergence is $\infty$, and so are those of its real and imaginary parts. By \cref{thm:termwise-diff}, applied to both parts, the entry is differentiable with the derivative
\[
\sum_{k\ge1}(A^k)_{ij}\frac{t^{k-1}}{(k-1)!}=\sum_{k\ge0}(A^{k+1})_{ij}\frac{t^k}{k!}=\bigl(Ae^{tA}\bigr)_{ij},
\]
where the last equation holds because $(A^{k+1})_{ij}=\sum_lA_{il}(A^k)_{lj}$ and convergent series may be added term by term. Finally $Ae^{tA}=e^{tA}A$ by \cref{prop:exp-basic}\,(ii).
\end{proof}

Matrix-valued functions of $t$ are differentiated entrywise, and the product rule $(FG)'=F'G+FG'$ holds for them, with the order of the factors kept, since each entry of $FG$ is a sum of products of entries; for a vector-valued $\mathbf x(t)$, $(F\mathbf x)'=F'\mathbf x+F\mathbf x'$. With these rules, the uniqueness argument of \cref{thm:exp-ode} extends from the equation $f'=\lambda f$ to systems.

\begin{theorem}[Linear systems with constant coefficients]
\label{thm:linear-system}
Let $A\in\mathbb C^{n\times n}$, $\mathbf x_0\in\mathbb C^n$, and let $I$ be an interval containing $0$. A differentiable function $\mathbf x\colon I\to\mathbb C^n$ satisfies
\[
\mathbf x'(t)=A\mathbf x(t)\quad\text{for all }t\in I,\qquad\mathbf x(0)=\mathbf x_0 ,
\]
if and only if $\mathbf x(t)=e^{tA}\mathbf x_0$ for all $t\in I$. If $A$ and $\mathbf x_0$ are real, so is $\mathbf x$.
\end{theorem}

\begin{proof}
By \cref{thm:exp-derivative}, $\mathbf x(t)=e^{tA}\mathbf x_0$ satisfies $\mathbf x'=Ae^{tA}\mathbf x_0=A\mathbf x$, and $\mathbf x(0)=I\mathbf x_0$. Conversely, first, $e^{tA}e^{-tA}=I$ for all $t$: the derivative of $H(t)\coloneqq e^{tA}e^{-tA}$ is $Ae^{tA}e^{-tA}-e^{tA}Ae^{-tA}=0$, since $e^{tA}$ commutes with $A$, so each entry of $H$ is constant by \cref{cor:zero-derivative}, applied to its real and imaginary parts, and $H(0)=I$. Now let $\mathbf x$ be a solution and $\mathbf g(t)\coloneqq e^{-tA}\mathbf x(t)$. Then
\[
\mathbf g'(t)=-Ae^{-tA}\mathbf x(t)+e^{-tA}A\mathbf x(t)=\mathbf 0 ,
\]
so $\mathbf g$ is constant, $\mathbf g(t)=\mathbf g(0)=\mathbf x_0$, and $\mathbf x(t)=e^{tA}e^{-tA}\mathbf x(t)=e^{tA}\mathbf g(t)=e^{tA}\mathbf x_0$. If $A$ is real, the series shows that $e^{tA}$ is real.
\end{proof}

The theorem is the system version of \cref{thm:exp-ode}, proved by the same device: the combination $e^{-tA}\mathbf x(t)$ has derivative zero. Applied to each column, it says that the matrix equation $X'=AX$, $X(0)=X_0$, has the unique solution $X(t)=e^{tA}X_0$. The solution of such systems by eigenvalues, the stability of the zero solution, the phase plane and the inhomogeneous system $\mathbf x'=A\mathbf x+\mathbf f(t)$ are the subject of \cref{ch:linear-ode-systems}; here we draw the algebraic consequences.


\begin{corollary}[Exponential laws]
\label{cor:exp-laws}
Let $A,B\in\mathbb C^{n\times n}$.
\begin{enumerate}
\item $e^{(s+t)A}=e^{sA}e^{tA}$ for all real $s$ and $t$, and $(e^A)^{-1}=e^{-A}$.
\item If $AB=BA$, then $e^{A+B}=e^Ae^B=e^Be^A$.
\end{enumerate}
\end{corollary}

\begin{proof}
(i) Fix $s$. Both $X(t)\coloneqq e^{(s+t)A}$ and $Y(t)\coloneqq e^{tA}e^{sA}$ satisfy $X'=AX$, by \cref{thm:exp-derivative} and the chain rule, and both have the value $e^{sA}$ at $t=0$; by \cref{thm:linear-system}, applied to each column, they are equal, and $e^{tA}e^{sA}=e^{sA}e^{tA}$ because both are functions of $A$. For $s=1$ and $t=-1$ this gives $e^Ae^{-A}=I$. (ii) $B$ commutes with every power of $A$, hence with every partial sum of the series of $e^{tA}$ and with $e^{tA}$. So $X(t)\coloneqq e^{tA}e^{tB}$ satisfies
\[
X'(t)=Ae^{tA}e^{tB}+e^{tA}Be^{tB}=(A+B)X(t),\qquad X(0)=I ,
\]
and $X(t)=e^{t(A+B)}$ by \cref{thm:linear-system}. At $t=1$ this is the claim, and exchanging $A$ and $B$ gives the second equation.
\end{proof}

The hypothesis of (ii) cannot be dropped. For $A=N_2$ and $B=N_2^{\mathsf T}$, the series terminate, $e^A=I+A$ and $e^B=I+B$, while $A+B$ has the rows $(0,1)$ and $(1,0)$ and satisfies $(A+B)^2=I$, so that the even and the odd terms of its series give $\cosh1$ and $\sinh1$:
\[
e^Ae^B=\begin{bmatrix}1&1\\0&1\end{bmatrix}\begin{bmatrix}1&0\\1&1\end{bmatrix}=\begin{bmatrix}2&1\\1&1\end{bmatrix},\qquad e^{A+B}=\begin{bmatrix}\cosh1&\sinh1\\\sinh1&\cosh1\end{bmatrix},
\]
and $e^Ae^B\ne e^{A+B}$, which is symmetric, while $e^Be^A$ has the rows $(1,1)$ and $(1,2)$. The product $e^Ae^B$ is again an exponential when $A$ and $B$ are small, but of a matrix that involves the commutator $AB-BA$; the note at the end of this section describes it.

\subsection{The Exponential of a Jordan Block}

\begin{proposition}[Exponential of a Jordan block]
\label{prop:exp-jordan}
For $\lambda\in\mathbb C$, $k\ge1$ and real $t$,
\[
e^{tJ_k(\lambda)}=e^{\lambda t}\sum_{j=0}^{k-1}\frac{t^j}{j!}N_k^j=e^{\lambda t}\begin{bmatrix}1&t&t^2/2!&\cdots&t^{k-1}/(k-1)!\\&1&t&\cdots&t^{k-2}/(k-2)!\\&&\ddots&\ddots&\vdots\\&&&1&t\\&&&&1\end{bmatrix}.
\]
Consequently, for every $A\in\mathbb C^{n\times n}$ with the pairwise different eigenvalues $\mu_1,\ldots,\mu_k$, each entry of $e^{tA}$ has the form $\sum_ip_i(t)e^{\mu_it}$ with polynomials $p_i$ of degree less than $\operatorname{ind}(\mu_i)$.
\end{proposition}

\begin{proof}
The matrices $\lambda tI$ and $tN_k$ commute, so $e^{tJ_k(\lambda)}=e^{\lambda tI}e^{tN_k}$ by \cref{cor:exp-laws}\,(ii). The first factor is $e^{\lambda t}I$, and the series of the second terminates because $N_k^k=0$. For a general $A=SJS^{-1}$, $e^{tA}=Se^{tJ}S^{-1}$ by \cref{prop:exp-basic}\,(ii), and $e^{tJ}$ is block diagonal with the blocks just computed; each entry of $e^{tA}$ is a fixed combination of their entries.
\end{proof}

The proposition is the continuous counterpart of \cref{thm:power-entries}: where the powers of a Jordan block carry the binomial coefficients $m^{\underline j}/j!$, polynomials in $m$, its exponential carries the polynomials $t^j/j!$. A solution $\mathbf x(t)=e^{tA}\mathbf x_0$ of $\mathbf x'=A\mathbf x$ is therefore a combination of the functions $t^je^{\mu_it}$, and the factors $t^j$ occur only for defective eigenvalues. \Cref{ch:linear-higher-order} meets the same functions as solutions of linear differential equations of higher order with repeated characteristic roots, and \cref{ch:linear-ode-systems} reads off from this form when all solutions tend to zero.

\subsection{Examples}

For the diagonalisable matrix $B$ of \eqref{eq:two-by-two-examples}, with the rows $(1,1)$ and $(0,2)$, the spectral decomposition $B=P_1+2P_2$ of \cref{sec:diagonalisation} and \cref{thm:function-properties}\,(v) give
\[
e^{tB}=e^tP_1+e^{2t}P_2=e^t\begin{bmatrix}1&-1\\0&0\end{bmatrix}+e^{2t}\begin{bmatrix}0&1\\0&1\end{bmatrix}=\begin{bmatrix}e^t&e^{2t}-e^t\\0&e^{2t}\end{bmatrix},
\]
whose derivative at $t=0$ is $B$, as it must be. For the defective matrix \eqref{eq:schur-example}, $tA=2tI+t(A-2I)$ with commuting summands and $(A-2I)^2=0$, so
\[
e^{tA}=e^{2t}\bigl(I+t(A-2I)\bigr)=e^{2t}\begin{bmatrix}1+t&t\\-t&1-t\end{bmatrix},\qquad\det e^{tA}=e^{4t}=e^{t\operatorname{tr}A},
\]
and the solution of $\mathbf x'=A\mathbf x$ with $\mathbf x(0)=\mathbf e_1$ is $e^{2t}(1+t,-t)$, in which the factor $t$ comes from the Jordan block $J_2(2)$. For the rotation generator $C_{0,1}$, with the rows $(0,-1)$ and $(1,0)$, which satisfies $C_{0,1}^2=-I$, the even and the odd terms of the series give the cosine and the sine series (\cref{thm:sin-cos-series}):
\[
e^{\theta C_{0,1}}=\Bigl(1-\frac{\theta^2}{2!}+\cdots\Bigr)I+\Bigl(\theta-\frac{\theta^3}{3!}+\cdots\Bigr)C_{0,1}=\cos\theta\,I+\sin\theta\,C_{0,1}=R_\theta .
\]
Since $C_{a,b}=aI+bC_{0,1}$ with commuting summands, $e^{tC_{a,b}}=e^{at}R_{bt}$: the solutions of $\mathbf x'=C_{a,b}\mathbf x$ turn about the origin with the angular velocity $b$ while their length changes by the factor $e^{at}$. This is the continuous version of the rotation-scaling $C_{a,b}^k=r^kR_{k\theta}$ of \cref{sec:complex-eigenvalues}.

\subsection{Rotations and Unitary Matrices}

The rotation example is an instance of a general relation between antisymmetric and orthogonal matrices.

\begin{proposition}[Exponentials of antisymmetric and Hermitian matrices]
\label{prop:exp-orthogonal}
\leavevmode
\begin{enumerate}
\item If $K\in\mathbb R^{n\times n}$ is antisymmetric, then $e^K$ is orthogonal with determinant $1$. Conversely, every $Q\in\mathrm{SO}(n)$ is $e^K$ for some real antisymmetric $K$.
\item If $H\in\mathbb C^{n\times n}$ is Hermitian, then $e^{iH}$ is unitary. Conversely, every unitary matrix is $e^{iH}$ for some Hermitian $H$.
\end{enumerate}
\end{proposition}

\begin{proof}
(i) $K^{\mathsf T}=-K$ commutes with $K$, so $(e^K)^{\mathsf T}e^K=e^{-K}e^K=I$ by \cref{prop:exp-basic}\,(ii) and \cref{cor:exp-laws}, and $\det e^K=e^{\operatorname{tr}K}=e^0=1$. Conversely, let $Q\in\mathrm{SO}(n)$. By \cref{cor:orthogonal-form}\,(i), $U^{\mathsf T}QU=\operatorname{diag}(I_p,-I_q,R_{\varphi_1},\ldots,R_{\varphi_s})$ with an even $q$, since $\det Q=1$; grouping the entries $-1$ in pairs and using $-I_2=R_\pi$, we may write $U^{\mathsf T}QU=\operatorname{diag}(I_p,R_{\psi_1},\ldots,R_{\psi_l})$. Let $K_0\coloneqq\operatorname{diag}(0_p,\psi_1C_{0,1},\ldots,\psi_lC_{0,1})$, with $0_p$ the zero matrix of order $p$. It is antisymmetric, and $e^{K_0}=\operatorname{diag}(I_p,R_{\psi_1},\ldots,R_{\psi_l})$ by the block rule and the example above. So $Q=Ue^{K_0}U^{\mathsf T}=e^{UK_0U^{\mathsf T}}$, and $UK_0U^{\mathsf T}$ is antisymmetric. (ii) $(e^{iH})^{\mathsf H}=e^{(iH)^{\mathsf H}}=e^{-iH}$, and $e^{-iH}e^{iH}=I$. Conversely, a unitary $V$ is normal with eigenvalues $e^{i\theta_j}$, $\theta_j$ real (\cref{cor:schur-normal}), so $V=Q\operatorname{diag}(e^{i\theta_j})Q^{\mathsf H}=e^{iH}$ with the Hermitian $H\coloneqq Q\operatorname{diag}(\theta_j)Q^{\mathsf H}$, by \cref{prop:exp-basic}\,(ii).
\end{proof}

By \cref{thm:linear-system}, the solutions of $\mathbf x'=K\mathbf x$ with an antisymmetric $K$ are $\mathbf x(t)=e^{tK}\mathbf x_0$, and the proposition says that they keep their length: a velocity field $K\mathbf x$, which is perpendicular to $\mathbf x$ since $\langle\mathbf x,K\mathbf x\rangle=-\langle K\mathbf x,\mathbf x\rangle=0$, produces a rotation. In $\mathbb R^3$ the rotation has an explicit form.

\begin{proposition}[Rodrigues' formula]
\label{prop:rodrigues}
Let $\mathbf w\in\mathbb R^3$ be a unit vector and $K$ the matrix with $K\mathbf x=\mathbf w\times\mathbf x$. Then $K^2=\mathbf w\mathbf w^{\mathsf T}-I$, $K^3=-K$, and for every real $\theta$
\[
e^{\theta K}=I+\sin\theta\,K+(1-\cos\theta)\,K^2 .
\]
This is the rotation by $\theta$ about the axis spanned by $\mathbf w$: it fixes $\mathbf w$, and it maps every unit vector $\mathbf u\perp\mathbf w$ to $\cos\theta\,\mathbf u+\sin\theta\,(\mathbf w\times\mathbf u)$.
\end{proposition}

\begin{proof}
$K$ is antisymmetric, $K\mathbf w=\mathbf w\times\mathbf w=\mathbf 0$, and its eigenvalues are $0$ and $\pm i\|\mathbf w\|=\pm i$ (\cref{sec:complex-eigenvalues}). By \cref{cor:orthogonal-form}\,(ii), $K=U\operatorname{diag}(0,C_{0,1})U^{\mathsf T}$ for an orthogonal $U$, whose first column spans the kernel of $K$ and may be taken to be $\mathbf w$. Since $C_{0,1}^2=-I_2$, $K^2=U\operatorname{diag}(0,-1,-1)U^{\mathsf T}=\mathbf w\mathbf w^{\mathsf T}-I$, and $K^3=K\mathbf w\mathbf w^{\mathsf T}-K=-K$. Hence $K^{2j+1}=(-1)^jK$ for $j\ge0$ and $K^{2j}=(-1)^{j-1}K^2$ for $j\ge1$, and splitting the series of $e^{\theta K}$ into odd and even terms gives
\[
e^{\theta K}=I+\Bigl(\sum_{j\ge0}\frac{(-1)^j\theta^{2j+1}}{(2j+1)!}\Bigr)K+\Bigl(\sum_{j\ge1}\frac{(-1)^{j-1}\theta^{2j}}{(2j)!}\Bigr)K^2=I+\sin\theta\,K+(1-\cos\theta)\,K^2
\]
by \cref{thm:sin-cos-series}. On $\mathbf w$ the last two terms vanish. For a unit vector $\mathbf u\perp\mathbf w$, $K\mathbf u=\mathbf w\times\mathbf u$ and $K^2\mathbf u=\langle\mathbf w,\mathbf u\rangle\mathbf w-\mathbf u=-\mathbf u$, so $e^{\theta K}\mathbf u=\cos\theta\,\mathbf u+\sin\theta\,(\mathbf w\times\mathbf u)$; the vectors $\mathbf u$ and $\mathbf w\times\mathbf u$ are orthonormal and perpendicular to $\mathbf w$ (\cref{prop:cross-product}), so this is the rotation by $\theta$ in the plane perpendicular to $\mathbf w$.
\end{proof}

By \cref{prop:rotation-axis}\,(iii), every rotation of space is of this form. For $\mathbf w=(1,1,1)/\sqrt3$ and $\theta=2\pi/3$, so that $\sin\theta=\sqrt3/2$ and $1-\cos\theta=3/2$, the formula gives
\[
e^{\theta K}=I+\frac12\begin{bmatrix}0&-1&1\\1&0&-1\\-1&1&0\end{bmatrix}+\frac32\Bigl(\frac13\begin{bmatrix}1&1&1\\1&1&1\\1&1&1\end{bmatrix}-I\Bigr)=\begin{bmatrix}0&0&1\\1&0&0\\0&1&0\end{bmatrix},
\]
the cyclic permutation matrix $P$ of \cref{sec:real-schur}, now written as the exponential of an antisymmetric matrix. The formula carries the name of Rodrigues, who studied rotations and their composition in 1840~\cite{rodrigues1840}.

\subsection{The Exponential as a Limit}

\Cref{thm:exp-limit} obtained $e^x$ as the limit of $(1+x/n)^n$, with an error bound. The same holds for matrices, with the same bound in terms of the norm.

\begin{theorem}[Compound interest for matrices]
\label{thm:exp-limit-matrix}
Let $A\in\mathbb C^{n\times n}$, let $\|\cdot\|$ be one of the norms $\|\cdot\|_1$, $\|\cdot\|_2$, $\|\cdot\|_\infty$, and $a\coloneqq\|A\|$. For every integer $m\ge1$,
\[
\Bigl\|e^A-\Bigl(I+\frac Am\Bigr)^m\Bigr\|\le e^a-\Bigl(1+\frac am\Bigr)^m ,
\]
and the right side is at most $e^aa^2/(m+a)$ when $m>a$. In particular $(I+A/m)^m\to e^A$.
\end{theorem}

\begin{proof}
By the binomial theorem for the commuting matrices $I$ and $A/m$,
\[
\Bigl(I+\frac Am\Bigr)^m=\sum_{k=0}^m\binom mk\frac{A^k}{m^k},
\]
so
\[
e^A-\Bigl(I+\frac Am\Bigr)^m=\sum_{k=0}^\infty\gamma_kA^k,\qquad\gamma_k\coloneqq\frac1{k!}-\binom mk\frac1{m^k},
\]
where the binomial coefficient is $0$ for $k>m$. Each $\gamma_k$ is nonnegative, since $k!$ times the coefficient of $A^k$ in the binomial expansion is $m^{\underline k}/m^k=\prod_{j<k}(1-j/m)\le1$. Hence, by the triangle inequality and $\|A^k\|\le a^k$, applied to the partial sums and passing to the limit as in \cref{prop:exp-basic},
\[
\Bigl\|\sum_k\gamma_kA^k\Bigr\|\le\sum_k\gamma_ka^k=e^a-\Bigl(1+\frac am\Bigr)^m ,
\]
the same identity for the number $a$ in place of $A$. The bound for $m>a$ is \cref{thm:exp-limit} with $x=a$.
\end{proof}

The theorem says that compounding the small step $I+A/m$ a large number of times approximates the exponential, and that the error is at most of the order of $a^2/m$; it is the convergence of Euler's method for $\mathbf x'=A\mathbf x$ with the step $1/m$, since $(I+A/m)^m\mathbf x_0$ is the result of $m$ steps $\mathbf x\mapsto\mathbf x+A\mathbf x/m$. As a method of computing $e^A$ it is poor, because the error decreases only like $1/m$; the next section describes a better use of the same idea.

\begin{note}{Products of exponentials and Lie groups}
For matrices that do not commute, $e^{A+B}$ is still a limit of products: the \emph{Lie product formula} $e^{A+B}=\lim_{m\to\infty}(e^{A/m}e^{B/m})^m$ holds for all $A,B\in\mathbb C^{n\times n}$. Conversely, for $A$ and $B$ small enough, $e^Ae^B=e^C$ with a matrix $C$ given by the \emph{Baker--Campbell--Hausdorff series}
\[
C=A+B+\tfrac12[A,B]+\tfrac1{12}[A,[A,B]]-\tfrac1{12}[B,[A,B]]+\cdots,\qquad[A,B]\coloneqq AB-BA ,
\]
in which every term after $A+B$ is built from commutators. Both results belong to the theory of \emph{Lie groups}, of which $\mathrm{SO}(n)$ and $\mathrm U(n)$ are examples: the real antisymmetric matrices form the \emph{Lie algebra} of $\mathrm{SO}(n)$, a vector space closed under the commutator, and \cref{prop:exp-orthogonal} says that the exponential maps the Lie algebra of $\mathrm{SO}(n)$ onto the group. Hall's book~\cite{hall2015} develops this theory with matrices, including proofs of both formulas.
\end{note}

\section{Computing Matrix Functions: A First Look}
\label{sec:computing-matrix-functions}

Besides the routines \code{linalg.schur}, \code{linalg.hessenberg} and \code{linalg.solve\_sylvester} met above, the stardust library provides two routines for functions of matrices: \code{linalg.expm} computes the exponential, and \code{linalg.matrix\_power} computes integer powers. This section asks how $e^A$ should be computed and tests the answers. Moler and Van Loan surveyed the possibilities in a paper with the title \emph{Nineteen dubious ways to compute the exponential of a matrix}~\cite{moler1978}, and returned to it twenty-five years later~\cite{moler2003}; we try three of the ways. The experiments were run on the machine and with the version of the library described in \cref{sec:systems-practice}; the test matrices were formed as described below, without random numbers. Exact values of $e^A$ for the stored binary64 matrices were computed with the ball arithmetic of python-flint at a precision of $300$ bits, and errors are measured in the Frobenius norm relative to $\|e^A\|_{\mathrm F}$.

\subsection{Summing the Series}

The definition suggests summing the series $\sum_kA^k/k!$ until the terms no longer change the partial sum. \Cref{sec:expansions} found that this fails for the scalar $e^{-20}$: the terms grow to $4.3\times10^7$ before they decrease, they alternate in sign, and their sum is $2.1\times10^{-9}$, so the rounding errors of the large terms swamp the result. A matrix with eigenvalues of negative real part and a large norm behaves in the same way, and the cancellation can occur even when the eigenvalues are of moderate size, because the size of the terms is governed by $\|A\|$, not by the eigenvalues. Moler and Van Loan used the matrix
\[
A_0=\begin{bmatrix}-49&24\\-64&31\end{bmatrix}=S\begin{bmatrix}-1&0\\0&-17\end{bmatrix}S^{-1},\qquad S=\begin{bmatrix}1&3\\2&4\end{bmatrix},
\]
whose exponential is, by \cref{thm:function-properties}\,(v),
\[
e^{tA_0}=e^{-t}\begin{bmatrix}-2&3/2\\-4&3\end{bmatrix}+e^{-17t}\begin{bmatrix}3&-3/2\\4&-2\end{bmatrix},
\]
with entries of modulus at most $4e^{-t}$ for $t\ge0$, while $\|A_0\|_2=89.6$. We summed the series for $tA_0$ in binary64, forming each term $T_k=T_{k-1}(tA_0)/k$ with the matrix product of the library and adding it, and stopped when the partial sum no longer changed. \Cref{tab:exp-mvl} shows the result together with the other methods of this section. For $t=4$ the terms reach the norm $9.2\times10^{28}$ before they decay, the computed sum has an error of $6.1\times10^{14}$ relative to a result of norm $0.10$, and not a single digit is correct. In every row the error of the series lies between $3$ and $7$ times $u$ times the ratio of the largest term to $\|e^{tA_0}\|_{\mathrm F}$, the amount by which the rounding errors of the largest terms, of relative size $u$, can exceed the result.

\begin{table}
\sdfloat{%
\begin{sdtable}{colspec={rllrllll}, colsep=4pt, row{1}={font=\small\bfseries}}
\toprule
$t$ & $\|tA_0\|_2$ & largest term & terms & series & $Se^{t\Lambda}S^{-1}$ & Alg.~\ref{alg:scaling-squaring} & \code{expm} \\
\midrule
$0.25$ & $22.4$ & $7.6\times10^{1}$ & $34$ & $5.8\times10^{-15}$ & $9.0\times10^{-17}$ & $1.5\times10^{-15}$ & $7.5\times10^{-16}$ \\
$0.5$ & $44.8$ & $3.8\times10^{3}$ & $48$ & $5.3\times10^{-13}$ & $5.9\times10^{-16}$ & $5.0\times10^{-15}$ & $3.4\times10^{-15}$ \\
$1$ & $89.6$ & $1.3\times10^{7}$ & $75$ & $2.8\times10^{-9}$ & $1.5\times10^{-15}$ & $1.3\times10^{-14}$ & $7.2\times10^{-15}$ \\
$2$ & $179$ & $2.2\times10^{14}$ & $125$ & $2.2\times10^{-1}$ & $3.2\times10^{-15}$ & $2.9\times10^{-14}$ & $1.3\times10^{-14}$ \\
$4$ & $359$ & $9.2\times10^{28}$ & $192$ & $6.1\times10^{14}$ & $6.8\times10^{-15}$ & $5.6\times10^{-14}$ & $2.7\times10^{-14}$ \\
\bottomrule
\end{sdtable}}{Relative errors of $e^{tA_0}$ for the matrix $A_0$ of Moler and Van Loan, computed in binary64 by summing the Taylor series, by the eigendecomposition computed with \code{linalg.eig}, by \cref{alg:scaling-squaring} with $m=18$, and by \code{linalg.expm}. The third column gives the largest $\|T_k\|_2$ of the terms $T_k=(tA_0)^k/k!$ as computed, and the fourth the number of terms summed.\label{tab:exp-mvl}}
\end{table}

\subsection{Diagonalising}

The second way uses the first definition: compute an eigendecomposition $A=V\Lambda V^{-1}$ and form $Ve^\Lambda V^{-1}$. We computed $V$ and $\Lambda$ with \code{linalg.eig}, formed $Ve^\Lambda$ with the exponential of the library, and solved for the product with $V^{-1}$ with \code{linalg.solve}. For $A_0$ the method is the most accurate in \cref{tab:exp-mvl}: its eigenvectors $(1,2)$ and $(3,4)$ are far from parallel, and the computed $V$, with columns of unit length, has $\kappa_2(V)=11.1$, the value $\cot(\theta/2)$ of the proof of \cref{prop:diagonalising-condition} for the angle $\theta$ between them. Near a defective matrix the method fails, as \cref{sec:jordan-unstable} predicts. \Cref{tab:exp-eig} shows it for the matrices $B_\delta\coloneqq RA_\delta R^{\mathsf T}$, where $A_\delta$ is the matrix of \cref{prop:diagonalising-condition} and $R=R_{\pi/6}$ is the rotation by $30$ degrees, with the products formed in binary64. For $\delta$ from $10^{-2}$ to $10^{-6}$ the computed $\kappa_2(V)$ is $2.00/\delta$ to three digits, the bound of the proposition, and the error of the method grows roughly in proportion to it; in every row it is below $u\,\kappa_2(V)$, the amplification that the condition number permits. For $\delta=10^{-8}$ and $10^{-10}$ the computed condition numbers, $1.39\times10^8$ and $9.69\times10^7$, are below $2/\delta$, and for $\delta=10^{-10}$ the error of the method is $2.2\times10^{-16}$. The matrices $B_\delta$ have norms below $1.7$, and the series, the scaling and squaring of \cref{alg:scaling-squaring} and \code{expm} were accurate to within $4\times10^{-16}$ for all of them.

\begin{table}
\sdfloat{%
\begin{sdtable}{colspec={rlllll}, row{1}={font=\small\bfseries}}
\toprule
$\delta$ & $\kappa_2(V)$ & $Ve^\Lambda V^{-1}$ & series & Alg.~\ref{alg:scaling-squaring} & \code{expm} \\
\midrule
$10^{-2}$ & $2.00\times10^{2}$ & $7.1\times10^{-15}$ & $1.1\times10^{-16}$ & $2.3\times10^{-16}$ & $2.7\times10^{-16}$ \\
$10^{-4}$ & $2.00\times10^{4}$ & $2.1\times10^{-13}$ & $1.1\times10^{-16}$ & $2.0\times10^{-16}$ & $3.8\times10^{-16}$ \\
$10^{-6}$ & $2.00\times10^{6}$ & $1.2\times10^{-10}$ & $1.1\times10^{-16}$ & $3.6\times10^{-16}$ & $1.1\times10^{-16}$ \\
$10^{-8}$ & $1.39\times10^{8}$ & $1.4\times10^{-9}$ & $1.4\times10^{-16}$ & $1.6\times10^{-16}$ & $1.1\times10^{-16}$ \\
$10^{-10}$ & $9.69\times10^{7}$ & $2.2\times10^{-16}$ & $2.0\times10^{-16}$ & $3.7\times10^{-16}$ & $6.1\times10^{-17}$ \\
\bottomrule
\end{sdtable}}{Relative errors of $e^{B_\delta}$ for the nearly defective matrices $B_\delta=R_{\pi/6}A_\delta R_{\pi/6}^{\mathsf T}$ formed in binary64, with the condition number of the eigenvector matrix $V$ computed by \code{linalg.eig}, measured by \code{linalg.cond}.\label{tab:exp-eig}}
\end{table}

\subsection{Scaling and Squaring}

The series is accurate when $\|A\|$ is small: then its terms decrease from the start, and few of them are needed. The functional equation $e^{2B}=(e^B)^2$ (\cref{cor:exp-laws}) reduces a large matrix to a small one: with $B\coloneqq A/2^s$,
\[
e^A=\bigl(e^{B}\bigr)^{2^s},
\]
and the power is computed by squaring $s$ times. The error of the truncated series of $e^B$ is bounded as follows. Let $T_m(B)\coloneqq\sum_{k=0}^mB^k/k!$ and $b\coloneqq\|B\|\le1$ in an operator norm. Then
\[
\|e^B-T_m(B)\|\le\sum_{k>m}\frac{b^k}{k!}\le\frac{b^{m+1}}{(m+1)!}\sum_{i\ge0}\Bigl(\frac b{m+2}\Bigr)^i=\frac{b^{m+1}}{(m+1)!}\,\frac{m+2}{m+2-b},
\]
since $k!\ge(m+1)!\,(m+2)^{k-m-1}$ for $k>m$. Moreover $\|e^B\|\ge e^{-b}$, because $1=\|e^Be^{-B}\|\le\|e^B\|\,e^{b}$. For $m=18$ and $b\le1$ the relative truncation error is therefore at most $e\cdot(20/19)/19!=2.4\times10^{-17}$, below the unit roundoff $u=1.1\times10^{-16}$. \Cref{alg:scaling-squaring} states the method with these choices, evaluating $T_m(B)$ by Horner's scheme
\[
T_m(B)=I+B\Bigl(I+\frac B2\Bigl(I+\frac B3\Bigl(\cdots\Bigl(I+\frac Bm\Bigr)\cdots\Bigr)\Bigr)\Bigr).
\]

\begin{algorithm}
\begin{algorithmic}[1]
\Require $A\in\mathbb C^{n\times n}$; the degree $m$ of the Taylor polynomial ($m=18$)
\Ensure an approximation $X$ of $e^A$
\Name{ExpScalingSquaring}{A,m}
\State $s\gets$ the smallest integer $s\ge0$ with $\|A\|_1\le2^s$ \Comment{then $\|B\|_1\le1$}
\State $B\gets A/2^s$
\State $X\gets I$
\For{$k=m,m-1,\ldots,1$}
  \State $X\gets I+BX/k$ \Comment{Horner's scheme for $T_m(B)$}
\EndFor
\For{$i=1,\ldots,s$}
  \State $X\gets X^2$ \Comment{$X\approx e^{2^iB}$}
\EndFor
\State \Return $X$
\end{algorithmic}
\caption{The exponential by scaling and squaring with a Taylor polynomial. The cost is $m+s$ matrix products; the truncation error of the Taylor polynomial is below the unit roundoff, and the rounding errors of the squarings are not bounded by this simple analysis.\label{alg:scaling-squaring}}
\end{algorithm}

We implemented \cref{alg:scaling-squaring} with the matrix product of the library. In \cref{tab:exp-mvl,tab:exp-eig} it is accurate for every matrix, with errors between $1.5\times10^{-15}$ and $5.6\times10^{-14}$ for $tA_0$, growing with $t$, and within $4\times10^{-16}$ for the matrices $B_\delta$. The routine \code{linalg.expm} uses the same idea in a refined form. According to its source code, it implements the algorithm of Al-Mohy and Higham~\cite{almohy2009}: the Taylor polynomial is replaced by a Pad\'e approximant $r_m(B)=(V-U)^{-1}(V+U)$, a rational function whose numerator and denominator are polynomials of degree $m\in\{3,5,7,9,13\}$, and the degree and the number $s$ of squarings are chosen, from estimates of the norms $\|A^k\|^{1/k}$, as the cheapest pair whose bound on the backward error is below $u$; triangular matrices receive a special treatment of the diagonal. Its errors in \cref{tab:exp-mvl} are about half of those of \cref{alg:scaling-squaring}, and they grow with $t$ in the same way. The theory of the method, including the choice of the Pad\'e degree and the analysis of the squaring phase~\cite{higham2005}, is the subject of Volume~VIII. Some growth with $\|A\|$ is unavoidable for any method: the relative condition number of the exponential at $A$ is at least $\|A\|_2$~\cite[Sect.~10.2]{higham2008}, and $u\|tA_0\|_2$ is $4.0\times10^{-14}$ for $t=4$.

\subsection{Powers by Repeated Squaring}

Squaring is also the way integer powers are computed. To form $A^m$, write $m$ in binary, $m=(b_r\cdots b_0)_2$ with $b_r=1$, start with $X=A$, and for each further bit from the left replace $X$ by $X^2$ and then, if the bit is $1$, by $XA$; this uses $r=\lfloor\log_2m\rfloor$ squarings and one product for each further bit $1$, instead of the $m-1$ products of repeated multiplication. According to its source code, \code{linalg.matrix\_power} computes positive powers in exactly this way, with $\lfloor\log_2m\rfloor$ squarings and one product fewer than the number of ones in the binary expansion of $m$; for negative $m$ it first inverts the matrix. For the Fibonacci matrix $F$ of \cref{sec:matrix-multiplication}, whose powers have the entries $F_{k+1}$, $F_k$ and $F_{k-1}$, the routine returned the Fibonacci numbers exactly for every $k\le77$, and inexact values from $k=78$ on, the first power with an entry, $F_{79}=14\,472\,334\,024\,676\,221$, above $2^{53}$, where not all integers are binary64 numbers. Volume~IV treats repeated squaring as an algorithm on numbers, and Volume~VIII the accuracy of matrix powers in floating-point arithmetic.

\subsection{Transient Growth}

The last experiment concerns the size of $e^{tA}$ rather than its computation. By \cref{prop:exp-jordan}, $e^{tA}\to0$ as $t\to\infty$ when every eigenvalue has negative real part, as \cref{ch:linear-ode-systems} shows. But before it decays, $\|e^{tA}\|$ may grow, just as the powers of $A_2$ did in \cref{fig:neumann}. Two numbers bound its behaviour: the \term{spectral abscissa} $\alpha(A)\coloneqq\max\{\operatorname{Re}\lambda:\lambda\in\operatorname{spec}A\}$ and the \term{numerical abscissa} $\omega(A)\coloneqq\lambda_n\bigl((A+A^{\mathsf H})/2\bigr)$, the largest eigenvalue of the Hermitian part of $A$.

\begin{proposition}[Bounds for the exponential]
\label{prop:exp-norm-bounds}
For $A\in\mathbb C^{n\times n}$ and $t\ge0$,
\[
e^{t\alpha(A)}\le\|e^{tA}\|_2\le e^{t\omega(A)} ,
\]
and if $A=S\Lambda S^{-1}$ with a diagonal $\Lambda$, then $\|e^{tA}\|_2\le\kappa_2(S)e^{t\alpha(A)}$. The constant $\omega(A)$ is the best possible in the upper bound for small $t$: $\|e^{tA}\|_2\ge1+t\,\omega(A)+o(t)$ as $t\to0^+$.
\end{proposition}

\begin{proof}
The eigenvalues of $e^{tA}$ are the numbers $e^{t\lambda}$ (\cref{prop:exp-basic}), of modulus $e^{t\operatorname{Re}\lambda}$, and the spectral radius is at most the norm (\cref{prop:spectral-radius-norm}); this is the lower bound. For the upper bound, let $\mathbf x(t)\coloneqq e^{tA}\mathbf x_0$ and $\varphi(t)\coloneqq\|\mathbf x(t)\|^2=\mathbf x^{\mathsf H}\mathbf x$. By the product rule and \cref{thm:linear-system},
\[
\varphi'=\mathbf x'^{\mathsf H}\mathbf x+\mathbf x^{\mathsf H}\mathbf x'=\mathbf x^{\mathsf H}(A^{\mathsf H}+A)\mathbf x\le2\omega(A)\varphi ,
\]
by \cref{thm:rayleigh} applied to the Hermitian matrix $(A+A^{\mathsf H})/2$. So the function $e^{-2\omega t}\varphi(t)$, with $\omega=\omega(A)$, has a nonpositive derivative and does not increase (\cref{cor:monotone-derivative}), and $\varphi(t)\le e^{2\omega t}\varphi(0)$, that is, $\|e^{tA}\mathbf x_0\|\le e^{\omega t}\|\mathbf x_0\|$ for every $\mathbf x_0$. If $A=S\Lambda S^{-1}$, then $e^{tA}=Se^{t\Lambda}S^{-1}$ and $\|e^{t\Lambda}\|_2=\max_i|e^{t\lambda_i}|=e^{t\alpha(A)}$. Finally, for a unit eigenvector $\mathbf x_0$ of $(A+A^{\mathsf H})/2$ for $\omega$, $\varphi(0)=1$ and $\varphi'(0)=2\omega$, so $\varphi(t)=1+2\omega t+o(t)$ and $\|e^{tA}\|_2\ge\varphi(t)^{1/2}=1+\omega t+o(t)$.
\end{proof}

For a normal matrix $A=Q\Lambda Q^{\mathsf H}$ the bounds coincide: the Hermitian part is $Q\operatorname{diag}(\operatorname{Re}\lambda_i)Q^{\mathsf H}$, so that $\omega(A)=\alpha(A)$, and indeed $\|e^{tA}\|_2=e^{t\alpha(A)}$ by \cref{prop:singular-eigen}\,(iii). Bounds of this kind are the starting point of the study of transient behaviour by Trefethen and Embree~\cite{trefethen2005}. For a nonnormal matrix the norm starts with the slope $\omega(A)$, which can be positive although every eigenvalue lies in the left half-plane, and it eventually decays at the rate $\alpha(A)$. \Cref{fig:transient} shows the example
\[
A=\begin{bmatrix}-1&10\\0&-2\end{bmatrix},\qquad e^{tA}=\begin{bmatrix}e^{-t}&10(e^{-t}-e^{-2t})\\0&e^{-2t}\end{bmatrix},
\]
with $\alpha(A)=-1$, $\omega(A)=(\sqrt{101}-3)/2=3.52$, the departure from normality $\operatorname{dep}(A)=10$, and eigenvector matrices with $\kappa_2(S)\ge\sqrt{101}+10=20.05$, by the argument of \cref{prop:diagonalising-condition}. We computed $\|e^{tA}\|_2$ from the entries of $e^{tA}$ above and the formula for the singular values of a $2\times2$ matrix, and checked two values against \code{linalg.expm} and \code{linalg.matrix\_norm}. The norm rises to $2.563$ at $t=0.66$, although the solutions of the normal system with the same eigenvalues decay from the start like $e^{-t}$; afterwards it decays, and $e^t\|e^{tA}\|_2$ approaches $\sqrt{101}=10.05$, a value between the two bounds of the proposition; it is $9.55$ at $t=3$ and $10.03$ at $t=6$. \Cref{ch:linear-ode-systems} meets this behaviour in the stability of linear systems, and Volume~IX in iterative methods.

\begin{figure}
\sdfloat{%
\begin{tikzpicture}
\begin{axis}[sdplot, width=0.6\textwidth, height=0.48\textwidth, ymode=log,
  xmin=0, xmax=6, ymin=1e-3, ymax=60, xtick={0,1,2,3,4,5,6},
  ytick={1e-3,1e-2,1e-1,1,10}, xlabel={$t$}, legend style={at={(0.03,0.03)}, anchor=south west}]
\addplot[black, line width=0.9pt] coordinates {(0.00,1) (0.05,1.1897) (0.10,1.396) (0.15,1.6027) (0.20,1.7969) (0.25,1.971) (0.30,2.1213) (0.35,2.2471) (0.40,2.3491) (0.45,2.4288) (0.50,2.4881) (0.55,2.5289) (0.60,2.5533) (0.65,2.5631) (0.70,2.5602) (0.75,2.5462) (0.80,2.5226) (0.85,2.4909) (0.90,2.4521) (0.95,2.4076) (1.00,2.3582) (1.10,2.2482) (1.20,2.1281) (1.30,2.0026) (1.40,1.8751) (1.50,1.7484) (1.60,1.6244) (1.70,1.5046) (1.80,1.3899) (1.90,1.2809) (2.00,1.1781) (2.10,1.0817) (2.20,0.99154) (2.30,0.90768) (2.40,0.8299) (2.50,0.75796) (2.60,0.69159) (2.70,0.6305) (2.80,0.57436) (2.90,0.52287) (3.00,0.4757) (3.10,0.43255) (3.20,0.39313) (3.30,0.35714) (3.40,0.32432) (3.50,0.29441) (3.60,0.26717) (3.70,0.24239) (3.80,0.21984) (3.90,0.19935) (4.00,0.18073) (4.10,0.16382) (4.20,0.14847) (4.30,0.13453) (4.40,0.12189) (4.50,0.11042) (4.60,0.10001) (4.70,0.090583) (4.80,0.082034) (4.90,0.074286) (5.00,0.067264) (5.10,0.060902) (5.20,0.055138) (5.30,0.049917) (5.40,0.045188) (5.50,0.040905) (5.60,0.037027) (5.70,0.033515) (5.80,0.030335) (5.90,0.027456) (6.00,0.02485)};
\addlegendentry{$\|e^{tA}\|_2$}
\addplot[black, densely dotted, line width=0.8pt, domain=0:1.16, samples=30] {exp(3.524938*x)};
\addlegendentry{$e^{t\omega(A)}$}
\addplot[black, dashdotted, line width=0.7pt, domain=0:6, samples=30] {20.049876*exp(-x)};
\addlegendentry{$\kappa_2(S)\,e^{-t}$}
\addplot[black, dashed, line width=0.7pt, domain=0:6, samples=30] {exp(-x)};
\addlegendentry{$e^{t\alpha(A)}=e^{-t}$}
\end{axis}
\end{tikzpicture}}{Transient growth of $\|e^{tA}\|_2$ for the nonnormal matrix $A$ with the rows $(-1,10)$ and $(0,-2)$ (solid), between the lower bound $e^{t\alpha(A)}=e^{-t}$, which is also $\|e^{tD}\|_2$ for the normal matrix $D=\operatorname{diag}(-1,-2)$ with the same eigenvalues (dashed), and the upper bounds $e^{t\omega(A)}$ with $\omega(A)=3.52$ (dotted) and $\kappa_2(S)e^{-t}$ with $\kappa_2(S)=20.05$ (dash-dotted) of \cref{prop:exp-norm-bounds}. The norm rises to $2.56$ at $t\approx0.66$ before it decays, approaching $\sqrt{101}\,e^{-t}$.\label{fig:transient}}
\end{figure}

\bigskip
Every complex square matrix is unitarily similar to an upper triangular matrix with its eigenvalues on the diagonal in any order, the Schur form; it is diagonal exactly for normal matrices, and its off-diagonal part measures the departure from normality, which is also the gap in Schur's inequality $\sum|\lambda_i|^2\le\|A\|_{\mathrm F}^2=\sum\sigma_i^2$. Schur's theorem gives short proofs of the trace and determinant formulas, of the spectral mapping and the Cayley--Hamilton theorems, of the simultaneous triangularisation of commuting matrices, of the solvability of the Sylvester equation, of the density of diagonalisable matrices, of Weyl's inequalities $|\lambda_1\cdots\lambda_k|\le\sigma_1\cdots\sigma_k$, and of the existence of norms within $\varepsilon$ of the spectral radius; real matrices have a quasi-triangular real Schur form, and real normal matrices, among them the orthogonal and the antisymmetric ones, a real block diagonal form. Every nilpotent operator has a basis of Jordan chains, every complex matrix splits the space into generalised eigenspaces of dimensions $\operatorname{am}(\lambda)$, and so every matrix is similar to a Jordan form, unique up to the order of the blocks and determined by the ranks of the powers of $A-\lambda I$, which are therefore the complete invariants of similarity; the largest block for $\mu$ gives the exponent of $x-\mu$ in the minimal polynomial. The powers of a Jordan block contain binomial coefficients, which describes the entries of $A^m$, proves the formula for the solutions of linear recurrences, and decides when the powers decay, stay bounded or grow. The Jordan form is discontinuous and its transformations can be arbitrarily ill-conditioned, so computations use the Schur form. Functions of matrices are defined by the Jordan form, agree with the definitions by diagonalisation and by power series, and are polynomials in the matrix, the Hermite interpolation polynomials of the function on the spectrum. The exponential $e^{tA}$ solves $\mathbf x'=A\mathbf x$, satisfies $\det e^A=e^{\operatorname{tr}A}$ and $e^{A+B}=e^Ae^B$ for commuting matrices, has entries $\sum p_i(t)e^{\mu_it}$, maps antisymmetric matrices onto rotations, and is the limit of $(I+A/m)^m$; it is computed by scaling and squaring, not by its series or by diagonalisation, and for nonnormal matrices its norm can grow before it decays.

This chapter closes the part of the book on linear algebra. \Cref{ch:vectors-matrices,ch:linear-systems} introduced the language of vectors and matrices and the elimination that solves linear systems, \cref{ch:vector-spaces,ch:linear-maps} the spaces, bases and linear maps behind them, and \cref{ch:determinants} the determinant. \Cref{ch:inner-products} added lengths and angles, orthogonal projections and least squares, \cref{ch:eigenvalues} the eigenvalues and the powers of matrices, \cref{ch:symmetric-matrices} the spectral theorem and quadratic forms, \cref{ch:svd} the singular value decomposition, and this chapter the canonical forms and the functions of matrices. Part~IV turns to differential equations. \Cref{ch:first-order-equations} treats equations of first order, \cref{ch:linear-higher-order} linear equations of higher order, among them those with constant coefficients, whose solutions $t^je^{\lambda t}$ for repeated characteristic roots the Jordan blocks have already shown, and \cref{ch:linear-ode-systems} linear systems $\mathbf x'=A\mathbf x$, in which the matrix exponential $e^{tA}$ of this chapter is the central object: there the eigenvalues and the Jordan form of $A$ decide the shape of the solutions in the phase plane and their stability.

\chapternotes

Schur proved in 1909 that every matrix is unitarily similar to a triangular one, and derived from it the inequality $\sum|\lambda_i|^2\le\sum|a_{ij}|^2$~\cite{schur1909}. The classification of matrices up to similarity grew out of Weierstrass's theory of elementary divisors of 1868, developed for pairs of bilinear and quadratic forms~\cite{weierstrass1868}, and out of the canonical form that Jordan published in 1870 in his treatise on substitutions, for linear substitutions over finite fields~\cite{jordan1870}; Hawkins traces how the two approaches developed and were recognised as equivalent~\cite{hawkins1977}. Sylvester gave his formula for functions of matrices with distinct eigenvalues in 1883~\cite{sylvester1883}, and Higham~\cite[Chap.~1]{higham2008} surveys the history of matrix functions. Henrici introduced the departure from normality in 1962~\cite{henrici1962}, Weyl proved his inequalities between eigenvalues and singular values in 1949~\cite{weyl1949}, and Bartels and Stewart published their algorithm for the Sylvester equation in 1972~\cite{bartels1972}. Golub and Wilkinson~\cite{golubwilkinson1976} analysed the computation of the Jordan form, Moler and Van Loan~\cite{moler1978,moler2003} the computation of the exponential, and Higham~\cite{higham2005} and Al-Mohy and Higham~\cite{almohy2009} the scaling and squaring method in the form used by stardust; Trefethen and Embree~\cite{trefethen2005} develop pseudospectra and the behaviour of nonnormal matrices.

For the theory, Horn and Johnson~\cite[Chaps.~2 and~3]{horn2013} prove Schur's theorem with its consequences and the Jordan form with its real and Weyr variants, and Axler~\cite[Chap.~8]{axler2024} constructs the Jordan form from generalised eigenvectors and nilpotent operators without determinants, much as \cref{sec:nilpotent,sec:jordan-form} do. Hoffman and Kunze~\cite[Chaps.~6 and~7]{hoffman1971} treat the primary decomposition and the rational and Jordan forms through polynomials, and Gantmacher~\cite[Chaps.~V--VII]{gantmacher1959} functions of matrices and the canonical forms in the classical style. Meyer~\cite[Chap.~7]{meyer2000} presents the Jordan form and matrix functions with applications to differential equations. Higham's book~\cite{higham2008} is the standard reference for the theory and computation of matrix functions, Golub and Van Loan~\cite[Chaps.~7 and~9]{golub2013} treat the computation of the Schur form and of matrix functions, and Trefethen and Bau~\cite[Lectures~24--29]{trefethen1997} give a short introduction to eigenvalue algorithms. Hall~\cite{hall2015} develops the exponential map and Lie groups from matrices.

\chapterreferences
